\documentclass[11pt]{article}
\usepackage{fontspec}
% Charter and Menlo (macOS); elsewhere, XCharter and DejaVu Sans Mono, from which they derive.
\IfFontExistsTF{Charter}{\setmainfont{Charter}}{\setmainfont{XCharter}}
\IfFontExistsTF{Menlo}{\setmonofont{Menlo}[Scale=0.82]}{\setmonofont{DejaVu Sans Mono}[Scale=0.82]}
\usepackage{polyglossia}\setmainlanguage{english}
\usepackage{microtype}
\usepackage[a4paper,margin=27mm]{geometry}
\usepackage{setspace}\setstretch{1.1}
\usepackage{xcolor}
\usepackage[hidelinks]{hyperref}
\usepackage{amsmath,amssymb}
\usepackage{etoolbox}
\usepackage[most]{tcolorbox}
\definecolor{frise}{HTML}{2F4DA8}

\newcommand{\q}[1]{``#1''}
% The video of the current section; each entry points to its second.
\newcommand{\videoid}{}
\newcommand{\video}[1]{\renewcommand{\videoid}{#1}}
\newcommand{\topic}[1]{\section{#1}}
% A session: its date as a small heading; the topics (sections) cut across sessions.
\newcommand{\session}[1]{\par\bigskip\noindent{\small\textbf{#1}}\par\nopagebreak\smallskip}
% Tape and side: a small italic label.
\newcommand{\face}[1]{\par\medskip\noindent{\footnotesize\itshape #1}\par\nopagebreak\smallskip}
\usepackage{titlesec}
\titleformat{\section}{\large\bfseries}{}{0pt}{}
\titlespacing*{\section}{0pt}{1.6ex plus 0.5ex minus 0.3ex}{0.8ex}
\setcounter{tocdepth}{1}
% An entry: the moment on the tape (link to YouTube), then the statement.
\newlength{\horo}\settowidth{\horo}{\small\texttt{00:00:00}\quad}
% Kinds of entry, as in a course: \entree[def]{…}, [thm], [prop], [ex], [cex], [rem], [q], [disc]; none = plain.
\definecolor{kdef}{HTML}{2E6B3F}\definecolor{kthm}{HTML}{2F4DA8}\definecolor{kex}{HTML}{8A5A00}
\definecolor{kcex}{HTML}{A8322F}\definecolor{krem}{HTML}{5E5E5E}\definecolor{kq}{HTML}{6B3FA0}\definecolor{kdisc}{HTML}{0F766E}
\makeatletter
\def\kind@def#1{\textcolor{kdef}{\textbf{Definition\kindname{#1}.}}\enspace}
\def\kind@thm#1{\textcolor{kthm}{\textbf{Theorem\kindname{#1}.}}\enspace\itshape}
\def\kind@prop#1{\textcolor{kthm}{\textbf{Proposition\kindname{#1}.}}\enspace\itshape}
\def\kind@ex#1{\textcolor{kex}{\textbf{Example\kindname{#1}.}}\enspace}
\def\kind@cex#1{\textcolor{kcex}{\textbf{Counterexample\kindname{#1}.}}\enspace}
\def\kind@rem#1{\textcolor{krem}{\textbf{Remark\kindname{#1}.}}\enspace}
\def\kind@q#1{\textcolor{kq}{\textbf{Question\kindname{#1}.}}\enspace}
\def\kind@disc#1{\textcolor{kdisc}{\textbf{Discussion\kindname{#1}.}}\enspace}
\newcommand{\kindlabel}[2]{\ifcsname kind@#1\endcsname\csname kind@#1\endcsname{#2}\fi}
\newcommand{\kindname}[1]{\if\relax\detokenize{#1}\relax\else\ (#1)\fi}
\makeatother
% The moment on the tape sits in the right margin, small and pale (a link to YouTube), the statement runs full width.
\definecolor{frisepale}{HTML}{7F90C4}
\definecolor{kread}{HTML}{6B5B3E}
\def\kind@read#1{\textcolor{kread}{\textbf{Reading.}}\enspace}
\tcbset{thmbox/.style={enhanced,breakable,colback=kthm!4,colframe=kthm!55,boxrule=0.4pt,arc=1pt,left=2.2mm,right=2.2mm,top=0.8mm,bottom=0.8mm,before skip=3pt,after skip=6pt}}
\NewDocumentCommand{\entree}{O{}O{}mmm}{\ifstrequal{#1}{thm}{%
\begin{tcolorbox}[thmbox,overlay unbroken and first={\node[anchor=north west,inner sep=0pt,xshift=3.5mm,yshift=-1.4mm] at (frame.north east) {\href{https://www.youtube.com/watch?v=\videoid&t=#3s}{\textcolor{frisepale}{\texttt{\footnotesize #4}}}};}]\raggedright\kindlabel{#1}{#2}#5\end{tcolorbox}%
}{\par\raggedright\noindent\rlap{\hspace{\dimexpr\linewidth+3.5mm\relax}\href{https://www.youtube.com/watch?v=\videoid&t=#3s}{\textcolor{frisepale}{\texttt{\footnotesize #4}}}}{\kindlabel{#1}{#2}#5\par}\smallskip}}
\raggedbottom\emergencystretch=1.5em

\title{\LARGE The Buffalo Course on Algebraic Geometry (1973):\\ Index of Results}
\author{Michel Hua\\ \small Independent researcher, Paris\\ \small \texttt{michel@commutator.io}}
\date{\small November 2026}

\begin{document}
\maketitle\thispagestyle{empty}
\newpage
\tableofcontents
\newpage

\section*{Introduction}\addcontentsline{toc}{section}{Introduction}

\emph{The course.} From May to July 1973, Alexandre Grothendieck gave a course on algebraic geometry, in English, at the State University of New York at Buffalo, alongside two other courses, on topos theory and on algebraic groups. Lecture notes for its first part were written by Federico Gaeta~\cite{gaeta}. The sessions were recorded on magnetic tape. This document relies on the recordings published by the YouTube channel \q{Erwin Schulhoff} in four videos; sessions are cited by date, tape and side, as the video descriptions give them, and moments by their position in the video.

\emph{What this index is.} For each result or definition treated in a session, the index gives one line: the statement, in today's mathematical language and in our own words, and the moment of the tape where it is treated, as a link to the video accurate to the second. An entry corresponds to a statement, not to a minute of the course. Words in quotation marks are a name or a term coined in the course, kept when it matters. When what is said is sketched or incomplete, the entry gives the correct statement, with the hypotheses it requires; a claim that is false as heard is not reproduced. A passage whose audio is still doubtful is left out. Entries marked \emph{Discussion} state no result: they point to an exchange with the audience, in our words, where something is settled, corrected or left open.

\emph{What it is not.} This index is not a transcription or a summary of the course. It leaves out the explanations, the worked examples and the digressions, follows the order of the exposition only as a succession of statements, and does not quote the course. The course as delivered is a protected oral work, whose rights belong to the heirs of Alexandre Grothendieck; only the recordings are authoritative, and the reader is referred to them through the links. The statements here are standard mathematical statements, written by us.

\emph{The division into sessions.} Three sessions begin inside a side that the video description attributes to the preceding session (6, 11 and 13 June); they are indexed under their own date, as a \q{continuation} of the tape. Some sessions of this course were published together with those of the course on algebraic groups: the end of the second tape of 11 June and the next two sides of the first video of that course, and two sides of 11 July of the second. They are indexed here, at their date, with a note of the video that contains them. Two dates given by the descriptions are probably off by one day; they are flagged without being corrected.


\session{29 May 1973}
\noindent\textit{Video 1/4 (\href{https://www.youtube.com/watch?v=51wmmpHand0}{\texttt{51wmmpHand0}}): tape 1, side A (0:00:00--0:44:40); tape 1, side B (0:44:40--1:34:31).}\par\medskip
\video{51wmmpHand0}
\face{Tape 1, side A}
\topic{Affine algebraic spaces as functors}
\entree[read]{207}{0:03:27}{For the commutative algebra needed he recommends Bourbaki, which he finds more than adequate, and Zariski and Samuel, slightly old-fashioned but substantial. \cite{bourbaki,zariskisamuel}}
\entree[read]{402}{0:06:42}{For sheaves he again points to Godement’s book. \cite{godement}}
\entree[read]{467}{0:07:47}{Generalized sheaves are in exposés I and II of SGA 4, an easy reading. \cite{sga4}}
\entree[cex]{791}{0:13:11}{For an arbitrary set $I$, let $P_I=k[T_i]_{i\in I}$; a polynomial $f\in P_I$ defines a map $\tilde f\colon k^I\to k$, which in general does not determine $f$.}
\entree[prop]{920}{0:15:20}{For every $k$-algebra $k'$, $u\mapsto(u(T_i))_{i\in I}$ is a bijection from $\mathrm{Hom}_k(P_I,k')$ onto $k'^I$ (universal property of the polynomial algebra).}
\entree[prop]{983}{0:16:23}{A point $x\in k^I$ is equivalent to a homomorphism $\varepsilon_x\colon P_I\to k$, and $f(x)=\varepsilon_x(f)$.}
\entree[prop]{1047}{0:17:27}{For fixed $x$, $f\mapsto f(x)$ is a homomorphism of $k$-algebras $P_I\to k$: $(f+g)(x)=f(x)+g(x)$, $(fg)(x)=f(x)g(x)$, $(\lambda f)(x)=\lambda f(x)$.}
\entree[def]{1570}{0:26:10}{The standard affine space $E^I$ is the functor from the category of $k$-algebras to sets, $k'\mapsto k'^I$, a homomorphism $u$ acting by $u^I$.}
\entree[def]{1703}{0:28:23}{A system of equations $S=(f_j)_{j\in J}$ in $P_I$ defines a subfunctor $V_S\subset E^I$: $V_S(k')$ is the set of solutions in $k'^I$.}
\entree[prop]{2102}{0:35:02}{The solutions of $S$ in every $k$-algebra $k'$ are those of the ideal $J$ generated by $S$: $V_S=V_J$, where $V_J(k')=\{x\in k'^I : f(x)=0\ \forall f\in J\}$.}
\entree[prop]{2298}{0:38:18}{The map $J\mapsto V_J$ from ideals of $P_I$ to subfunctors of $E^I$ is injective; its values are the subfunctors defined by equations.}
\entree[prop]{2430}{0:40:30}{A functor from $k$-algebras to sets is isomorphic to some $V_S$ if and only if it is representable: $k'\mapsto\mathrm{Hom}_k(A,k')$ for a $k$-algebra $A$, denoted $V_A$.}
\entree[prop]{2499}{0:41:39}{The functor $V_A$ embeds as a subfunctor $V_S$ of $E^I$ if and only if $A$ is a quotient of $P_I$, that is, generated as a $k$-algebra by a family indexed by $I$.}
\entree[prop]{2565}{0:42:45}{The isomorphisms from $V_A$ onto some $V_S\subset E^I$ correspond to the families $(a_i)_{i\in I}$ of generators of the $k$-algebra $A$.}
\face{Tape 1, side B}
\entree[prop]{2819}{0:46:59}{For an ideal $J$ of $P_I$, $V_J(k')$ is identified, via $\mathrm{Hom}_k(P_I,k')\simeq k'^I$, with the set of homomorphisms $P_I\to k'$ vanishing on $J$, hence with $\mathrm{Hom}_k(P_I/J,k')$.}
\entree[prop]{2889}{0:48:09}{The functor $V_J$ is represented by the $k$-algebra $P_I/J$, the isomorphism being functorial in $k'$: $V_J\simeq V_{P_I/J}$.}
\entree[def]{2954}{0:49:14}{The functors $V_J\subset E^I$ are the closed algebraic subspaces of $E^I$; a functor isomorphic to one of them is an affine algebraic space over $k$.}
\entree[prop]{3290}{0:54:50}{Every $k$-algebra $A$ is isomorphic to a quotient $P_I/J$: take a generating family $(a_i)_{i\in I}$ of $A$ (for instance $I=A$) and $J$ the kernel of $T_i\mapsto a_i$.}
\entree[prop]{3423}{0:57:03}{The affine algebraic spaces over $k$ are exactly the representable functors; an embedding into $E^I$ amounts to a family of generators of the algebra indexed by $I$.}
\entree[def]{3851}{1:04:11}{A morphism $V_J\to V_{J'}$ is a natural transformation, that is, a family of maps $V_J(k')\to V_{J'}(k')$ compatible with the homomorphisms of $k'$.}
\entree[prop]{4057}{1:07:37}{A family $(F_{i'})_{i'\in I'}$ of polynomials of $P_I$ defines a morphism $E^I\to E^{I'}$, since $E^{I'}=(E^1)^{I'}$ and each polynomial gives $k'^I\to k'$ functorially in $k'$.}
\entree[prop]{4188}{1:09:48}{For functors represented by $A$ and $A'$, the morphisms from $V_A$ to $V_{A'}$ correspond bijectively to the homomorphisms of $k$-algebras $A'\to A$.}
\entree[prop]{4571}{1:16:11}{Every morphism $V_J\to V_{J'}$ is induced by a polynomial map $E^I\to E^{I'}$ sending $V_J$ into $V_{J'}$; it corresponds to a homomorphism $P_{I'}\to P_I$ sending $J'$ into $J$.}
\entree[prop]{4837}{1:20:37}{The functor $E^I$ is represented by $P_I$; a morphism $V_A\to E^I$ amounts to a homomorphism $P_I\to A$, that is, to a family of elements of $A$ indexed by $I$.}
\entree[prop]{4970}{1:22:50}{The morphism $V_A\to E^I$ defined by $P_I\to A$ is a monomorphism if and only if $P_I\to A$ is an epimorphism of rings; surjectivity is sufficient but not necessary.}
\entree[ex]{5353}{1:29:13}{The inclusion $k[T]\to k[T,T^{-1}]$ is a non-surjective epimorphism of rings ($k\neq0$): a homomorphism from $k[T,T^{-1}]$ is determined by the image of $T$, which must be invertible.}
\entree[prop]{5417}{1:30:17}{$\mathrm{Hom}_k(k[T,T^{-1}],k')$ is in bijection with the set of invertible elements of $k'$.}
\entree[cex]{5552}{1:32:32}{The monomorphism $V_{k[T,T^{-1}]}\to E^1$ is an isomorphism onto no closed subspace $V_J$ of $E^1$: otherwise $k[T]\to k[T,T^{-1}]$ would be surjective.}
\session{30 May 1973}
\noindent\textit{Video 1/4 (\href{https://www.youtube.com/watch?v=51wmmpHand0}{\texttt{51wmmpHand0}}): tape 1, side A (1:34:31--2:07:40); tape 1, side B (2:07:40--2:40:50); tape 2 (2:40:50--3:07:15).}\par\medskip
\video{51wmmpHand0}
\face{Tape 1, side A}
\entree[prop]{5809}{1:36:49}{The functor $E^I$ is the $I$-th cartesian power of $E^1$, which is represented by $k[t]$.}
\entree[def]{5871}{1:37:51}{The elements of $k'^I$ are called points of $E^I$ with values in $k'$ (or $k'$-points); the $k'$-points of $V_S$ are those that satisfy the equations of $S$.}
\entree[prop]{6129}{1:42:09}{The functors $X$ isomorphic to some $V_S$ are exactly the representable functors; if $V_S$ is defined by $J$, it is represented by $P_I/J$.}
\entree[prop]{6522}{1:48:42}{For an ideal $J$ of $P_I$: $f\in J$ if and only if $f(x')=0$ for every $k$-algebra $k'$ and every point $x'\in V_J(k')$.}
\entree[prop]{6651}{1:50:51}{It suffices to take $k'=P_I/J$ and the point corresponding to the canonical surjection $P_I\to P_I/J$ to recover $J$ as the ideal of the $f$ vanishing on $V_J$.}
\entree[prop]{6782}{1:53:02}{The map $J\mapsto V_J$ is an order-reversing bijection in both directions from the set of ideals of $P_I$ onto the set of closed subspaces of $E^I$: $J\subset J'$ if and only if $V_{J'}\subset V_J$.}
\entree[prop]{6911}{1:55:11}{The morphisms between affine algebraic spaces $X=V_A$ and $X'=V_{A'}$ correspond to the homomorphisms of $k$-algebras $A'\to A$; the category of affine algebraic spaces is anti-equivalent to that of $k$-algebras.}
\topic{Radicals and the Nullstellensatz}
\entree[prop]{7174}{1:59:34}{If one tests only on fields $k'$, then $V'_J=V'_{\widetilde J}$, where $\widetilde J=\{f : \exists n,\ f^n\in J\}$ is the radical of $J$; indeed $f(x)^n=0$ in a field implies $f(x)=0$.}
\entree[prop]{7240}{2:00:40}{The radical $\widetilde J$ is the preimage in $P_I$ of the nilradical of $P_I/J$; $P_I/\widetilde J$ is the ring $P_I/J$ modulo its nilradical, with no nonzero nilpotents.}
\entree[prop]{7367}{2:02:47}{Over fields, $V'_J\subset V'_{J'}$ if and only if $\widetilde J\supset\widetilde{J'}$; $J\mapsto V'_J$ is an order-reversing bijection between radical ideals of $P_I$ and subfunctors of $E'^I$ (restricted to fields) defined by equations.}
\entree[prop]{7502}{2:05:02}{The radical $\widetilde J$ is the set of $f\in P_I$ vanishing on $V_J(k')$ for every field $k'$ over $k$.}
\entree[prop]{7566}{2:06:06}{If $A$ is a ring and $f\in A$ satisfies $u(f)=0$ for every homomorphism $u$ from $A$ to a field, then $f$ is nilpotent.}
\face{Tape 1, side B}
\entree[prop]{7663}{2:07:43}{Every nonzero ring has a maximal ideal (Krull); it follows that the intersection of the kernels of the homomorphisms of a ring into fields is its nilradical.}
\entree[read]{7663}{2:07:43}{He comes back to Bourbaki’s commutative algebra. \cite{bourbaki}}
\entree[prop]{7725}{2:08:45}{The kernels of the homomorphisms from $A$ to fields are the prime ideals of $A$; the nilradical of $A$ is the intersection of its prime ideals.}
\entree[thm]{7993}{2:13:13}{Nullstellensatz: if $k$ is an algebraically closed field and $I$ is finite, the radical of $J\subset k[T_1,\dots,T_n]$ is the set of polynomials vanishing at the points of $k^n$ where $J$ vanishes.}
\entree[thm]{8200}{2:16:40}{Algebraic form: if $A$ is a finitely generated algebra over an algebraically closed field $k$ and $f\in A$ is killed by every $k$-homomorphism $A\to k$, then $f$ is nilpotent.}
\entree[thm]{8263}{2:17:43}{Statement (a): if $A$ is finitely generated over an algebraically closed field $k$ and $A$ is a field, then $A=k$.}
\entree[thm]{8327}{2:18:47}{Statement (b): if $A$ is a finitely generated algebra over a field $k$, its nilradical is the intersection of its maximal ideals.}
\entree[prop]{8452}{2:20:52}{If $A$ is finitely generated over an algebraically closed field $k$, $u\mapsto\ker u$ is a bijection from $\mathrm{Hom}_k(A,k)$ onto the set of maximal ideals of $A$.}
\entree[thm]{8519}{2:21:59}{The Nullstellensatz is equivalent to the conjunction of (a) and (b); (b) holds over every field, whereas (a) requires $k$ to be algebraically closed.}
\entree[thm]{8585}{2:23:05}{Statement (a$'$): over an arbitrary field $k$, a finitely generated $k$-algebra that is a field is a finite extension of $k$.}
\entree[prop]{8650}{2:24:10}{If $A$ is finitely generated over a field $k$ and $f\in A$ is killed by every $k$-homomorphism from $A$ to a finite extension of $k$, then $f$ is nilpotent.}
\entree[prop]{8717}{2:25:17}{Over a field $k$ and in finite dimension, the correspondence between closed subspaces and radical ideals holds if one tests only on finite extensions of $k$ (on $k$ itself if $k$ is algebraically closed).}
\topic{Algebras of functions}
\entree[prop]{9373}{2:36:13}{For $X=V_A$, the map $\mathrm{Hom}(X,E^1)\to A$ is a bijection of sets: the morphisms $V_A\to E^1$ correspond to the homomorphisms $k[t]\to A$, hence to the elements of $A$.}
\entree[prop]{9437}{2:37:17}{The functor $E^1$ is an algebra object over $k$: addition and multiplication $E^1\times E^1\to E^1$, and the action of each scalar of $k$.}
\entree[prop]{9504}{2:38:24}{The ring structure of $\mathrm{Hom}(X,E^1)$ deduced from that of $E^1$ is that of $A$: $A=\mathrm{Hom}(V_A,E^1)$ as a $k$-algebra.}
\face{Tape 2}
\entree[prop]{9652}{2:40:52}{An element of $A$ is identified with a natural transformation $X(k')\to k'$; the operations of $A$ correspond to pointwise operations, computed in each algebra $k'$.}
\entree[prop]{9780}{2:43:00}{A morphism $E^I\to E^1$ is given by a unique polynomial of $P_I$: the polynomial is determined by the family of functions $k'^I\to k'$, for all $k$-algebras $k'$.}
\entree[prop]{9910}{2:45:10}{For $X=V_J$, restricting polynomials to $V_J(k')$ identifies $P_I/J$ with the set of morphisms of functors $V_J\to E^1$.}
\topic{Products, sums and fibre products}
\entree[prop]{10536}{2:55:36}{The category $\mathcal{C}$ of affine algebraic spaces is opposite to that of $k$-algebras; its projective limits correspond to inductive limits of $k$-algebras.}
\entree[def]{10603}{2:56:43}{A product of the objects $X_i$ is an object $X$ equipped with projections $p_i$ such that $\mathrm{Hom}(Z,X)\to\prod_i\mathrm{Hom}(Z,X_i)$ is bijective for every $Z$; the empty family gives the final object.}
\entree[def]{10732}{2:58:52}{A sum of the $X_i$ is an object $X$ equipped with universal $p_i\colon X_i\to X$; the products of $\mathcal{C}$ are the sums of $k$-algebras.}
\entree[prop]{10867}{3:01:07}{The sum of two $k$-algebras $A'$ and $A''$ is $A'\otimes_kA''$, with the homomorphisms $x'\mapsto x'\otimes1$ and $x''\mapsto1\otimes x''$; likewise for a finite family.}
\entree[prop]{11119}{3:05:19}{For an infinite family $(A_i)$, the sum is the filtered inductive limit of the tensor products $A_{I'}=\bigotimes_{i\in I'}A_i$, with $I'$ ranging over the finite subsets of $I$.}
\entree[prop]{11182}{3:06:22}{Filtered inductive limits of $k$-algebras are computed on the underlying sets, unlike sums.}
\session{4 June 1973}
\noindent\textit{Video 1/4 (\href{https://www.youtube.com/watch?v=51wmmpHand0}{\texttt{51wmmpHand0}}): tape 1, side A (3:07:15--3:56:17); tape 1, side B (3:56:17--4:44:38); tape 2 (4:44:38--4:55:19).}\par\medskip
\video{51wmmpHand0}
\face{Tape 1, side A}
\entree[prop]{11436}{3:10:36}{Projective limits of representable functors are representable: the category $\mathcal{C}$ has all projective limits, which are computed as limits of functors.}
\entree[prop]{11506}{3:11:46}{The product of an arbitrary family $(X_i)$ of $\mathcal{C}$ is the functor $k'\mapsto\prod_iX_i(k')$; if $X_i=V_{A_i}$, it is represented by $\bigotimes_kA_i$.}
\entree[prop]{11633}{3:13:53}{For $X=V_A$ and $Y=V_B$, the product $X\times Y$ is represented by $A\otimes_kB$.}
\entree[prop]{11778}{3:16:18}{The final object of $\mathcal{C}$ is $V_k=\mathrm{Spec}\,k$ (denoted $E^0_k$), since $k$ is the initial object of the category of $k$-algebras.}
\entree[prop]{11778}{3:16:18}{To have all projective limits in a category, it suffices to have products and fibre products.}
\entree[prop]{12039}{3:20:39}{For $S=V_A$, the category $\mathcal{C}/S$ is equivalent to the opposite of the category of $A$-algebras, that is, to the category of affine spaces over $A$.}
\entree[prop]{12165}{3:22:45}{If $S$, $X$, $Y$ are represented by $A$, $B$, $C$ (with $A\to B$, $A\to C$), the fibre product $X\times_SY$ is represented by $B\otimes_AC$.}
\entree[prop]{12293}{3:24:53}{The product $X\times Y$ is the fibre product $X\times_{e}Y$ over the final object $e$.}
\entree[prop]{12487}{3:28:07}{For $S'\to S$, the base change $X\mapsto X'=X\times_SS'$ is right adjoint to the functor $\mathcal{C}/S'\to\mathcal{C}/S$ given by composition: $\mathrm{Hom}_S(Z',X)\simeq\mathrm{Hom}_{S'}(Z',X')$.}
\entree[prop]{12811}{3:33:31}{Viewed as a functor on $\mathcal{C}/S'$, $X'$ is the restriction of $X$ along $\mathcal{C}/S'\to\mathcal{C}/S$: $X'(Z')=X(Z')$.}
\entree[prop]{12939}{3:35:39}{For $S=V_A$ and $S'=V_{A'}$, the functor of $X'$ on $A'$-algebras is the composite of the functor of $X$ on $A$-algebras with the restriction of scalars $A\to A'$.}
\entree[prop]{13009}{3:36:49}{For $X=V_B$ with $B$ an $A$-algebra, the base change $A\to A'$ is represented by $B'=B\otimes_AA'$, viewed as an $A'$-algebra (also denoted $B_{A'}$).}
\topic{Kernels of pairs of arrows}
\entree[def]{13203}{3:40:03}{The kernel of a pair of arrows $u,v\colon X\to Y$ is the functor $K(Z)=\{x\in X(Z) : u(x)=v(x)\}$; it is representable.}
\entree[prop]{13331}{3:42:11}{The kernel of $(u,v)$ is the fibre product of $(u,v)\colon X\to Y\times Y$ and the diagonal $\Delta_Y\colon Y\to Y\times Y$.}
\entree[prop]{13584}{3:46:24}{The diagonal $Y=V_B\to Y\times Y$ corresponds to the augmentation $B\otimes_kB\to B$, $x\otimes y\mapsto xy$, which is surjective with kernel the ideal generated by the $x\otimes1-1\otimes x$.}
\entree[prop]{13847}{3:50:47}{If $u,v\colon B\to A$ correspond to $X\rightrightarrows Y$, the kernel of $(u,v)$ is represented by $A/I$, where $I$ is the ideal of $A$ generated by the $u(x)-v(x)$, $x\in B$.}
\face{Tape 1, side B}
\topic{Filtered limits, finite presentation and noetherian rings}
\entree[disc]{13983}{3:53:03}{A question on the Nullstellensatz leads to a warning to distinguish finite generation as module, algebra and field extension; $k(T)$ is finitely generated as a field but not as a ring.}
\entree[prop]{14373}{3:59:33}{Filtered projective limits in $\mathcal{C}$ correspond to filtered inductive limits of $k$-algebras, which are computed on the underlying sets.}
\entree[prop]{14500}{4:01:40}{Every $k$-algebra $A$ is the filtered union of its finitely generated subalgebras; every affine space is therefore a filtered projective limit of affine spaces of finite type.}
\entree[prop]{14628}{4:03:48}{A $k$-algebra is finitely generated if and only if it is a quotient of a polynomial ring in finitely many variables.}
\entree[def]{14693}{4:04:53}{$A$ is finitely presented if $A\simeq k[T_1,\dots,T_n]/J$ with $J$ generated by finitely many elements.}
\entree[prop]{14761}{4:06:01}{If $A$ is finitely presented, then for every presentation of $A$ as a quotient of a polynomial ring in finitely many variables, the ideal is finitely generated.}
\entree[thm]{14825}{4:07:05}{Hilbert basis theorem: if $k$ is a field, every ideal of $k[T_1,\dots,T_n]$ is finitely generated.}
\entree[def]{14889}{4:08:09}{A ring is noetherian if all its ideals are finitely generated, that is, if every increasing sequence of ideals is stationary.}
\entree[prop]{14955}{4:09:15}{If $k$ is noetherian, then $k[T]$ is, hence so is $k[T_1,\dots,T_n]$.}
\entree[prop]{15019}{4:10:19}{Over a noetherian ring $k$, every finitely generated algebra is finitely presented; every closed subspace of $E^n$ is defined by finitely many equations.}
\topic{Group objects, kernels and stabilizers}
\entree[prop]{15218}{4:13:38}{Giving a section of an object $X$ over the final object amounts to giving a point $x\in X(k)$, compatible with every $X(Z)$; a group object has one, the unit section.}
\entree[prop]{15347}{4:15:47}{The kernel of a homomorphism of affine groups $G\to H$, $K(Z)=\ker(G(Z)\to H(Z))$, is an affine group: it is the fibre product of $G$ and the final object over $H$ via the unit section.}
\entree[prop]{15479}{4:17:59}{If $G=V_A$, $H=V_B$ and $B^+$ is the augmentation ideal of $B$, the kernel of $G\to H$ is represented by $A\otimes_Bk=A/B^+A$.}
\entree[prop]{15555}{4:19:15}{If $B\to A$ is a homomorphism and $J$ an ideal of $B$, then $A\otimes_B(B/J)\simeq A/JA$.}
\entree[prop]{15814}{4:23:34}{If a group $G$ acts on $X$ and $x\in X(k)$, the stabilizer $G_x(Z)=\{g\in G(Z) : g\cdot x(Z)=x(Z)\}$ is representable; it is the preimage of $x$ under $g\mapsto g\cdot x$.}
\entree[prop]{15944}{4:25:44}{The stabilizer of a $k$-rational point is therefore an affine subgroup; the intersection of finitely many such stabilizers $G_{x_i}$ is again one.}
\entree[prop]{16010}{4:26:50}{The intersection of the $G_{x_i}$, $x_i\in X_i(k)$, is the stabilizer of the point $(x_i)$ of $\prod_iX_i$.}
\entree[prop]{16207}{4:30:07}{For subobjects $X_i\to X$ (monomorphisms), the intersection subfunctor is the fibre product of the $X_i$ over $X$; it is therefore representable, even for an infinite family.}
\topic{Inductive limits and quotients}
\entree[prop]{16465}{4:34:25}{The category $\mathcal{C}$ has all inductive limits, since the category of $k$-algebras has all projective limits.}
\entree[prop]{16527}{4:35:27}{Projective limits of $k$-algebras are computed on the underlying sets, with the induced algebra structure.}
\entree[def]{16658}{4:37:38}{A double arrow $R\rightrightarrows X$, that is, $R\to X\times X$, is an equivalence relation if it is a monomorphism and each $R(Z)$ is the graph of an equivalence relation on $X(Z)$.}
\entree[prop]{16789}{4:39:49}{The cokernel of a double arrow $X\rightrightarrows Y$ of $\mathcal{C}$, corresponding to $u,v\colon B\to A$, is represented by the subalgebra $\{x\in B : u(x)=v(x)\}$.}
\entree[ex]{16981}{4:43:01}{For a subgroup $H$ of $G$, the arrows $(h,g)\mapsto g$ and $(h,g)\mapsto hg$ from $H\times G$ to $G$ define an equivalence relation on $G$, whose classes are the $Hg$.}
\face{Tape 2}
\entree[disc]{17146}{4:45:46}{Discussion of why $H\times G\to G\times G$, $(h,g)\mapsto(g,hg)$, composed with an isomorphism, is a monomorphism, so that the relation defining the quotient $G/H$ is a genuine equivalence relation.}
\entree[ex]{17211}{4:46:51}{The morphism $(h,g)\mapsto(g,hg)$ from $H\times G$ to $G\times G$ is a monomorphism, and defines an equivalence relation because $H$ is a subgroup.}
\entree[def]{17406}{4:50:06}{The standard Borel subgroup $B\subset\mathrm{GL}_n$ is the functor $k'\mapsto$ group of invertible triangular matrices over $k'$, with invertible diagonal coefficients $\lambda_i$.}
\entree[prop]{17469}{4:51:09}{The group $\mathrm{GL}_n$ is represented by $k[T_{ij}]$ with the inverse of the determinant adjoined.}
\entree[prop]{17536}{4:52:16}{Over a field $k$, the functions on $\mathrm{GL}_n$ invariant under translation by $B$ are constant: the cokernel of $\mathrm{GL}_n\times B\rightrightarrows\mathrm{GL}_n$ in affine spaces is the final object.}
\entree[prop]{17536}{4:52:16}{The quotient $\mathrm{GL}_n/B$ in the sense of schemes is a complete variety (flag variety), whose global functions are the constants.}
\entree[disc]{17536}{4:52:16}{Functions on $G$ invariant under a Borel subgroup are constant, so the affine quotient collapses to a point while the true quotient is a complete flag variety; for a scalar subgroup it is affine.}
\entree[prop]{17599}{4:53:19}{The quotient of $\mathrm{GL}_n$ by the subgroup of scalars is an affine variety.}
\session{5 June 1973}
\noindent\textit{Video 1/4 (\href{https://www.youtube.com/watch?v=51wmmpHand0}{\texttt{51wmmpHand0}}): tape 1, side A (4:55:19--5:43:08); tape 1, side B (5:43:08--6:31:22); tape 2 (6:31:22--6:38:51).}\par\medskip
\video{51wmmpHand0}
\face{Tape 1, side A}
\topic{Closed subspaces, flatness and dual numbers}
\entree[disc]{17791}{4:56:31}{A question about $K[T]$ versus $K[T,T^{-1}]$ and the line minus the origin: the right functorial notion is invertible elements of a ring, not nonzero elements, which would not give a functor.}
\entree[prop]{17855}{4:57:35}{The functor represented by $K[T]$ is $K'\mapsto K'$; that of $K[T,T^{-1}]$ is $K'\mapsto K'^{*}$ (invertible elements), the inclusion $K'^{*}\subset K'$ coming from $K[T]\to K[T,T^{-1}]$.}
\entree[prop]{18250}{5:04:10}{If $X$ is represented by $A$, the point $\mathrm{id}_A\in X(A)$ is universal: every point of $X(K')=\mathrm{Hom}(A,K')$ is the image of $\mathrm{id}_A$ under a unique homomorphism $A\to K'$.}
\entree[prop]{18383}{5:06:23}{An epimorphism of $K$-algebras $K'\to K''$ corresponds to a monomorphism of functors, hence to a subfunctor.}
\entree[cex]{18446}{5:07:26}{An epimorphism of rings is not necessarily surjective; the surjective homomorphisms correspond to closed immersions.}
\entree[prop]{18708}{5:11:48}{For $X$ with ring $A$, the subfunctor $Y_J\colon K'\mapsto\{u\in\mathrm{Hom}(A,K')\mid u(J)=0\}$ is a closed subspace, and $J\mapsto Y_J$ is a bijection from the ideals of $A$ onto the closed subspaces of $X$.}
\entree[prop]{18773}{5:12:53}{The ideal $J$ is recovered as the set of $f\in A$ such that $f(x)=0$ for every $K'$ and every $x\in Y(K')$.}
\entree[disc]{18837}{5:13:57}{A doubt whether ideals correspond to all subfunctors of an affine space; the lecturer answers that the bijection is only with the closed algebraic subspaces, defined by vanishing conditions.}
\entree[prop]{18901}{5:15:01}{This correspondence concerns all ideals and not only radical ideals, because the test algebras $K'$ are arbitrary.}
\entree[prop]{18901}{5:15:01}{The ring $A$ is identified with $\mathrm{Hom}(X,\mathbb{A}^1_K)$: $f\in A$ defines $x\mapsto f(x)=x(f)\in K'$, functorially in $K'$.}
\entree[prop]{19097}{5:18:17}{A $K$-module $M$ is flat if and only if $M\otimes_K K'\to M\otimes_K K''$ is injective for every injective homomorphism of $K$-algebras $K'\to K''$.}
\entree[disc]{19097}{5:18:17}{Question on why flatness of a module can be tested by tensoring with injections of algebras only; the lecturer sketches how this implies the usual definition.}
\entree[def]{19161}{5:19:21}{The algebra of dual numbers $D_K(N)=K\times N$ with $(\lambda,x)(\lambda',x')=(\lambda\lambda',\lambda x'+\lambda'x)$; $N$ is an ideal of square zero in it, and $D_K(N)$ is functorial in $N$.}
\entree[prop]{19362}{5:22:42}{One has $M\otimes_K D_K(N)\simeq M\oplus(M\otimes_K N)$, functorially in $N$; this reduces the flatness of $M$ to injectivity for algebras.}
\entree[read]{19425}{5:23:45}{The principle of idealization is in Nagata’s book. \cite{nagata}}
\entree[prop]{19491}{5:24:51}{Projective limits of affine spaces correspond to inductive limits of algebras: products and fibre products correspond to the tensor products $\otimes_K$ and $A\otimes_C B$.}
\entree[prop]{19491}{5:24:51}{The kernel of a double arrow $u,v\colon A\rightrightarrows B$ of algebras corresponds to the quotient of $B$ by the ideal generated by the $u(x)-v(x)$.}
\entree[prop]{19557}{5:25:57}{The kernel of a homomorphism of affine algebraic groups and the stabilizer of a rational point of an affine space equipped with an action of $G$ are closed algebraic subgroups of $G$.}
\entree[rem]{19684}{5:28:04}{The embedding of affine spaces into functors on algebras commutes with projective limits but not with inductive limits.}
\topic{Sums and connectedness}
\entree[prop]{19812}{5:30:12}{The finite sum of affine spaces $X_i$ with rings $A_i$ is the affine space with ring $\prod_i A_i$; the canonical injections correspond to the projections $\prod_j A_j\to A_i$.}
\entree[prop]{19886}{5:31:26}{The morphisms $X_i\to\coprod_j X_j$ are closed immersions (the projections are surjective), hence injective on points with values in any $K'$.}
\entree[def]{20208}{5:36:48}{An object $X$ of a category is connected if every decomposition $X=X'\sqcup X''$ forces $X'$ or $X''$ to be initial; in topological spaces one recovers connectedness.}
\entree[prop]{20350}{5:39:10}{The functor $K'\mapsto\mathrm{Hom}(0,K')$ of the zero algebra has a single element if $K'=0$ and is empty otherwise; the corresponding affine space, \q{empty}, is the initial object.}
\entree[prop]{20474}{5:41:14}{A decomposition of an affine space with ring $A$ as a sum of two affine spaces is equivalent to a decomposition of $A$ as a product of two algebras $A'\times A''$.}
\entree[prop]{20539}{5:42:19}{A decomposition $A=A'\times A''$ is equivalent to the data of orthogonal idempotents $e',e''$ of $A$ with sum $1$ (the images of the units of $A'$ and $A''$).}
\face{Tape 1, side B}
\entree[prop]{20656}{5:44:16}{A decomposition of $A$ as a product of $n$ rings is equivalent to $n$ pairwise orthogonal idempotents with sum $1$; one recovers $A_i=A/(1-e_i)$ and $A\simeq\prod_i A_i$.}
\entree[prop]{20722}{5:45:22}{If $e$ is idempotent, so is $1-e$ and $e(1-e)=0$: a pair of orthogonal idempotents with sum $1$ is equivalent to a single idempotent.}
\entree[prop]{20786}{5:46:26}{A ring $A$ is connected (its affine space is) if and only if its only idempotents are $0$ and $1$.}
\entree[def]{20850}{5:47:30}{An object is $0$-connected if it is connected and not initial; a ring is $0$-connected if it is connected and nonzero.}
\entree[prop]{21041}{5:50:41}{For $Z$ affine, a morphism $Z\to\coprod_{i}X_i$ (finite sum) is equivalent to a decomposition $Z=\coprod_i Z_i$ together with morphisms $Z_i\to X_i$.}
\entree[prop]{21173}{5:52:53}{If $K'$ is $0$-connected, the map $\coprod_i X_i(K')\to X(K')$ for $X=\coprod_i X_i$ (finite sum) is bijective.}
\entree[prop]{21173}{5:52:53}{An integral domain is $0$-connected.}
\entree[prop]{21237}{5:53:57}{A local ring is $0$-connected.}
\entree[rem]{21237}{5:53:57}{Announced: an affine space is connected (resp. empty) if and only if the topological space $\mathrm{Spec}\,A$ is.}
\entree[cex]{21337}{5:55:37}{For $K'=0$ the hypothesis \q{nonzero} is necessary: $X(0)$ is a singleton whereas the sum of the $X_i(0)$ has $n$ elements.}
\topic{Points with values in fields and the spectrum}
\entree[def]{21723}{6:02:03}{Points $u'\in X(K')$, $u''\in X(K'')$, with values in fields over $K$, are equivalent if they have the same image in $X(K''')$ for some $K$-homomorphisms $K',K''\to K'''$ into a field $K'''$.}
\entree[prop]{21859}{6:04:19}{Two homomorphisms $u'\colon A\to K'$, $u''\colon A\to K''$ into fields are equivalent if and only if $\ker u'=\ker u''$; this is therefore an equivalence relation.}
\entree[prop]{21989}{6:06:29}{The kernel of a homomorphism from $A$ into a field is a prime ideal $P$; the homomorphisms $A\to K'$ with kernel $P$ correspond to the field homomorphisms $\kappa(P)=\mathrm{Frac}(A/P)\to K'$.}
\entree[prop]{22052}{6:07:32}{Two extensions $K'$, $K''$ of a field $\kappa$ have a common extension: $K'\otimes_\kappa K''$ is nonzero, hence has a quotient that is a field.}
\entree[prop]{22123}{6:08:43}{The set of classes of points with values in fields is identified with the set $\mathrm{Spec}(A)$ of prime ideals of $A$.}
\topic{The Zariski topology on the spectrum}
\entree[prop]{22510}{6:15:10}{For $S\subset A$, $V(S)$ is the set of prime ideals containing $S$; $V(S)=V(J_S)=V(\sqrt{J_S})$, where $J_S$ is the ideal generated by $S$.}
\entree[prop]{22575}{6:16:15}{The radical of an ideal is the intersection of the prime ideals containing it; $J\mapsto V(J)$ is therefore a bijection from radical ideals onto the subsets of the form $V(S)$.}
\entree[prop]{22644}{6:17:24}{For radical $J$, $J'$, $V(J)\subset V(J')$ if and only if $J\supset J'$.}
\entree[prop]{22774}{6:19:34}{For every family of ideals, $\bigcap_i V(J_i)=V\bigl(\sum_i J_i\bigr)$.}
\entree[prop]{22837}{6:20:37}{For a finite family of ideals, $\bigcup_i V(J_i)=V\bigl(\bigcap_i J_i\bigr)=V\bigl(\prod_i J_i\bigr)$.}
\entree[prop]{23031}{6:23:51}{The $V(J)$ are the closed sets of a topology on $\mathrm{Spec}(A)$, in bijection with the radical ideals of $A$.}
\entree[prop]{23031}{6:23:51}{The space $\mathrm{Spec}(A)$ is quasi-compact.}
\entree[prop]{23167}{6:26:07}{One has $V(J)=\varnothing$ if and only if $J=A$ (every proper ideal is contained in a prime ideal).}
\entree[prop]{23294}{6:28:14}{The closure of the point $x$ with prime ideal $P$ is $V(P)$; the point $y$ with ideal $Q$ is in the closure of $x$ if and only if $Q\supset P$.}
\entree[prop]{23357}{6:29:17}{Specialization on $\mathrm{Spec}(A)$ corresponds to inclusion of prime ideals in the reverse direction; the space $\mathrm{Spec}(A)$ is sober: every irreducible closed set is the closure of a unique point.}
\face{Tape 2}
\topic{Irreducible spaces}
\entree[prop]{23485}{6:31:25}{A space is irreducible if it is nonempty and is not the union of two strictly smaller closed sets; the closure of an irreducible subset is irreducible.}
\entree[prop]{23485}{6:31:25}{The closure of a point is irreducible; the image of an irreducible space under a continuous map is irreducible.}
\entree[prop]{23547}{6:32:27}{A Hausdorff space is irreducible if and only if it is reduced to a point.}
\entree[prop]{23611}{6:33:31}{The spectrum of $A$, with its topology, does not change when $A$ is replaced by $A/\mathrm{nil}(A)$, because every prime ideal contains the nilradical.}
\entree[prop]{23676}{6:34:36}{One has $\mathrm{Spec}(A)=\varnothing$ if and only if $A=0$.}
\entree[prop]{23738}{6:35:38}{For $A$ reduced and $J'$, $J''$ radical, $\mathrm{Spec}(A)=V(J')\cup V(J'')$ if and only if $J'J''=0$.}
\entree[prop]{23801}{6:36:41}{The space $\mathrm{Spec}(A)$ is irreducible if and only if $A/\mathrm{nil}(A)$ is an integral domain.}
\session{6 June 1973}
\noindent\textit{Video 1/4 (\href{https://www.youtube.com/watch?v=51wmmpHand0}{\texttt{51wmmpHand0}}): tape 2, cont. (6:38:51--7:19:33); tape 1 (7:19:33--8:08:05); tape 2 (8:08:05--8:10:17).}\par\medskip
\video{51wmmpHand0}
\face{Tape 2, cont.}
\entree[def]{24256}{6:44:16}{The subsets $V(S)$ are the closed sets of a topology on $\mathrm{Spec}(A)$, the spectral (Zariski) topology; $V(J)=V(J')$ if and only if $\sqrt J=\sqrt{J'}$.}
\entree[prop]{24326}{6:45:26}{The nilradical of every ring is the intersection of its prime ideals.}
\entree[prop]{24456}{6:47:36}{The points with values in a field $K'$ of class $x\in\mathrm{Spec}(A)$ correspond to the $K$-homomorphisms $\kappa(x)\to K'$.}
\entree[prop]{24518}{6:48:38}{A nonempty space is irreducible if and only if two nonempty open sets meet, if and only if every nonempty open set is dense.}
\entree[prop]{24714}{6:51:54}{For every ideal $J$ of $A$, the closed set $V(J)$, with the induced topology, is homeomorphic to $\mathrm{Spec}(A/J)$.}
\entree[prop]{25035}{6:57:15}{The closed set $V(J)$ is irreducible if and only if $\sqrt J$ is prime; $P\mapsto V(P)$ is a bijection from $\mathrm{Spec}(A)$ onto the irreducible closed sets.}
\entree[prop]{25167}{6:59:27}{This bijection is $x\mapsto\overline{\{x\}}$: $\mathrm{Spec}(A)$ is sober, every irreducible closed set having a unique generic point.}
\entree[ex]{25297}{7:01:37}{The spectrum of an integral domain $A$ is Hausdorff if and only if $A$ is a field; those of $\mathbb{Z}$, of $k[T]$ or of a discrete valuation ring are not.}
\topic{Irreducible components}
\entree[def]{25433}{7:03:53}{An irreducible component of a space is an irreducible closed set that is maximal among irreducible closed sets.}
\entree[prop]{25498}{7:04:58}{Every irreducible subset of a topological space is contained in an irreducible component (Zorn's lemma); the space is therefore the union of its irreducible components.}
\entree[prop]{25626}{7:07:06}{A filtered union of irreducible subsets is irreducible.}
\entree[prop]{25626}{7:07:06}{In a Hausdorff space, the irreducible components are the points.}
\entree[prop]{25689}{7:08:09}{The irreducible components of $\mathrm{Spec}(A)$ correspond to the minimal prime ideals of $A$ (minimal among all primes).}
\entree[prop]{25949}{7:12:29}{For $X$ affine of finite type over $\mathbb{C}$, with ring $A$, the points of $X(\mathbb{C})$ correspond to the maximal ideals of $A$, that is, to the closed points of $\mathrm{Spec}(A)$.}
\entree[prop]{26013}{7:13:33}{For $X$ of finite type over $\mathbb{C}$, $X$ is connected if and only if $X(\mathbb{C})$, with its complex topology, is connected.}
\entree[prop]{26144}{7:15:44}{A space is irreducible if and only if it has exactly one irreducible component; the empty space has none.}
\entree[prop]{26208}{7:16:48}{In a sober space, the irreducible components correspond to the maximal points for the specialization relation (generic points).}
\face{Tape 1}
\entree[prop]{26376}{7:19:36}{In a quasi-compact space, every irreducible closed set contains a minimal irreducible closed set (Zorn's lemma).}
\entree[prop]{26508}{7:21:48}{A nonempty subset equal to the closure of each of its points is a minimal irreducible closed set, and conversely.}
\entree[prop]{26698}{7:24:58}{In a sober space, the minimal irreducible closed sets are the closed points; if the space is moreover quasi-compact, every nonempty closed set contains a closed point.}
\entree[prop]{26762}{7:26:02}{The closed points of $\mathrm{Spec}(A)$ are the maximal ideals; the existence of a closed point for $A\neq0$ is the existence of a maximal ideal (Krull).}
\topic{Noetherian spaces}
\entree[def]{26826}{7:27:06}{A space is noetherian if every increasing sequence of open sets is stationary, or equivalently every decreasing sequence of closed sets.}
\entree[prop]{26888}{7:28:08}{A space is noetherian if and only if all its open sets are quasi-compact.}
\entree[prop]{27083}{7:31:23}{In a noetherian space, every subset is quasi-compact and noetherian.}
\entree[prop]{27151}{7:32:31}{A noetherian Hausdorff space is finite discrete.}
\entree[prop]{27151}{7:32:31}{If $A$ is a noetherian ring, $\mathrm{Spec}(A)$ is a noetherian space.}
\entree[prop]{27214}{7:33:34}{Fields and principal rings are noetherian; if $A$ is noetherian, so is $A[T]$; a ring is noetherian if and only if its ideals are finitely generated.}
\entree[prop]{27282}{7:34:42}{A noetherian space has finitely many irreducible components; a noetherian ring therefore has finitely many minimal prime ideals.}
\entree[disc]{27345}{7:35:45}{Several people reconstruct together a proof that a Noetherian space has finitely many irreducible components, arriving at Noetherian induction on a minimal closed set lacking the property.}
\entree[prop]{27535}{7:38:55}{In a noetherian space, every nonempty family of closed sets has a minimal element (noetherian induction).}
\entree[cex]{27747}{7:42:27}{A non-noetherian ring can have a noetherian spectrum: $K\oplus V$, with $V$ infinite-dimensional and $V^2=0$, has a spectrum reduced to a point.}
\entree[disc]{27747}{7:42:27}{Can a non-Noetherian ring have a Noetherian spectrum? Yes, for instance a ring with a single prime whose nilradical is not finitely generated; chain conditions on radical versus prime ideals are left open.}
\entree[prop]{27809}{7:43:29}{The space $\mathrm{Spec}(A)$ is noetherian if and only if the radical ideals of $A$ satisfy the ascending chain condition.}
\entree[prop]{28019}{7:46:59}{The points of $\mathrm{Spec}(A)$ with values in $k$ correspond to the maximal ideals $M$ whose residue field $\kappa(M)$ is $k$; if $k$ is algebraically closed and $A$ finitely generated, these are all the maximal ideals.}
\entree[prop]{28083}{7:48:03}{For $A$ finitely generated over a field $k$, one has $X(k)\subset X_0\subset\mathrm{Spec}(A)$, where $X_0$ is the set of maximal ideals, with $X(k)=X_0$ if $k$ is algebraically closed.}
\topic{Very dense subsets}
\entree[def]{28151}{7:49:11}{A subset $X_0$ of $X$ is very dense if $U\mapsto U\cap X_0$, from the open sets of $X$ to those of $X_0$, is injective (hence bijective).}
\entree[prop]{28339}{7:52:19}{A subset $X_0$ is very dense if and only if every nonempty locally closed subset of $X$ meets $X_0$.}
\entree[prop]{28403}{7:53:23}{A subset $X_0$ is very dense if and only if, for every closed set $T$ of $X$, $T\cap X_0$ is dense in $T$.}
\entree[prop]{28472}{7:54:32}{If $X_0$ is very dense in $X$, the lattices of open sets are isomorphic and $\mathrm{Sh}(X)\to\mathrm{Sh}(X_0)$ is an equivalence of categories.}
\entree[prop]{28537}{7:55:37}{Conversely, if the inverse image of sheaves from $X$ to $X_0$ is faithful or conservative, then $X_0$ is very dense in $X$.}
\entree[prop]{28667}{7:57:47}{If $X$ is sober and $X_0$ very dense, $X$ is the sober space associated to $X_0$ (points: irreducible closed sets of $X_0$), hence determined by $X_0$ up to homeomorphism.}
\entree[prop]{28733}{7:58:53}{For $A$ finitely generated over a field, the set of closed points of $\mathrm{Spec}(A)$ is very dense (a consequence of the Nullstellensatz).}
\entree[prop]{28988}{8:03:08}{For $A$ finitely generated over a field, the nilradical of $A$ is the intersection of its maximal ideals; for any ring it is the intersection of the prime ideals.}
\entree[prop]{29114}{8:05:14}{For $k$ algebraically closed, $X(k)$, with the induced topology, has the same open sets, closed sets, locally closed sets, irreducible closed sets and components as $\mathrm{Spec}(A)$, but lacks the generic points.}
\face{Tape 2}
\entree[rem]{29287}{8:08:07}{The topology of $\mathrm{Spec}(A)$ does not suffice to reconstruct $A$; one must equip $\mathrm{Spec}(A)$ with a sheaf of rings.}
\session{11 June 1973}
\noindent\textit{Video 1/4 (\href{https://www.youtube.com/watch?v=51wmmpHand0}{\texttt{51wmmpHand0}}): tape 2, cont. (8:10:17--8:56:07); tape 1 (8:56:07--9:45:14); tape 2 (9:45:14--10:33:10).}\par\medskip
\noindent\textit{Video 1/2 of the groups course (\href{https://www.youtube.com/watch?v=_lTf_L0yvrU}{\texttt{\_lTf\_L0yvrU}}): tape 2, from 3:45:53 (3:45:53--3:47:24); tape 1 (3:47:24--4:34:37); tape 2 (4:34:37--5:03:20).}\par\medskip
\video{51wmmpHand0}
\face{Tape 2, cont.}
\topic{Group objects and Hopf algebras}
\entree[def]{29480}{8:11:20}{A group in a category embedded in functors is an object $G$ equipped with a group structure on each $G(Z)=\mathrm{Hom}(Z,G)$, functorial in $Z$; this requires neither products nor a final object.}
\entree[prop]{29738}{8:15:38}{A reduced affine group of finite type over an algebraically closed field is smooth (its ring is regular); if it is connected, its function field is a purely transcendental extension.}
\entree[disc]{29804}{8:16:44}{A question on where the group law lives for $GL_n$: it sits on the points with values in every algebra, functorially, so on the ring it appears as a co-group structure.}
\entree[ex]{29867}{8:17:47}{For $A=K[T_{ij}][\Delta^{-1}]$, $\Delta=\det(T_{ij})$, the affine space represented has as points with values in $K'$ the $n\times n$ matrices with invertible determinant: it is $\mathrm{GL}_n$.}
\entree[prop]{29996}{8:19:56}{A group structure on the affine space with ring $A$ is equivalent to a comultiplication $A\to A\otimes_K A$ (with counit) satisfying axioms; $A$ is a bialgebra.}
\entree[prop]{30063}{8:21:03}{Inversion $g\mapsto g^{-1}$ corresponds to an automorphism of the $K$-algebra $A$, the antipode.}
\topic{Complements on the spectrum}
\entree[prop]{30192}{8:23:12}{The $V(J)$ are the closed sets of $\mathrm{Spec}(A)$, which is quasi-compact; the closed sets correspond to radical ideals.}
\entree[prop]{30258}{8:24:18}{Prime ideals correspond to irreducible closed sets and to points; $\mathrm{Spec}(A)$ is sober, and empty if and only if $A=0$.}
\entree[prop]{30323}{8:25:23}{The space $\mathrm{Spec}(A)$ is irreducible if and only if the radical of $A$ is prime, that is, $A/\mathrm{rad}(A)$ is an integral domain.}
\entree[prop]{30451}{8:27:31}{A space is noetherian if and only if (a) every closed set has finitely many irreducible components and (b) every decreasing sequence of irreducible closed sets is stationary.}
\entree[prop]{30645}{8:30:45}{For $A$ finitely generated over a field $K$, the inclusion of the set $X_0$ of closed points into $\mathrm{Spec}(A)$ is a quasi-homeomorphism ($X_0$ very dense); if $K$ is algebraically closed, $X_0=X(K)$.}
\entree[prop]{30897}{8:34:57}{The irreducible components of an open set $U$ of a space are the intersections with $U$ of the components of the space that meet $U$.}
\entree[prop]{30897}{8:34:57}{If $X$ has finitely many irreducible components, $X$ is locally connected; $x,y$ are in the same connected component if and only if a chain of irreducible components joins them, two consecutive ones meeting.}
\topic{Functoriality of Spec and basic open sets}
\entree[prop]{31421}{8:43:41}{A homomorphism $u\colon A\to B$ induces $\varphi\colon\mathrm{Spec}(B)\to\mathrm{Spec}(A)$, $\mathfrak p\mapsto u^{-1}(\mathfrak p)$, which is prime because $A/u^{-1}(\mathfrak p)\subset B/\mathfrak p$ is an integral domain.}
\entree[prop]{31557}{8:45:57}{One has $\varphi^{-1}(V(J))=V(JB)$; the map $\varphi$ is therefore continuous.}
\entree[prop]{31622}{8:47:02}{The map $A\mapsto\mathrm{Spec}(A)$ is a contravariant functor from $K$-algebras to topological spaces.}
\entree[prop]{31688}{8:48:08}{If $u\colon A\to B\simeq A/J$ is surjective, $\varphi$ is a homeomorphism of $\mathrm{Spec}(B)$ onto the closed set $V(J)$: a closed immersion.}
\entree[prop]{31819}{8:50:19}{The map $A\to A_f$ is an epimorphism of rings; the homomorphisms $A_f\to K'$ are those of $A$ that send $f$ to a unit, and $\mathrm{Spec}(A_f)\to\mathrm{Spec}(A)$ is injective.}
\entree[prop]{31882}{8:51:22}{The map $\mathrm{Spec}(A_f)\to\mathrm{Spec}(A)$ induces isomorphisms between residue fields at corresponding points.}
\entree[prop]{31945}{8:52:25}{Its image is $Y_f=Y\setminus V(fA)$, the set of prime ideals not containing $f$, and it is a homeomorphism onto its image.}
\entree[prop]{32072}{8:54:32}{The open sets $Y_f$, $f\in A$, form a base of the topology of $Y=\mathrm{Spec}(A)$.}
\face{Tape 1 (8:56:07--9:45:14)}
\entree[prop]{32243}{8:57:23}{One has $Y_1=Y$ and $Y_f\cap Y_g=Y_{fg}$; the open sets $Y_f$ are stable under finite intersections.}
\entree[prop]{32306}{8:58:26}{A map from a space to $\mathrm{Spec}(A)$ is continuous if and only if the preimage of each $Y_f$ is open.}
\entree[prop]{32371}{8:59:31}{For $\varphi\colon\mathrm{Spec}(B)\to\mathrm{Spec}(A)$ induced by $u$, one has $\varphi^{-1}(Y_f)=X_{u(f)}$.}
\entree[prop]{32699}{9:04:59}{The map $\mathrm{Spec}(A_f)\to\mathrm{Spec}(A)$ is an open embedding with image $Y_f$: the basic open sets $Y_f$ are identified with the spectra $\mathrm{Spec}(A_f)$.}
\entree[prop]{32774}{9:06:14}{If $A=\prod_i A_i$ (finite family), $\mathrm{Spec}(A)$ is the topological sum of the $\mathrm{Spec}(A_i)$, each being open and closed.}
\topic{Spec and fibre products}
\entree[prop]{33240}{9:14:00}{The canonical map $\mathrm{Spec}(A\otimes_K B)\to\mathrm{Spec}(A)\times_{\mathrm{Spec}(K)}\mathrm{Spec}(B)$ is surjective, but almost never injective: $\mathrm{Spec}$ does not commute with fibre products.}
\entree[prop]{33308}{9:15:08}{If $x\mapsto y$ under $\mathrm{Spec}(u)$, the injection $A/\mathfrak q\to B/\mathfrak p$ induces an inclusion $\kappa(y)\to\kappa(x)$ of residue fields.}
\entree[prop]{33500}{9:18:20}{The fibre of this map over $(x,y)$, with $x,y$ over $s$, is in bijection with $\mathrm{Spec}\bigl(\kappa(x)\otimes_{\kappa(s)}\kappa(y)\bigr)$.}
\entree[ex]{33562}{9:19:22}{If $\kappa(x)=\kappa(s)$ or $\kappa(y)=\kappa(s)$, the fibre is a point; if $\kappa(x)$ and $\kappa(y)$ are transcendental over $\kappa(s)$, it is infinite.}
\entree[ex]{34080}{9:28:00}{For $K$ algebraically closed, $\mathrm{Spec}\,K[T]$ consists of the maximal ideals $(T-a)$, $a\in K$, and the generic point $\eta=(0)$.}
\entree[ex]{34212}{9:30:12}{The closed sets of $\mathrm{Spec}\,K[T]$ are the finite subsets of $K$ and the whole space; the nonempty open sets are the complements of finite subsets of $K$.}
\entree[ex]{34277}{9:31:17}{For $\mathrm{Spec}\,K[U,V]\to\mathrm{Spec}\,K[U]\times\mathrm{Spec}\,K[V]$, the fibre over any point of $K\times K$, $K\times\{\eta\}$ or $\{\eta\}\times K$ is a single point.}
\entree[ex]{34542}{9:35:42}{The fibre over $(\eta,\eta)$ is $\mathrm{Spec}(K(U)\otimes_K K(V))$: the whole plane and the irreducible curves $f=0$ that are neither of the form $U=\lambda$ nor $V=\lambda$.}
\entree[prop]{34737}{9:38:57}{For a continuous map $\varphi$, the closure of $\varphi(x)$ is the closure of the image of the closure of $x$.}
\entree[cex]{34866}{9:41:06}{Continuous maps between spectra are not characterized as coming from ring homomorphisms.}
\entree[cex]{34995}{9:43:15}{The topology of $\mathrm{Spec}(A\otimes_K B)$ is not the preimage of the product topology of $\mathrm{Spec}(A)\times\mathrm{Spec}(B)$.}
\face{Tape 2 (9:45:14--10:33:10)}
\entree[prop]{35116}{9:45:16}{The $K$-points of $X\times Y$ are $X(K)\times Y(K)$; the spectral topology of the product, restricted to these points, is finer than the product topology.}
\topic{Faithfully flat descent and \v{C}ech cohomology}
\entree[thm]{35250}{9:47:30}{If $A\to A'$ is faithfully flat, $0\to M\to M\otimes_A A'\rightrightarrows M\otimes_A A'\otimes_A A'$ is exact: $M$ is the submodule of the elements of $M\otimes_A A'$ having the same image under the two base changes (descent).}
\entree[prop]{35570}{9:52:50}{If $X'\to X$ admits finite cartesian powers $X'^I$ over $X$, one obtains an augmented simplicial object of $\mathcal{C}/X$ (index $I$ a nonempty finite set, augmentation $X$).}
\entree[prop]{36021}{10:00:21}{For a contravariant functor $F$ with values in abelian groups, the alternating sum of the faces gives a complex; its cohomology is the \v{C}ech cohomology $\check H^i(X'/X,F)$, with an augmentation $F(X)\to\check H^0$.}
\entree[def]{36152}{10:02:32}{If $X'\to X$ admits a section, the augmentation is a homotopy equivalence of simplicial objects: $\check H^i(X'/X,F)=0$ for $i\geq1$ and $F(X)\simeq\check H^0$.}
\entree[prop]{36474}{10:07:54}{If $A\to A'$ admits a retraction $A'\to A$, then $H^i(A'/A,M)=0$ for $i>0$ and $H^0=M$, where $F_M(B)=M\otimes_A B$.}
\entree[prop]{36607}{10:10:07}{For $B'=A'\otimes_A B$, the \v{C}ech complex of $B'/B$ with coefficients in $M_B$ is that of $A'/A$ with coefficients in $M$ tensored by $B$ over $A$.}
\entree[prop]{36737}{10:12:17}{If $B$ is flat over $A$, $H^i(B'/B,M_B)=H^i(A'/A,M)\otimes_A B$; if $B$ is faithfully flat, the exactness of the augmentation can be checked after extension to $B$.}
\entree[thm]{36802}{10:13:22}{If $A'$ is faithfully flat over $A$, then $H^i(A'/A,M)=0$ for $i>0$ and $M\to H^0(A'/A,M)$ is an isomorphism (take $B=A'$, where the diagonal is split).}
\entree[prop]{36939}{10:15:39}{If the $f_i$ generate the unit ideal of $A$, the \v{C}ech complex of the covering of $\mathrm{Spec}(A)$ by the $X_{f_i}$, with coefficients in $M_f$ on $X_f$, is exact in degrees $>0$ and has $H^0=M$.}
\entree[prop]{37068}{10:17:48}{Two morphisms $X'\to X''$ over $X$ induce homotopic morphisms between the simplicial complexes; if $X'$ and $X''$ dominate each other, they have the same \v{C}ech cohomology.}
\entree[prop]{37263}{10:21:03}{If $f_1$ to $f_n$ generate the unit ideal of $A$, the algebra $A\to\prod_i A_{f_i}$ is faithfully flat.}
\entree[prop]{37457}{10:24:17}{The presheaf $X_f\mapsto M_f$ is a sheaf $\widetilde{M}$ on $\mathrm{Spec}(A)$, and $H^i(\mathrm{Spec}(A),\widetilde{M})=0$ for $i>0$: an affine scheme is acyclic for $\widetilde{M}$.}
\entree[prop]{37586}{10:26:26}{The diagonal $\mathrm{Spec}(A)\to\mathrm{Spec}(A)\times_{\mathrm{Spec}(K)}\mathrm{Spec}(A)$ is a closed immersion, because $A\otimes_K A\to A$, $x\otimes y\mapsto xy$, is surjective.}
\entree[prop]{37650}{10:27:30}{A module $M$ of finite presentation is flat if and only if it is finitely generated projective, if and only if $\widetilde{M}$ is locally free on $\mathrm{Spec}(A)$.}
\entree[disc]{37650}{10:27:30}{Asked about the geometric meaning of flat modules, the lecturer explains that finitely presented flat modules are projective, hence locally free sheaves, a version of vector bundles.}
\entree[rem]{37782}{10:29:42}{To $X$ affine of finite type over $\mathbb{C}$ is associated the analytic space $X^{\mathrm{an}}$, with a morphism of ringed spaces $X^{\mathrm{an}}\to X$ universal for morphisms of analytic spaces to $X$.}
\entree[prop]{37848}{10:30:48}{For $M$ finitely generated over $A$ finitely generated over $\mathbb{C}$, $M$ is flat if and only if the inverse image of $\widetilde{M}$ on $X^{\mathrm{an}}$ is locally free, i.e. a holomorphic vector bundle.}
\video{_lTf_L0yvrU}
\face{Tape 2, from 3:45:53}
\noindent\textit{Recorded on video 1/2 of the groups course (\href{https://www.youtube.com/watch?v=_lTf_L0yvrU}{\texttt{\_lTf\_L0yvrU}}); by its content it belongs to the algebraic geometry course.}\par\smallskip
\topic{Sheaves on a spectrum}
\entree[prop]{13553}{3:45:53}{The spectrum $X=\mathrm{Spec}\,A$ of a ring, equipped with its structure sheaf, is a locally ringed space.}
\entree[prop]{13617}{3:46:57}{The functor $A\mapsto\mathrm{Spec}\,A$ from rings to locally ringed spaces turns tensor products into fibre products: $\mathrm{Spec}(A\otimes_kB)=\mathrm{Spec}\,A\times_{\mathrm{Spec}\,k}\mathrm{Spec}\,B$.}
\face{Tape 1 (3:47:24--4:34:37)}
\noindent\textit{Recorded on video 1/2 of the groups course (\href{https://www.youtube.com/watch?v=_lTf_L0yvrU}{\texttt{\_lTf\_L0yvrU}}); by its content it belongs to the algebraic geometry course.}\par\smallskip
\entree[prop]{13646}{3:47:26}{A sheaf on $X=\mathrm{Spec}\,A$ is equivalent to a contravariant functor on the open sets $X_f$ satisfying the sheaf condition for coverings of an $X_f$ by $X_{f_i}$.}
\entree[prop]{13976}{3:52:56}{$X_f\subset X_g$ if and only if there exist an integer $n\ge0$ and $h\in A$ such that $f^n=gh$; this relation is transitive.}
\entree[prop]{14363}{3:59:23}{$X_f=\bigcup_iX_{f_i}$ if and only if $f_i^{n_i}\in fA$ for all $i$ and there exist $n$ and $g_i\in A$ (almost all zero) with $f^n=\sum_ig_if_i$.}
\topic{The sheaf associated to a module}
\entree[def]{14427}{4:00:27}{To an $A$-module $M$ one associates a sheaf $\widetilde{M}$ on $X=\mathrm{Spec}\,A$; for $M=A$, it is a sheaf of rings.}
\entree[def]{14490}{4:01:30}{One sets $\widetilde{M}(X_f)=M_f$, the localization of $M$ at $f$, which is the inductive limit of $M\xrightarrow{f}M\xrightarrow{f}M\to\cdots$.}
\entree[prop]{14557}{4:02:37}{In $M_f$, multiplication by $f$ is bijective, and $M\to M_f$ is the canonical map.}
\entree[prop]{14682}{4:04:42}{Universal property: for every $A$-module $N$ on which $f$ acts bijectively, every homomorphism $M\to N$ factors uniquely through $M_f$; in particular $M\mapsto M_f$ is a functor.}
\entree[prop]{14941}{4:09:01}{Description by fractions: $M_f$ is the set of $x/f^n$, with $x/f^n=x'/f^{n'}$ if and only if there exists $m$ such that $f^m(f^{n'}x-f^nx')=0$.}
\entree[prop]{15076}{4:11:16}{$A_f$ is a ring, $A\to A_f$ a ring homomorphism, and $\frac{x}{f^n}\cdot\frac{y}{f^{n'}}=\frac{xy}{f^{n+n'}}$.}
\entree[prop]{15142}{4:12:22}{$M_f$ is an $A_f$-module via $\frac{\lambda}{f^n}\cdot\frac{x}{f^m}=\frac{\lambda x}{f^{n+m}}$.}
\entree[prop]{15206}{4:13:26}{Universal property of $A_f$: every homomorphism from $A$ to a ring $R$ (even noncommutative) sending $f$ to a unit factors uniquely through $A_f$.}
\entree[prop]{15270}{4:14:30}{Giving an $A_f$-module is equivalent to giving an $A$-module on which $f$ acts bijectively.}
\entree[prop]{15575}{4:19:35}{If $X_f\subset X_g$, say $f^n=gh$, then $g$ is invertible in $M_f$, whence a restriction $M_g\to M_f$, functorial in $M$ and compatible with composition: $f\mapsto M_f$ is a presheaf.}
\entree[thm]{15707}{4:21:47}{$f\mapsto M_f$ is a sheaf on the base of open sets $X_f$, that is, $M_f\to\prod_iM_{f_i}\rightrightarrows\prod_{i,j}M_{f_if_j}$ is exact for every covering of $X_f$ by $X_{f_i}$.}
\entree[prop]{15842}{4:24:02}{$\widetilde{A}$ is a sheaf of rings on $X$ and $\widetilde{M}$ is a sheaf of modules over $\widetilde{A}$.}
\face{Tape 2 (4:34:37--5:03:20)}
\noindent\textit{Recorded on video 1/2 of the groups course (\href{https://www.youtube.com/watch?v=_lTf_L0yvrU}{\texttt{\_lTf\_L0yvrU}}); by its content it belongs to the algebraic geometry course.}\par\smallskip
\entree[prop]{16542}{4:35:42}{The open sets $X_f$ are quasi-compact: a covering of $X_f$ by $X_{f_i}$ admits a finite subcover, because $f^n=\sum_ig_if_i$ is a finite sum; the sheaf condition can be checked on this subcover.}
\entree[prop]{16989}{4:43:09}{One has $\widetilde{M}(X)=M_1=M$ and $\widetilde{M}(X_f)=M_f$: the localization $M_f$ is the module of sections of $\widetilde{M}$ over the open set $X_f$.}
\entree[prop]{17115}{4:45:15}{Localization is identified with a base change: the canonical map $M\otimes_AA_f\to M_f$ is an isomorphism.}
\entree[prop]{17311}{4:48:31}{If $f_1,\ldots,f_r$ generate the unit ideal and $A'=\prod_iA_{f_i}$, then $A'\otimes_AA'=\prod_{i,j}A_{f_if_j}$; the sheaf condition is equivalent to the exactness of $0\to M\to M\otimes_AA'\rightrightarrows M\otimes_AA'\otimes_AA'$.}
\topic{Flat and faithfully flat modules}
\entree[thm]{17443}{4:50:43}{If $A'$ is a faithfully flat $A$-algebra, the sequence $0\to M\to M\otimes_AA'\rightrightarrows M\otimes_AA'\otimes_AA'$ is exact for every $A$-module $M$; it can fail to be exact without the hypothesis (for example for $A'=0$).}
\entree[def]{17573}{4:52:53}{An $A$-module $M$ is flat if the functor $N\mapsto N\otimes_AM$ is exact, that is (this functor being right exact) if it sends injections to injections.}
\entree[def]{17766}{4:56:06}{$M$ is faithfully flat if it is flat and $N\mapsto M\otimes_AN$ is faithful; for this exact functor, that is equivalent to: $M\otimes_AN=0$ implies $N=0$.}
\entree[prop]{17830}{4:57:10}{If $M$ is faithfully flat, a homomorphism $N\to N'$ is injective (resp. surjective, bijective) if and only if $M\otimes N\to M\otimes N'$ is.}
\entree[def]{17894}{4:58:14}{An $A$-algebra $A'$ is faithfully flat if it is so as a module; extension of scalars from $A$-modules to $A'$-modules is then exact and faithful.}
\entree[prop]{17958}{4:59:18}{For $A'$ faithfully flat over $A$, the complex $0\to M\to M\otimes_AA'\to M\otimes_AA'\otimes_AA'\to\cdots$ is exact: $H^0(A'/A,M)=M$ and $H^i(A'/A,M)=0$ for $i>0$.}
\session{12 June 1973}
\noindent\textit{Video 2/4 (\href{https://www.youtube.com/watch?v=thSz9ZHDK9o}{\texttt{thSz9ZHDK9o}}): tape 1 (0:00:00--0:47:50); tape 2 (0:47:50--1:36:26); tape 2 (1:36:26--2:09:27).}\par\medskip
\video{thSz9ZHDK9o}
\face{Tape 1}
\entree[def]{66}{0:01:06}{A module $M$ is flat if the functor $-\otimes M$ is exact.}
\entree[prop]{129}{0:02:09}{The diagonal of an affine algebraic space $X=\mathrm{Spec}\,A$ over $K$ is a closed immersion, because $A\otimes_K A\to A$, $x\otimes y\mapsto xy$, is surjective.}
\topic{Disjoint and universal sums; constant objects}
\entree[prop]{336}{0:05:36}{Finite sums of affine algebraic spaces over $K$ correspond to finite products of $K$-algebras.}
\entree[def]{401}{0:06:41}{A sum $X=\coprod X_i$ is disjoint if the $u_i\colon X_i\to X$ are monomorphisms and $X_i\times_X X_j$ is an initial object for $i\neq j$.}
\entree[def]{464}{0:07:44}{A sum $X=\coprod X_i$ is universal if, for every base change $X'\to X$, $\coprod X'_i\to X'$ is an isomorphism, $X'_i=X_i\times_X X'$.}
\entree[prop]{527}{0:08:47}{A morphism $u_i$ is a monomorphism if and only if $X_i\times_X X_i\simeq X_i$ via the diagonal.}
\entree[prop]{593}{0:09:53}{In the category of affine algebraic spaces over $K$, finite sums are disjoint and universal.}
\entree[prop]{658}{0:10:58}{For $A=\prod_i A_i$, the $X_i\to X$ are closed immersions ($A\to A_i$ surjective) and $A_i\otimes_A A_j=0$ for $i\neq j$.}
\entree[prop]{918}{0:15:18}{If sums are universal, $(\coprod_i X_i)\times(\coprod_j Y_j)\simeq\coprod_{i,j}X_i\times Y_j$ (distributivity of the product over the sum).}
\entree[def]{1176}{0:19:36}{For a finite set $I$, the constant object $I_C=\coprod_I e$ is the sum of $I$ copies of the final object $e$.}
\entree[prop]{1239}{0:20:39}{If finite sums are disjoint and universal, $I_C\times J_C\simeq(I\times J)_C$; the functor $I\mapsto I_C$ commutes with fibre products and the final object, hence is left exact.}
\entree[prop]{1369}{0:22:49}{If $I$ is a finite group, $I_K$ is an affine algebraic group over $K$ (constant group); more generally $I_K$ inherits any structure defined by finite limits on $I$.}
\entree[prop]{1437}{0:23:57}{A functor between categories with finite projective limits is left exact if and only if it commutes with fibre products and sends final object to final object.}
\entree[prop]{1502}{0:25:02}{One has $I_K=\mathrm{Spec}\,K^I$, where $K^I$ is the $K$-algebra of maps from $I$ to $K$.}
\entree[prop]{1568}{0:26:08}{The law $I\times I\to I$ induces on $K^I$ the comultiplication $K^I\to K^{I\times I}\simeq K^I\otimes_K K^I$ transpose to the law.}
\topic{Monomorphisms, fibres and the spectrum}
\entree[prop]{1836}{0:30:36}{If $Y\to S$ is a monomorphism of affine algebraic spaces, $\mathrm{Spec}\,Y\to\mathrm{Spec}\,S$ is injective.}
\entree[prop]{1899}{0:31:39}{If $Y\to S$ is a monomorphism, $\mathrm{Spec}(X\times_S Y)$ is, in $\mathrm{Spec}\,X$, the preimage of $\mathrm{Spec}\,Y$ under $\mathrm{Spec}\,X\to\mathrm{Spec}\,S$.}
\entree[prop]{2093}{0:34:53}{If $Y\to S$ is a closed immersion ($A\to C$ surjective), so is $Z=X\times_S Y\to X$, and the topology of $\mathrm{Spec}\,Z$ is induced by that of $\mathrm{Spec}\,X$.}
\entree[prop]{2156}{0:35:56}{If $Y=S_f$ ($A\to A_f$), then $B\otimes_A A_f=B_f$ and the topology of $\mathrm{Spec}\,Z$ is induced by that of $\mathrm{Spec}\,X$.}
\entree[prop]{2288}{0:38:08}{If $A\to C$ is a quotient $A/J$ followed by a localization, the topology of $\mathrm{Spec}(B\otimes_A C)$ is induced by $\mathrm{Spec}\,B$ and $\mathrm{Spec}$ commutes with this fibre product.}
\entree[prop]{2416}{0:40:16}{For two subobjects $X$, $Y$ of $S$, one has $\mathrm{Spec}(X\cap Y)=\mathrm{Spec}\,X\cap\mathrm{Spec}\,Y$ in $\mathrm{Spec}\,S$ (set-theoretically).}
\entree[prop]{2613}{0:43:33}{For a prime ideal $\mathfrak p$ of $A$ and $K'=\kappa(\mathfrak p)$, the fibre $X\times_S\mathrm{Spec}\,K'$ has as spectrum the fibre of $\mathrm{Spec}\,X\to\mathrm{Spec}\,S$ at $\mathfrak p$, with the induced topology.}
\entree[prop]{2810}{0:46:50}{For $u\colon A\to B$ and $\varphi=\mathrm{Spec}\,u$, the closure of $\varphi(V(J))$ is $V(u^{-1}(J))$, for every ideal $J$ of $B$.}
\face{Tape 2 (0:47:50--1:36:26)}
\entree[prop]{3134}{0:52:14}{The morphism $\varphi=\mathrm{Spec}\,u$ is dominant if and only if $\mathrm{Ker}\,u$ is contained in the nilradical of $A$; if $A$ is reduced, if and only if $u$ is injective.}
\topic{The structure sheaf and quasi-coherent modules}
\entree[prop]{3405}{0:56:45}{For an $A$-module $M$, $X_f\mapsto M_f$ defines a sheaf $\widetilde{M}$ of $A$-modules on $X=\mathrm{Spec}\,A$.}
\entree[prop]{3469}{0:57:49}{$\widetilde{A}$ is a sheaf of $A$-algebras with $\Gamma(X,\widetilde{A})=A$; $\widetilde{M}$ is a sheaf of $\widetilde{A}$-modules.}
\entree[def]{3533}{0:58:53}{The sheaf of rings $\mathcal{O}_X=\widetilde{A}$ is the structure sheaf of $X=\mathrm{Spec}\,A$, which thus becomes a ringed space; the ring $A$ is recovered as $\Gamma(X,\mathcal{O}_X)$.}
\entree[prop]{3728}{1:02:08}{The functor $M\mapsto\widetilde{M}$ is exact and commutes with arbitrary inductive limits and with the tensor product.}
\entree[prop]{3792}{1:03:12}{If $M$ is finitely presented, $\widetilde{\mathrm{Hom}_A(M,N)}\simeq\underline{\mathrm{Hom}}_{\mathcal{O}_X}(\widetilde{M},\widetilde{N})$; for $A$ noetherian, finitely presented is equivalent to finitely generated.}
\entree[prop]{3856}{1:04:16}{The functor $M\mapsto\widetilde{M}$ is fully faithful: $\mathrm{Hom}_A(M,N)\simeq\mathrm{Hom}_{\mathcal{O}_X}(\widetilde{M},\widetilde{N})$.}
\entree[prop]{4112}{1:08:32}{One has $\widetilde{M}\otimes_{\mathcal{O}_X}\widetilde{N}\simeq\widetilde{M\otimes_A N}$.}
\entree[prop]{4308}{1:11:48}{One has $\widetilde{M}\simeq M\otimes_A\mathcal{O}_X$: it is the inverse image of $M$ under the morphism of ringed spaces $(X,\mathcal{O}_X)\to(\mathrm{pt},A)$.}
\entree[def]{4703}{1:18:23}{An $\mathcal{O}_X$-module $\mathcal{M}$ is quasi-coherent if there exists an open covering $(X_i)$ such that $\mathcal{M}|_{X_i}$ is the cokernel of a morphism $\mathcal{O}_{X_i}^{(J_1)}\to\mathcal{O}_{X_i}^{(J)}$, with arbitrary index sets.}
\entree[thm]{4841}{1:20:41}{On $X=\mathrm{Spec}\,A$, an $\mathcal{O}_X$-module is quasi-coherent if and only if it is isomorphic to some $\widetilde{M}$.}
\entree[prop]{4903}{1:21:43}{Every $\widetilde{M}$ is the cokernel of a morphism of free $\mathcal{O}_X$-modules $\mathcal{O}_X^{(J_1)}\to\mathcal{O}_X^{(J)}$; conversely such a cokernel is of the form $\widetilde{M}$.}
\topic{Gluing and descent}
\entree[prop]{5103}{1:25:03}{If $X=\bigcup_i X_{f_i}$ (finite family), modules $M_i$ over $A_{f_i}$ with gluing isomorphisms on the $A_{f_if_j}$ satisfying the cocycle condition come from an $A$-module $M$.}
\entree[prop]{5165}{1:26:05}{For $A'=\prod_i A_{f_i}$, one has $A'\otimes_A A'\simeq\prod_{i,j}A_{f_if_j}$; the gluing datum is a descent datum from $A'$ to $A$.}
\entree[prop]{5297}{1:28:17}{If $X=\bigcup_i X_{f_i}$ (finite family), $A\to\prod_i A_{f_i}$ is faithfully flat.}
\entree[def]{5551}{1:32:31}{$X'\to X$ is of effective descent for a fibred category $F$ if $F(X)\to D_{X'/X}(F(X'))$ (objects with descent datum) is an equivalence; faithful or fully faithful are weaker cases.}
\entree[prop]{5614}{1:33:34}{If $X'\to X$ admits a section, it is a morphism of effective descent for every fibred category.}
\entree[prop]{5680}{1:34:40}{If $A\to A'$ is faithfully flat, $\mathrm{Spec}\,A'\to\mathrm{Spec}\,A$ is of effective descent for modules: an $A$-module is equivalent to an $A'$-module equipped with a descent datum.}
\face{Tape 2 (1:36:26--2:09:27)}
\entree[thm]{5791}{1:36:31}{By faithfully flat descent, $M\mapsto\widetilde{M}$ is an equivalence between $A$-modules and quasi-coherent $\mathcal{O}_X$-modules on $X=\mathrm{Spec}\,A$.}
\topic{Stalks and localization}
\entree[prop]{5989}{1:39:49}{The fibre of $\widetilde{M}$ at the point $x$ associated to $\mathfrak p$ is $M_{\mathfrak p}=\varinjlim_{f\notin\mathfrak p}M_f$ (filtered system).}
\entree[prop]{6117}{1:41:57}{The fibre of $\mathcal{O}_X$ at $x$ is $\varinjlim_{f\notin\mathfrak p}A_f=A_{\mathfrak p}$.}
\entree[def]{6181}{1:43:01}{For a multiplicative subset $T\subset A$, $MT^{-1}=\varinjlim_{f\in T}M_f$; it is a module over the algebra $AT^{-1}=\varinjlim_{f\in T}A_f$; for $T=\{f^n\}$ one recovers $M_f$.}
\entree[prop]{6246}{1:44:06}{$M\to MT^{-1}$ is universal among the $A$-linear maps $M\to N$ such that multiplication by each $f\in T$ is bijective on $N$.}
\entree[prop]{6375}{1:46:15}{$A\to AT^{-1}$ is universal among the homomorphisms $A\to B$ that make the elements of $T$ invertible.}
\entree[prop]{6568}{1:49:28}{One has $MT^{-1}\simeq M\otimes_A AT^{-1}$.}
\entree[prop]{6632}{1:50:32}{The ring $A_{\mathfrak p}$ is local, with residue field $\kappa(\mathfrak p)=\mathrm{Frac}(A/\mathfrak p)$.}
\topic{Local rings}
\entree[def]{6702}{1:51:42}{A ring $O$ is local if it has a unique maximal ideal $\mathfrak m$; the field $O/\mathfrak m$ is its residue field.}
\entree[prop]{6764}{1:52:44}{An element is invertible if and only if it belongs to no maximal ideal; in a local ring, $O^*=O\setminus\mathfrak m$.}
\entree[prop]{6963}{1:56:03}{If $\mathfrak m$ is an ideal of $O$ with $O^*=O\setminus\mathfrak m$, then $O$ is local with maximal ideal $\mathfrak m$.}
\entree[prop]{7027}{1:57:07}{A ring $O\neq0$ is local if and only if, for every $f\in O$, $f$ or $1-f$ is invertible.}
\entree[prop]{7093}{1:58:13}{The maximal ideal of $A_{\mathfrak p}$ is $\mathfrak p A_{\mathfrak p}$, and $A_{\mathfrak p}/\mathfrak p A_{\mathfrak p}\simeq\kappa(\mathfrak p)$; thus $A_{\mathfrak p}^*=A_{\mathfrak p}\setminus\mathfrak pA_{\mathfrak p}$.}
\entree[def]{7631}{2:07:11}{A ringed space is locally ringed if all its fibres of $\mathcal{O}_X$ are local rings; $\mathrm{Spec}\,A$ is locally ringed.}
\session{13 June 1973}
\noindent\textit{Video 2/4 (\href{https://www.youtube.com/watch?v=thSz9ZHDK9o}{\texttt{thSz9ZHDK9o}}): tape 2, cont. (2:09:27--2:25:35); tape 1 (2:25:35--3:14:52); tape 2 (3:14:52--4:03:12).}\par\medskip
\video{thSz9ZHDK9o}
\face{Tape 2, cont.}
\entree[cex]{7837}{2:10:37}{The functor $M\mapsto\widetilde{M}$ does not commute with infinite products in general: $(\prod_i M_i)^{\sim}\to\prod_i\widetilde{M_i}$ is not always an isomorphism.}
\entree[cex]{8093}{2:14:53}{If $A$ is an integral domain and not a field, $f\neq0$ is a non-unit and $I$ is infinite, $(A^I)_f\to(A_f)^I$ is not surjective.}
\entree[cex]{8093}{2:14:53}{For $M=A^{(I)}$ with $I$ infinite, $\widetilde{\mathrm{Hom}_A(M,N)}\to\underline{\mathrm{Hom}}(\widetilde{M},\widetilde{N})$ is not an isomorphism in general.}
\entree[prop]{8227}{2:17:07}{For $O\neq0$: $O$ is local if and only if the non-units form an ideal, if and only if $f+g=1$ implies $f$ or $g$ is invertible.}
\entree[prop]{8624}{2:23:44}{If $A$ is local noetherian, $\hat A$ is faithfully flat over $A$; an $\hat A$-module equipped with a descent datum relative to $A\to\hat A$ comes from an $A$-module $M$, unique up to canonical isomorphism.}
\entree[prop]{8624}{2:23:44}{For $A\to\hat A$ faithfully flat, $M$ is finitely generated if and only if $M\otimes_A\hat A$ is; then $M\otimes_A\hat A$ is the completion of $M$.}
\face{Tape 1}
\topic{Locally ringed spaces and morphisms of spectra}
\entree[read]{8802}{2:26:42}{The yoga of descent is in his Séminaire Bourbaki talks. \cite{fga}}
\entree[def]{8874}{2:27:54}{For $f\in\mathcal{O}_X(U)$, one sets $U_f=\{x\in U : f(x)\neq0\}$, where $f(x)$ is the image of $f_x$ in $\kappa(x)$; it is the set of $x$ where $f_x$ is invertible.}
\entree[prop]{9004}{2:30:04}{A ringed space with nonzero fibres is locally ringed if and only if, for $f_1,\dots,f_n\in\mathcal{O}_X(U)$ with $\sum f_i=1$, the $U_{f_i}$ cover $U$.}
\entree[prop]{9067}{2:31:07}{For $u\colon A\to B$, $\varphi\colon\mathrm{Spec}\,B\to\mathrm{Spec}\,A$, $\mathfrak p\mapsto u^{-1}(\mathfrak p)$, is continuous.}
\entree[prop]{9194}{2:33:14}{The homomorphism $u$ defines $\mathcal{O}_Y\to\varphi_*\mathcal{O}_X$, given on $Y_f$ by $A_f\to B_{u(f)}$, since $\varphi^{-1}(Y_f)=X_{u(f)}$; $\varphi$ is a morphism of ringed spaces.}
\entree[prop]{9257}{2:34:17}{For an $A$-algebra $B$ and $f\in A$, the localization of $B$ at the image of $f$ coincides with the localization of the $A$-module $B$ at $f$.}
\entree[prop]{9386}{2:36:26}{One recovers $u$ as $\Gamma(\varphi^{\#})\colon\Gamma(Y,\mathcal{O}_Y)\to\Gamma(X,\mathcal{O}_X)$; the map $\mathrm{Hom}(A,B)\to\mathrm{Hom}(\mathrm{Spec}\,B,\mathrm{Spec}\,A)$ is injective.}
\entree[prop]{9513}{2:38:33}{For a $B$-module $N$, one has $\varphi_*\widetilde{N}\simeq\widetilde{N_{[A]}}$: direct image corresponds to restriction of scalars.}
\entree[prop]{9833}{2:43:53}{For $M$ an $A$-module and $N$ a $B$-module, $\mathrm{Hom}_{\mathcal{O}_Y}(\widetilde{M},\varphi_*\widetilde{N})\simeq\mathrm{Hom}_A(M,N_{[A]})\simeq\mathrm{Hom}_B(M\otimes_A B,N)$.}
\entree[prop]{9961}{2:46:01}{One has $\varphi^*\widetilde{M}\simeq(M\otimes_A B)^{\sim}$: inverse image corresponds to extension of scalars.}
\entree[def]{10229}{2:50:29}{A morphism of ringed spaces induces homomorphisms $\mathcal{O}_{Y,\varphi(x)}\to\mathcal{O}_{X,x}$; definition: $u\colon O\to O'$ between local rings is local if $u^{-1}(\mathfrak m')=\mathfrak m$.}
\entree[prop]{10435}{2:53:55}{For $u\colon O\to O'$ between local rings, one always has $u^{-1}(\mathfrak m')\subset\mathfrak m$; $u$ is local if and only if $u(\mathfrak m)\subset\mathfrak m'$, if and only if $f\in O^*\Leftrightarrow u(f)\in O'^*$.}
\entree[def]{10502}{2:55:02}{A morphism of locally ringed spaces is a morphism of ringed spaces all of whose induced homomorphisms on fibres are local.}
\entree[prop]{10708}{2:58:28}{For a morphism of ringed spaces, $\varphi^{-1}(U_f)\subset(X_{\varphi^{-1}U})_{\varphi^*(f)}$; it is local if and only if this is always an equality, whence $\varphi^{-1}(Y_f)=X_{\varphi^*(f)}$.}
\face{Tape 2}
\entree[disc]{11890}{3:18:10}{A confusion about morphisms of locally ringed spaces being automatically local is resolved by a counterexample from a discrete valuation ring and its fraction field, with the point sent to the closed point.}
\entree[cex]{12148}{3:22:28}{For $V$ a discrete valuation ring with fraction field $K$, the morphism of ringed spaces $\mathrm{Spec}\,K\to\mathrm{Spec}\,V$ at the closed point, with $V\subset K$, is not local and does not come from any ring homomorphism.}
\entree[prop]{12408}{3:26:48}{A morphism of locally ringed spaces $\varphi\colon\mathrm{Spec}\,B\to\mathrm{Spec}\,A$ equals $\mathrm{Spec}\,u$, where $u=\Gamma(\varphi^{\#})$: one has $\varphi(\mathfrak p)=u^{-1}(\mathfrak p)$.}
\entree[prop]{12471}{3:27:51}{The map $\mathrm{Hom}(A,B)\to\mathrm{Hom}_{\mathrm{loc}}(\mathrm{Spec}\,B,\mathrm{Spec}\,A)$ is bijective.}
\topic{Schemes and their functor of points}
\entree[def]{12602}{3:30:02}{An affine $k$-scheme is a locally ringed space isomorphic to the spectrum of a $k$-algebra; they form a category equivalent to that of affine algebraic spaces over $k$.}
\entree[prop]{12665}{3:31:05}{The only global functions on complex projective space $\mathbb{P}^n$, $n\geq1$, are the constants; it is therefore not affine.}
\entree[def]{12727}{3:32:07}{A scheme is a locally ringed space $(X,\mathcal{O}_X)$ that is a union of open sets $X_i$ each of which is isomorphic to some $\mathrm{Spec}\,A_i$.}
\entree[prop]{12791}{3:33:11}{Giving a scheme $X$ a $k$-scheme structure is equivalent to giving a ring homomorphism $k\to\Gamma(X,\mathcal{O}_X)$.}
\entree[def]{13111}{3:38:31}{$X(k')=\mathrm{Hom}_k(\mathrm{Spec}\,k',X)$; for $X=\mathrm{Spec}\,A$, one recovers the set $\mathrm{Hom}_{k\text{-alg}}(A,k')$.}
\entree[prop]{13183}{3:39:43}{The functor associating to a $k$-scheme $X$ the functor $k'\mapsto X(k')$ is fully faithful from the category of $k$-schemes to that of functors from $k$-algebras to sets.}
\entree[prop]{13375}{3:42:55}{If $X=\mathrm{Spec}\,A$ is affine, $\mathrm{Hom}(X,Y)\to\mathrm{Hom}(X',Y')$ is bijective, since $\mathrm{Hom}(X',Y')=Y'(A)=Y(A)$ by Yoneda ($X'$ representable).}
\entree[prop]{13648}{3:47:28}{If $X=\bigcup X_i$ is an open covering, $\mathrm{Hom}(X,Y)\to\prod_i\mathrm{Hom}(X_i,Y)\rightrightarrows\prod_{i,j}\mathrm{Hom}(X_{ij},Y)$ is exact for every locally ringed space $Y$ ($X_{ij}=X_i\cap X_j$).}
\entree[prop]{14251}{3:57:31}{For $X$ a locally ringed space and $Y=\mathrm{Spec}\,A$, $\mathrm{Hom}_{\mathrm{loc}}(X,\mathrm{Spec}\,A)\simeq\mathrm{Hom}(A,\Gamma(X,\mathcal{O}_X))$.}
\session{18 June 1973}
\noindent\textit{Video 2/4 (\href{https://www.youtube.com/watch?v=thSz9ZHDK9o}{\texttt{thSz9ZHDK9o}}): tape 1 (4:03:12--4:51:20); tape 2, side A (4:51:20--5:39:53); tape 2, side B (5:39:53--6:28:25); tape 3 (6:28:25--7:00:11); tape 3, from 7:00:11 (7:00:11--7:17:49).}\par\medskip
\video{thSz9ZHDK9o}
\face{Tape 1}
\topic{Morphisms into an affine scheme}
\entree[prop]{14657}{4:04:17}{To a homomorphism $u\colon A\to\Gamma(X,\mathcal{O}_X)$, $X$ a locally ringed space, one associates $\varphi\colon X\to\mathrm{Spec}\,A$, $x\mapsto\ker(A\to\kappa(x))$, continuous, with $\varphi^{-1}(D(f))=X_{u(f)}$, the open set where $u(f)$ is invertible.}
\entree[prop]{14721}{4:05:21}{The homomorphisms $A_f\to\Gamma(X_{u(f)},\mathcal{O}_X)$ given by the universal property of $A_f$ complete $\varphi$ into a morphism of locally ringed spaces $X\to\mathrm{Spec}\,A$.}
\entree[thm]{14914}{4:08:34}{For every locally ringed space $X$, $\mathrm{Hom}(X,\mathrm{Spec}\,A)\simeq\mathrm{Hom}_{\mathrm{ring}}(A,\Gamma(X,\mathcal{O}_X))$, functorially: $\mathrm{Spec}$ is right adjoint to $X\mapsto\Gamma(X,\mathcal{O}_X)$ (opposite rings).}
\topic{Schemes as sheaves for the Zariski topology}
\entree[rem]{15108}{4:11:48}{Giving a scheme over a ring $k$ amounts to giving a scheme $X$ and a morphism $X\to\mathrm{Spec}\,k$; more generally one studies schemes over an arbitrary base scheme $S$.}
\entree[prop]{15240}{4:14:00}{The opposite of the category of $k$-algebras is equivalent to the category of affine schemes over $\mathrm{Spec}\,k$.}
\entree[prop]{15432}{4:17:12}{A $k$-scheme $X$ defines, by restricting $\mathrm{Hom}(-,X)$ to affine schemes, a functor from $k$-algebras to $\mathbf{Set}$; $X(k')$ is the set of points of $X$ with values in $k'$.}
\entree[prop]{15698}{4:21:38}{The functor represented by a scheme $X$ is a Zariski sheaf: for an open covering $Y=\bigcup Y_i$, $X(Y)\to\prod X(Y_i)\rightrightarrows\prod X(Y_{ij})$ is exact.}
\entree[prop]{15892}{4:24:52}{The families of open immersions $(Y_i\to Y)$ that are jointly surjective form a pretopology on the category of $k$-schemes: the Zariski topology.}
\entree[prop]{16149}{4:29:09}{Restriction to affine schemes is a fully faithful functor from Zariski sheaves on $k$-schemes to sheaves on affine $k$-schemes.}
\entree[prop]{16277}{4:31:17}{Every open subset of a scheme is a scheme; the affine open sets form a base of the topology (the $X_f$ do so for an affine scheme).}
\entree[prop]{16339}{4:32:19}{A sheaf $X$ is determined by its values on affines: $X(Y)$ is the projective limit of the $X(Y_{ijk})$, the $Y_{ijk}$ being affines covering the $Y_{ij}$ and $Y_i$.}
\entree[prop]{16404}{4:33:24}{The functor associating to a $k$-scheme its functor from $k$-algebras to $\mathbf{Set}$ is fully faithful: a scheme is determined up to isomorphism by its points with values in algebras.}
\topic{Sums of schemes}
\entree[prop]{16684}{4:38:04}{The disjoint topological sum of the $X_i$, equipped with the sheaf of rings inducing $\mathcal{O}_{X_i}$ on each $X_i$, is the sum of the $X_i$ in the categories of ringed spaces and of locally ringed spaces.}
\entree[prop]{16874}{4:41:14}{This sum is a scheme; it is therefore the sum of the $X_i$ in the category of schemes (and in that of schemes over $S$), the $X_i$ being disjoint open sets covering it.}
\entree[prop]{17064}{4:44:24}{A morphism $S'\to\coprod_iX_i$ amounts to a decomposition of $S'$ into disjoint open sets $S'_i$ and morphisms $S'_i\to X_i$; in general it factors through no $X_i$.}
\entree[prop]{17128}{4:45:28}{In the category of Zariski sheaves on schemes, the sheaf represented by $\coprod_iX_i$ is the sum of the sheaves represented by the $X_i$.}
\entree[prop]{17192}{4:46:32}{A family of monomorphisms $F_i\to F$ of sheaves makes $F$ the sum of the $F_i$ if $F_i\times_FF_j$ is initial for $i\neq j$ and the $F_i$ cover $F$ (jointly epimorphic).}
\entree[prop]{17324}{4:48:44}{For $i\neq j$, $X_i\times_XX_j$ is the empty scheme; the empty scheme represents the initial sheaf (a point over the empty scheme, the empty set elsewhere).}
\entree[prop]{17450}{4:50:50}{Restriction of Zariski sheaves on $k$-schemes to sheaves on affine $k$-schemes is an equivalence of categories (comparison lemma).}
\face{Tape 2, side A}
\topic{Epimorphisms of Zariski sheaves}
\entree[prop]{17483}{4:51:23}{A family of subsheaves $(F_i)$ of $F$ has $F$ as its supremum if and only if the family of inclusions $F_i\to F$ is epimorphic in Zariski sheaves.}
\entree[prop]{17546}{4:52:26}{A morphism $G\to F$ of Zariski sheaves is an epimorphism if and only if every section of $F$ over $S'$ lifts to $G$ locally on $S'$ for the Zariski topology.}
\entree[prop]{17740}{4:55:40}{Subsheaves $F_i$ cover $F$ if and only if every section of $F$ over $S'$ belongs, locally on $S'$, to one of the $F_i$.}
\entree[prop]{17805}{4:56:45}{For a Zariski sheaf $F$, $U\mapsto F(U)$ is an ordinary sheaf on each $S'$; a morphism of Zariski sheaves is an epimorphism if and only if it is one on each $S'$.}
\entree[prop]{17934}{4:58:54}{For $X=\coprod X_i$, the subsheaves $\widetilde{X_i}$ cover $\widetilde X$: every morphism $S'\to X$ factors locally on $S'$ through some $X_i$ (preimages of the $X_i$).}
\entree[prop]{18065}{5:01:05}{If $U=\bigcup U_i$ is a union of open sets of a scheme $X$, then $\widetilde U$ is the supremum of the subsheaves $\widetilde{U_i}$ of $\widetilde X$.}
\entree[disc]{18195}{5:03:15}{A question why a morphism into a sum of schemes need not factor through one summand, even for sheaves: it factors only locally, and the embedding into sheaves is fully faithful.}
\entree[prop]{18390}{5:06:30}{The inclusion of sheaves into presheaves commutes with all projective limits (products, fibre products), but not in general with inductive limits.}
\entree[prop]{18775}{5:12:55}{For a sum $F=\coprod_{i\in I}F_i$ of Zariski sheaves, a section of $F$ over $S'$ is a decomposition of $S'$ into disjoint open sets $S'_i$ with a section of $F_i$ over each $S'_i$.}
\entree[prop]{18972}{5:16:12}{Every decomposition of $\widetilde{S'}$ as a sum of disjoint subsheaves $F_i$ is of the form $F_i=\widetilde{S'_i}$, where $S'=\coprod S'_i$ is a decomposition of $S'$ into open sets.}
\entree[prop]{19171}{5:19:31}{For the sheaf condition, one may restrict to coverings of $S'$ by affine open sets, then by open sets $D(f)$, which form a base of the topology.}
\entree[prop]{19298}{5:21:38}{For $f_i\in A'$, one has $\mathrm{Spec}\,A'=\bigcup_iD(f_i)$ if and only if the $f_i$ generate the unit ideal.}
\entree[prop]{19365}{5:22:45}{A functor $F$ from algebras to $\mathbf{Set}$ is a Zariski sheaf if, for $f_i\in A'$ generating the unit ideal, $F(A')\to\prod F(A'_{f_i})\rightrightarrows\prod F(A'_{f_if_j})$ is exact.}
\entree[prop]{19428}{5:23:48}{A functor on $S$-algebras that is a sheaf extends, uniquely up to unique isomorphism, to a Zariski sheaf on $S$-schemes.}
\topic{Open immersions of sheaves}
\entree[def]{19840}{5:30:40}{A morphism of sheaves $F\to G$ is an open immersion if, for every $S'\to G$ with $S'$ a scheme, $F\times_GS'\to S'$ is isomorphic to the inclusion of an open set $U'$ of $S'$ (in the sense of $\widetilde{U'}\subset\widetilde{S'}$).}
\entree[prop]{19970}{5:32:50}{A property of morphisms of schemes that is stable under base change extends to morphisms of sheaves: $F\to G$ has it if its base change by every scheme $S'\to G$ has it.}
\entree[prop]{20034}{5:33:54}{If $U'\subset S'$ is open and $S''\to S'$ a morphism, the preimage $U''$ of $U'$ satisfies $U''=S''\times_{S'}U'$, also among locally ringed spaces: open immersions of schemes are stable under base change.}
\entree[prop]{20162}{5:36:02}{Open immersions of sheaves are stable under base change.}
\entree[prop]{20290}{5:38:10}{The composite of two open immersions of sheaves is an open immersion.}
\face{Tape 2, side B}
\topic{A representability criterion}
\entree[thm][representability criterion]{20468}{5:41:08}{A Zariski sheaf $F$ on $\mathrm{Sch}/S$ is representable by a scheme if and only if it has subsheaves $F_i$ with (a) $F_i\to F$ an open immersion, (b) $F=\sup F_i$, (c) $F_i$ representable.}
\entree[prop]{20672}{5:44:32}{In the criterion, (c) may be replaced by (c$'$): each $F_i$ is representable by an affine scheme, that is, for $S=\mathrm{Spec}\,k$, by a $k$-algebra.}
\entree[prop]{20737}{5:45:37}{For $S=\mathrm{Spec}\,k$, the $F_i$ cover $F$ if, for every $a\in F(k')$, there exist $f_\alpha\in k'$ generating the unit ideal and indices $i(\alpha)$ such that $a|_{k'_{f_\alpha}}\in F_{i(\alpha)}(k'_{f_\alpha})$.}
\entree[prop]{20995}{5:49:55}{An open set $U'$ of $\mathrm{Spec}\,k'$ is a union of open sets $D(g_\lambda)$ and $\widetilde{U'}$ is the supremum of the subfunctors represented by the localizations $k'_{g_\lambda}$.}
\entree[prop]{21254}{5:54:14}{The open set $\mathrm{Spec}\,k'\setminus V(J)$ represents the subfunctor $B\mapsto\{u\colon k'\to B\mid u(J)B=B\}$; if $J=(g_\lambda)$, the condition is that the $u(g_\lambda)$ generate the unit ideal of $B$.}
\entree[prop]{21454}{5:57:34}{For a scheme $X$ covered by affine open sets $X_i$, the $\widetilde{X_i}\to\widetilde X$ are open immersions that cover $\widetilde X$: the conditions of the criterion are necessary, with (c$'$).}
\entree[prop]{21516}{5:58:36}{If $F_j\to F$ is an open immersion, $F_i\cap F_j\to F_i$ is one too (base change of $F_j\to F$ by $F_i\to F$).}
\topic{Gluing schemes}
\entree[prop]{21714}{6:01:54}{Schemes $X_i$ over $S$, open sets $X_{ij}\subset X_i$ and isomorphisms $\varphi_{ij}\colon X_{ij}\to X_{ji}$ satisfying the cocycle condition on $X_{ij}\cap X_{ik}$ glue to an $S$-scheme $X$ covered by the $X_i$.}
\entree[prop]{22295}{6:11:35}{If a scheme $X$ is covered by open sets $X_i$, $\widetilde X$ is the gluing of the $\widetilde{X_i}$ along the $\widetilde{X_{ij}}$; likewise for the supremum of subsheaves $F_i\subset F$.}
\topic{Projective space}
\entree[rem]{22643}{6:17:23}{Over a field $k$, for $V$ of dimension $n$, the set of lines of $V$ is the quotient $(V\setminus\{0\})/k^{\times}$.}
\entree[prop]{22770}{6:19:30}{In a basis of $V$, the subsets $X_i=\{x_i\neq0\}$ cover this set of lines, and each is in bijection with $k^{n-1}$ (normalize $x_i=1$).}
\entree[ex]{22963}{6:22:43}{For $n=2$, the projective line is obtained by gluing two copies of $k$ along $k^{\times}$ by $\lambda\mapsto\lambda^{-1}$ (coordinates $y/x$ and $x/y$).}
\face{Tape 3}
\entree[prop]{23316}{6:28:36}{For $V$ finite-dimensional over a field, the lines of $V$ correspond bijectively to the hyperplanes of the dual $V^{\vee}$, and the hyperplanes of $V$ to the $1$-dimensional quotients of $V$.}
\entree[def]{23380}{6:29:40}{$\mathbf{P}(V)$ henceforth denotes the set of hyperplanes of $V$, that is, of $1$-dimensional quotients of $V$ (and not the set of lines).}
\topic{Invertible modules and the functor $\mathbf{P}(V)$}
\entree[cex]{23642}{6:34:02}{If the $f_i$ generate the unit ideal and each $L_{f_i}$ is free of rank $1$, the module $L$ is not necessarily free: being free of rank $1$ is not local on $\mathrm{Spec}\,k'$.}
\entree[def]{23770}{6:36:10}{An $A$-module $L$ is invertible if $\widetilde L$ is locally free of rank $1$, that is, if there exist $f_\alpha$ generating the unit ideal with $L_{f_\alpha}$ free of rank $1$ over $A_{f_\alpha}$.}
\entree[prop]{23835}{6:37:15}{$L$ is invertible if and only if there exists an $A$-module $M$ with $M\otimes_AL\simeq A$; $M$ is then unique up to unique isomorphism compatible with the isomorphisms $M\otimes_AL\simeq A$, and isomorphic to $\mathrm{Hom}_A(L,A)$, denoted $L^{-1}$.}
\entree[prop]{23907}{6:38:27}{$L$ is invertible if and only if it is finitely generated projective and $L_{\mathfrak{p}}$ is free of rank $1$ for every prime ideal $\mathfrak{p}$.}
\entree[prop]{24039}{6:40:39}{For $V$ free of rank $n$, the invertible quotients of $V_{k'}$ correspond to direct-summand submodules of rank $n-1$ at every prime; one must require them to be direct summands.}
\entree[prop]{24104}{6:41:44}{If $L$ is an invertible quotient of $V_{k'}$ and $k'\to k''$ a homomorphism, $L\otimes_{k'}k''$ is an invertible quotient of $V_{k''}$ (right exactness of $\otimes$, stability of invertible modules).}
\entree[def]{24168}{6:42:48}{For a module $V$ over a ring $k$, $\mathbf{P}(V)$ is the functor $k'\mapsto\{\text{invertible quotients of }V\otimes_kk'\}$ on $k$-algebras.}
\entree[thm]{24233}{6:43:53}{For every $k$-module $V$, with no finiteness hypothesis, the functor $\mathbf{P}(V)$ is representable by a scheme, of finite type over $k$ if $V$ is finitely generated.}
\entree[prop]{24492}{6:48:12}{$\mathbf{P}(V)(S')$ is the set of invertible quotients of the inverse image of $\widetilde V$ on $S'$; it is a Zariski sheaf, since the condition is local.}
\entree[prop]{24619}{6:50:19}{If $(e_i)$ generates $V$, $\mathbf{P}_{e_i}(k')$ is the set of quotients $L$ of $V_{k'}$ such that the image of $e_i$ is a basis of $L$; it is a subsheaf of $\mathbf{P}(V)$.}
\entree[prop]{24821}{6:53:41}{For every $e\in V$, $\mathbf{P}_e\to\mathbf{P}(V)$ is an open immersion: the locus where a section of an invertible sheaf is a basis is open.}
\entree[prop]{25011}{6:56:51}{If the $e_i$ generate $V$, the $\mathbf{P}_{e_i}$ cover $\mathbf{P}(V)$: locally on $S'$, one of the $e_i$ is a basis of the invertible quotient.}
\entree[prop]{25080}{6:58:00}{$\mathbf{P}_{e_i}$ is represented by $\mathrm{Sym}(V)/(e_i-1)$: a point is a homomorphism $V_{k'}\to k'$ sending $e_i$ to $1$.}
\entree[ex]{25145}{6:59:05}{If $(e_i)_{i\in I}$ is a basis of $V$, $\mathrm{Sym}(V)/(e_i-1)$ is the polynomial ring in the $T_j$, $j\neq i$: $\mathbf{P}_{e_i}$ is the affine space of dimension $|I|-1$.}
\face{Tape 3, from 7:00:11}
\noindent\textit{A new session, which the description does not date.}\par\smallskip
\topic{Zariski sheaves, reformulated}
\entree[prop]{25465}{7:04:25}{A contravariant functor $F$ on $\mathrm{Sch}/S$ is a Zariski sheaf if, for every $S'$ and every open covering $(S'_i)$, $F(S')\to\prod F(S'_i)\rightrightarrows\prod F(S'_i\cap S'_j)$ is exact.}
\entree[prop]{25594}{7:06:34}{To $F$ one associates, for each $S'$, the presheaf $F_{S'}\colon U'\mapsto F(U')$ on the space $S'$; $F$ is a Zariski sheaf if and only if all the $F_{S'}$ are ordinary sheaves.}
\entree[prop]{25794}{7:09:54}{Every $S$-morphism $f\colon S''\to S'$ induces a morphism of presheaves $F_{S'}\to f_*F_{S''}$, given by $F(U')\to F(f^{-1}U')$, functorial in $f$.}
\entree[prop]{26123}{7:15:23}{A contravariant functor on $\mathrm{Sch}/S$ is equivalent to presheaves $F_{S'}$ functorial in $S'$ such that $F_{S'}|_{S''}\to F_{S''}$ is an isomorphism for every open immersion $S''\to S'$; $F(S')=F_{S'}(S')$.}
\session{19 June 1973}
\noindent\textit{Video 2/4 (\href{https://www.youtube.com/watch?v=thSz9ZHDK9o}{\texttt{thSz9ZHDK9o}}): tape 1 (7:17:49--8:06:52); tape 2 (8:06:52--8:55:21).}\par\medskip
\video{thSz9ZHDK9o}
\face{Tape 1}
\topic{Limits and epimorphisms of Zariski sheaves}
\entree[prop]{26334}{7:18:54}{Projective limits of Zariski sheaves on $\mathrm{Sch}/S$ are computed argument by argument, and $(\varprojlim F_i)_{S'}=\varprojlim(F_i)_{S'}$.}
\entree[prop]{26397}{7:19:57}{For an inductive limit $\varinjlim F_i$ computed in Zariski sheaves, $(\varinjlim F_i)_{S'}=\varinjlim(F_i)_{S'}$, the limit being taken in sheaves on the space $S'$.}
\entree[prop]{26529}{7:22:09}{A morphism $F\to G$ of Zariski sheaves is an epimorphism if and only if $F_{S'}\to G_{S'}$ is one for every $S'$: every section of $G$ over $U'$ is locally the image of a section of $F$.}
\entree[prop]{26850}{7:27:30}{A morphism of $S$-schemes $X\to Y$ is an epimorphism of Zariski sheaves if and only if it admits local sections over $Y$: every point of $Y$ has an open neighbourhood $V$ and a section $V\to X$.}
\entree[prop]{26981}{7:29:41}{A family of $S$-schemes $X_i\to Y$ is epimorphic for Zariski sheaves if and only if every point of $Y$ has a neighbourhood over which one of the $X_i$ admits a section.}
\entree[prop]{27459}{7:37:39}{An open immersion of sheaves is a monomorphism.}
\topic{Projective bundles}
\entree[prop]{27788}{7:43:08}{Representability of a sheaf is local on the base: if $S=\bigcup S_i$ (open sets) and the restriction $F_i$ of $F$ to schemes over $S_i$ is representable for all $i$, then $F$ is.}
\entree[prop]{27917}{7:45:17}{In this case $F_i=S_i\times_SF\to F$ is an open immersion; if $F=\sup F_i$ with open subfunctors $F_i$, then $F$ is representable if and only if each $F_i$ is.}
\entree[prop]{27980}{7:46:20}{An $\mathcal{O}_S$-module is quasi-coherent if it is locally the cokernel of a morphism of free modules; it then comes from a module on each affine open set of $S$.}
\entree[prop]{28043}{7:47:23}{For $\mathcal{M}$ quasi-coherent on $S$, $\mathbf{P}(\mathcal{M})\colon S'\mapsto\{\text{invertible quotients of }\mathcal{M}_{S'}\}$ is a Zariski sheaf on $\mathrm{Sch}/S$.}
\entree[thm]{28112}{7:48:32}{The sheaf $\mathbf{P}(\mathcal{M})$ is representable by an $S$-scheme, the projective bundle associated to $\mathcal{M}$; by localization on $S$, one reduces to the affine case.}
\entree[prop]{28307}{7:51:47}{If $\mathcal{M}$ is generated by global sections $e_i$, $P_i(S')$ is the set of quotients $L$ of $\mathcal{M}_{S'}$ for which the image of $e_i$ is a basis of $L$; $P_i\subset\mathbf{P}(\mathcal{M})$.}
\entree[prop]{28503}{7:55:03}{The morphism $P_i\to\mathbf{P}(\mathcal{M})$ is an open immersion.}
\entree[prop]{28630}{7:57:10}{The set of points of $S'$ where a section of an invertible sheaf is a basis is open.}
\entree[prop]{28693}{7:58:13}{The $P_i$ cover $\mathbf{P}(\mathcal{M})$: if the $e_i$ generate $\mathcal{M}$, locally on $S'$ one of them generates the invertible quotient, hence is a basis of it.}
\entree[prop]{29021}{8:03:41}{For a $k$-module $V$, the functor $k'\mapsto\mathrm{Hom}_{k'}(V_{k'},k')=\mathrm{Hom}_k(V,k')$ is represented by $\mathrm{Sym}_k(V)$.}
\entree[prop]{29157}{8:05:57}{$P_i(k')$ is the set of homomorphisms $V_{k'}\to k'$ sending $e_i$ to $1$; it is represented by $\mathrm{Sym}(V)/(e_i-1)$, so $P_i$ is representable ($S$ affine case).}
\face{Tape 2}
\topic{The universal quotient and vector bundles}
\entree[prop]{29283}{8:08:03}{On $\mathbf{P}=\mathbf{P}(\mathcal{M})$ there is a universal invertible quotient $\mathcal{M}_{\mathbf{P}}\to\mathcal{O}_{\mathbf{P}}(1)$, corresponding to the identity point of the represented functor.}
\entree[def]{29348}{8:09:08}{One sets $\mathcal{O}_{\mathbf{P}}(n)=\mathcal{O}_{\mathbf{P}}(1)^{\otimes n}$ for $n\in\mathbb{Z}$, with $L^{\otimes(-m)}=(L^{\otimes m})^{-1}\simeq(L^{-1})^{\otimes m}$ for $m\geq0$.}
\entree[prop]{29412}{8:10:12}{Universal property: for every invertible quotient of $\mathcal{M}_{S'}$, there exists a unique morphism $S'\to\mathbf{P}$ whose inverse image of the universal quotient is this quotient.}
\entree[def]{29606}{8:13:26}{The vector bundle $\mathbf{V}(\mathcal{M})$ is the functor $S'\mapsto\mathrm{Hom}_{\mathcal{O}_{S'}}(\mathcal{M}_{S'},\mathcal{O}_{S'})$.}
\entree[prop]{29932}{8:18:52}{The subfunctor $\mathbf{V}^{*}$ of epimorphisms $\mathcal{M}_{S'}\to\mathcal{O}_{S'}$ (nonzero on each fibre) is represented by the open complement of the zero section of $\mathbf{V}(\mathcal{M})$.}
\entree[prop]{30064}{8:21:04}{The map $u\mapsto\mathcal{M}_{S'}/\ker u$ defines a morphism $\mathbf{V}^{*}\to\mathbf{P}(\mathcal{M})$, whose image is the free quotients of rank $1$.}
\entree[prop]{30256}{8:24:16}{The morphism $\mathbf{V}^{*}\to\mathbf{P}(\mathcal{M})$ is an epimorphism of sheaves: an invertible quotient is locally free of rank $1$, hence locally of the form $\mathcal{M}_{S'}/\ker u$.}
\topic{Projective space as a quotient by $\mathbb{G}_m$}
\entree[prop]{30385}{8:26:25}{The functor $\mathbb{G}_m\colon S'\mapsto\Gamma(S',\mathcal{O}_{S'})^{\times}$ is represented by the affine line minus the origin, with ring $k[t,t^{-1}]$ over $S=\mathrm{Spec}\,k$.}
\entree[prop]{30521}{8:28:41}{The group $\mathbb{G}_m$ acts on $\mathbf{V}^{*}$ by multiplication; this action is free ($\lambda u=u\Rightarrow\lambda=1$ for $u$ an epimorphism), whereas it is not free on $\mathbf{V}$ (the origin is fixed).}
\entree[prop]{30648}{8:30:48}{The epimorphisms $u$ and $\lambda u$, with $\lambda$ invertible, have the same kernel: $\mathbf{V}^{*}\to\mathbf{P}(\mathcal{M})$ is $\mathbb{G}_m$-invariant.}
\entree[prop]{30775}{8:32:55}{Two epimorphisms $u,v\colon\mathcal{M}_{S'}\to\mathcal{O}_{S'}$ define the same quotient if and only if $v=\lambda u$ with $\lambda\in\Gamma(S',\mathcal{O}_{S'})^{\times}$.}
\entree[prop]{30838}{8:33:58}{The vector bundle minus its zero section is a $\mathbb{G}_{m,\mathbf{P}}$-torsor over $\mathbf{P}(\mathcal{M})$, and $\mathbf{P}(\mathcal{M})=\mathbf{V}^{*}/\mathbb{G}_m$.}
\entree[prop]{30903}{8:35:03}{The group $\mathbb{G}_m=\mathrm{GL}_1$ is the functor $S'\mapsto\mathrm{Aut}(\mathcal{O}_{S'})$ of automorphisms of the module $\mathcal{O}_{S'}$.}
\entree[prop]{30966}{8:36:06}{$\mathbb{G}_m$-torsors over $S$ are equivalent to locally free sheaves of rank $1$ on $S$, via $L\mapsto\mathrm{Isom}_{\mathcal{O}_S}(\mathcal{O}_S,L)$.}
\entree[def]{31032}{8:37:12}{A $\mathbb{G}_m$-torsor over $S$ is a sheaf $F$ equipped with a free action of $\mathbb{G}_m$ such that $F\to S$ is an epimorphism; it is locally isomorphic to $\mathbb{G}_m$.}
\entree[prop]{31099}{8:38:19}{If a sheaf of groups $G$ acts freely on a sheaf $F$, then $F$ is a torsor under $G_T$ over $T=F/G$.}
\entree[prop]{31232}{8:40:32}{The $\mathbb{G}_m$-torsor $\mathbf{V}^{*}\to\mathbf{P}$ is that of trivializations of the universal invertible quotient: a point above $S'\to\mathbf{P}$ is an isomorphism of this quotient with $\mathcal{O}_{S'}$.}
\entree[def]{31617}{8:46:57}{Projective space $\mathbb{P}^n$ over a field $k$ is $\mathbf{P}(k^{n+1})$.}
\entree[prop]{31747}{8:49:07}{A morphism $X\to\mathbb{P}^n$ is equivalent to an invertible sheaf $L$ on $X$ and $n+1$ sections of $L$ that generate it at every point, that is, an invertible quotient of $\mathcal{O}_X^{\,n+1}$.}
\entree[read]{31942}{8:52:22}{The same yoga for analytic geometry is in his exposés at the Cartan seminar. \cite{cartan13}}
\entree[prop]{32079}{8:54:39}{$\mathbf{P}(\mathcal{M})(S')$ is the set of isomorphism classes of pairs $(L,\mathcal{M}_{S'}\to L)$, with $L$ invertible and $\mathcal{M}_{S'}\to L$ an epimorphism.}
\session{20 June 1973}
\noindent\textit{Video 2/4 (\href{https://www.youtube.com/watch?v=thSz9ZHDK9o}{\texttt{thSz9ZHDK9o}}): tape 1 (8:55:21--9:43:38); tape 2 (9:43:38--10:32:26).}\par\medskip
\video{thSz9ZHDK9o}
\face{Tape 1}
\topic{Fibres and base change for projective bundles}
\entree[prop]{32186}{8:56:26}{Base change commutes with the projective bundle: $\mathbf{P}(\mathcal{M})\times_SS'=\mathbf{P}(\mathcal{M}_{S'})$, as functors on schemes over $S'$.}
\entree[prop]{32313}{8:58:33}{The fibre $X_s$ of a scheme $X$ over $S$ at a point $s$ is obtained by base change along $\mathrm{Spec}\,\kappa(s)\to S$; it is a scheme over the field $\kappa(s)$.}
\entree[prop]{32379}{8:59:39}{The underlying space of $X_s$ is the set-theoretic fibre of $X\to S$ over $s$, equipped with the topology induced by that of $X$.}
\entree[prop]{32506}{9:01:46}{The fibre of $\mathbf{P}(\mathcal{M})$ at $s$ is the projective space of the $\kappa(s)$-vector space $\mathcal{M}(s)=\mathcal{M}_s\otimes\kappa(s)=\mathcal{M}_s/\mathfrak{m}_s\mathcal{M}_s$.}
\entree[disc]{32571}{9:02:51}{Asked whether the fibre of a projective bundle embeds as a closed subscheme, he says only for closed points, while a morphism from a projective bundle to a separated scheme has closed image.}
\entree[prop]{32636}{9:03:56}{If $\mathcal{M}$ is finitely generated and $X$ separated over $S$, every $S$-morphism $\mathbf{P}(\mathcal{M})\to X$ has closed image.}
\entree[prop]{32959}{9:09:19}{If $\mathcal{M}\to\mathcal{N}$ is an epimorphism of quasi-coherent modules, there is a morphism $\mathbf{P}(\mathcal{N})\to\mathbf{P}(\mathcal{M})$, which is a closed immersion.}
\entree[prop]{33023}{9:10:23}{For an arbitrary homomorphism $\mathcal{M}\to\mathcal{N}$, there is an open set $U\subset\mathbf{P}(\mathcal{N})$ of invertible quotients $L$ such that $\mathcal{M}_{S'}\to L$ is surjective, and a morphism $U\to\mathbf{P}(\mathcal{M})$.}
\entree[ex]{33086}{9:11:26}{One has $\mathbf{P}(0)=\varnothing$ and $\mathbf{P}(\mathcal{O}_S)=S$; for $\mathcal{M}=\mathcal{O}_S^{2}$, $\mathbf{P}(\mathcal{M})$ is the standard projective line over $S$.}
\topic{Relative spectrum and affine morphisms}
\entree[disc]{33086}{9:11:26}{A question on what happens for the zero module: the projective bundle is empty, since no nonempty scheme has an invertible quotient of zero; the first interesting case is the projective line.}
\entree[prop]{33553}{9:19:13}{Over $S=\mathrm{Spec}\,K$, a quasi-coherent algebra structure on $\widetilde A$ is equivalent to a $K$-algebra structure on the $K$-module $A$.}
\entree[def]{33684}{9:21:24}{For a quasi-coherent $\mathcal{O}_S$-algebra $\mathcal{A}$, $\mathrm{Spec}(\mathcal{A})$ is the functor $S'\mapsto\mathrm{Hom}_{\mathcal{O}_{S'}\text{-alg}}(\mathcal{A}_{S'},\mathcal{O}_{S'})$; over affine $S$, it is $\mathrm{Spec}\,A$.}
\entree[prop]{33941}{9:25:41}{The functor $\mathrm{Spec}(\mathcal{A})$ is a Zariski sheaf representable by an $S$-scheme; over an affine open set $U\subset S$, it is $\mathrm{Spec}\,\Gamma(U,\mathcal{A})$.}
\entree[prop]{34198}{9:29:58}{The functor $\mathcal{A}\mapsto\mathrm{Spec}(\mathcal{A})$, from the opposite category of quasi-coherent $\mathcal{O}_S$-algebras to $S$-schemes, is fully faithful.}
\entree[prop]{34263}{9:31:03}{An $S$-scheme $X$ is affine over $S$ if the preimage of every affine open set of $S$ is affine; the essential image of $\mathrm{Spec}$ consists of the schemes affine over $S$.}
\entree[prop]{34392}{9:33:12}{The spectrum commutes with base change: $\mathrm{Spec}(\mathcal{A})\times_SS'=\mathrm{Spec}(\mathcal{A}_{S'})$.}
\entree[prop]{34520}{9:35:20}{If $f\colon X\to S$ is affine, $f_*\mathcal{O}_X$ is a quasi-coherent $\mathcal{O}_S$-algebra; $X\mapsto f_*\mathcal{O}_X$ is quasi-inverse to $\mathrm{Spec}$: an anti-equivalence with the schemes affine over $S$.}
\entree[prop]{34651}{9:37:31}{For $\mathrm{Spec}\,A\to\mathrm{Spec}\,K$ and $M$ an $A$-module, the direct image of $\widetilde M$ is the sheaf $\widetilde M$ associated to $M$ viewed as a $K$-module.}
\topic{Fibre products of schemes}
\entree[thm]{34882}{9:41:22}{Finite projective limits exist in the category of schemes and in each $\mathrm{Sch}/S$, that is, a final object and fibre products.}
\entree[prop]{34949}{9:42:29}{The final object of schemes is $\mathrm{Spec}\,\mathbb{Z}$, that of $\mathrm{Sch}/S$ is $S$.}
\face{Tape 2}
\entree[prop]{35020}{9:43:40}{If $S=\bigcup S_i$ is an open covering and each $X_i\times_{S_i}Y_i$ is a scheme, then $X\times_SY$ is a scheme: the question is local on $S$.}
\entree[prop]{35088}{9:44:48}{If $X=\bigcup X_i$ and $Y=\bigcup Y_j$ are open coverings, $X\times_SY$ is a scheme if and only if the $X_i\times_SY_j$ are.}
\entree[prop]{35214}{9:46:54}{For open sets $U\subset X$ and $V\subset Y$, $U\times_SV\to X\times_SY$ is an open immersion, and the $X_i\times_SY_j$ cover $X\times_SY$.}
\entree[thm]{35549}{9:52:29}{The fibre product of two $S$-schemes exists; in the affine case, the product $\mathrm{Spec}\,A\times_{\mathrm{Spec}\,K}\mathrm{Spec}\,B$ is $\mathrm{Spec}(A\otimes_KB)$.}
\entree[prop]{35871}{9:57:51}{The morphism $X_s\to X$ is a monomorphism, and a homeomorphism onto its image equipped with the induced topology.}
\entree[prop]{35937}{9:58:57}{In $X\times_SY$, the product $X_i\times_SY_j$ is the intersection of the preimages of $X_i$ and $Y_j$ under the two projections.}
\entree[prop]{36147}{10:02:27}{If the $X_i=\mathrm{Spec}\,\mathcal{A}_i$ are affine over $S$ and form a diagram, their projective limit exists in schemes, is affine over $S$ and equals $\mathrm{Spec}(\varinjlim\mathcal{A}_i)$.}
\entree[rem]{36210}{10:03:30}{In commutative algebras, the sum of a family is the tensor product (possibly infinite, a filtered inductive limit of finite tensor products).}
\entree[prop]{36274}{10:04:34}{An inductive limit of quasi-coherent $\mathcal{O}_S$-modules is quasi-coherent.}
\entree[prop]{36401}{10:06:41}{If $X=\mathrm{Spec}\,\mathcal{A}$ is affine over $S$ and $Z$ is an $S$-scheme with structure morphism $g$, then $\mathrm{Hom}_S(Z,X)\simeq\mathrm{Hom}_{\mathcal{O}_S\text{-alg}}(\mathcal{A},g_*\mathcal{O}_Z)$.}
\entree[rem]{36923}{10:15:23}{For every locally ringed space $Z$, $\mathrm{Hom}(Z,\mathrm{Spec}\,A)=\mathrm{Hom}(A,\Gamma(Z,\mathcal{O}_Z))$; this is what makes the preceding constructions valid in locally ringed spaces.}
\topic{Grassmannians, flags and automorphism groups}
\entree[prop]{37118}{10:18:38}{For $\mathcal{M}$ quasi-coherent, the functor of locally free quotients of given rank of $\mathcal{M}_{S'}$ is representable: it is the Grassmannian bundle associated to $\mathcal{M}$.}
\entree[prop]{37184}{10:19:44}{Likewise, the functor of flags of locally free quotients of $\mathcal{M}_{S'}$ of given ranks $n_1,n_2,\ldots$ is representable (flag bundle).}
\entree[prop]{37247}{10:20:47}{If $\mathcal{M}$ is locally free of rank $n$, the functor $S'\mapsto\mathrm{Aut}(\mathcal{M}_{S'})$ is representable by a group scheme locally isomorphic to $\mathrm{GL}_n$.}
\topic{Localization at a point}
\entree[def]{37565}{10:26:05}{The localization of $X$ at a point $x$ is the affine scheme $\mathrm{Spec}(\mathcal{O}_{X,x})$.}
\entree[prop]{37701}{10:28:21}{For an affine open set $U=\mathrm{Spec}\,A\ni x$, with prime $\mathfrak{p}$, $\mathcal{O}_{X,x}=A_{\mathfrak{p}}$ gives $\mathrm{Spec}(\mathcal{O}_{X,x})\to U\subset X$, independent of the choice of $U$.}
\entree[prop]{37765}{10:29:25}{The morphism $\mathrm{Spec}(\mathcal{O}_{X,x})\to X$ is a monomorphism of schemes, hence injective on points.}
\entree[prop]{37831}{10:30:31}{It is a homeomorphism onto its image: the topology of $\mathrm{Spec}(\mathcal{O}_{X,x})$ is induced by that of $X$.}
\entree[prop]{37831}{10:30:31}{The spectrum of a local ring has a unique closed point, a common specialization of all its other points (the prime ideals are contained in the maximal ideal).}
\entree[prop]{37894}{10:31:34}{A continuous map preserves the specialization relation between points.}
\session{25 June 1973}
\noindent\textit{Video 3/4 (\href{https://www.youtube.com/watch?v=Oqhr-Gt66es}{\texttt{Oqhr-Gt66es}}): tape 1 (0:00:00--0:48:22); tape 2 (0:48:22--1:18:29); tape 2, from 1:18:29 (1:18:29--1:37:19); tape 3, side A (1:37:19--2:26:43); tape 3, side B (2:26:43--3:16:07).}\par\medskip
\video{Oqhr-Gt66es}
\face{Tape 1}
\entree[prop]{1}{0:00:01}{Localization: $\mathrm{Spec}\,\mathcal{O}_{X,x}\to X$ is a monomorphism, with image the set of generizations of $x$ (points $x'$ such that $x\in\overline{\{x'\}}$), the intersection of the open neighbourhoods of $x$.}
\entree[prop]{65}{0:01:05}{For $X$ affine, the prime ideals of $A_{\mathfrak p}$ correspond bijectively to the prime ideals $\mathfrak q\subset\mathfrak p$ of $A$.}
\entree[prop]{128}{0:02:08}{The structure sheaf of the localization is induced by $\mathcal{O}_X$: at every generization $x'$ of $x$, the local rings of $X$ and of the localization are isomorphic, since $(A_{\mathfrak p})_{\mathfrak qA_{\mathfrak p}}=A_{\mathfrak q}$.}
\topic{Local, noetherian and artinian schemes}
\entree[def]{325}{0:05:25}{A local scheme is an affine scheme whose ring is local.}
\entree[prop]{388}{0:06:28}{A scheme is local if and only if it is quasi-compact and has a unique closed point.}
\entree[prop]{784}{0:13:04}{The morphisms from a local scheme $\mathrm{Spec}\,\mathcal{O}$ to $X$ correspond to pairs consisting of a point $x\in X$ (the image of the closed point) and a local homomorphism $\mathcal{O}_{X,x}\to\mathcal{O}$.}
\entree[prop]{909}{0:15:09}{For a ring $A$ and a local ring $\mathcal{O}$, giving $A\to\mathcal{O}$ amounts to giving a prime $\mathfrak p$ of $A$ (the kernel of $A\to\mathcal{O}\to\kappa(\mathcal{O})$) and a local homomorphism $A_{\mathfrak p}\to\mathcal{O}$.}
\entree[prop]{974}{0:16:14}{A morphism $Y\to X$ sending $y$ to $x$ induces $\mathrm{Spec}\,\mathcal{O}_{Y,y}\to\mathrm{Spec}\,\mathcal{O}_{X,x}$: localization is functorial in the pair (scheme, point).}
\entree[prop]{1102}{0:18:22}{The morphisms $\mathrm{Spec}\,k\to X$, with $k$ a field, correspond to pairs of a point $x\in X$ and a field homomorphism $\kappa(x)\to k$.}
\entree[prop]{1232}{0:20:32}{For $x\in X$, $\mathrm{Spec}\,\kappa(x)$ is a closed subscheme of $\mathrm{Spec}\,\mathcal{O}_{X,x}$, and every morphism $\mathrm{Spec}\,k\to X$ with image $x$ factors uniquely through $\mathrm{Spec}\,\kappa(x)$.}
\entree[cex]{1295}{0:21:35}{A scheme reduced to a point is the spectrum of a local ring whose maximal ideal is nil; it is not necessarily the spectrum of a field.}
\entree[prop]{1429}{0:23:49}{A noetherian local ring whose spectrum is a point has a nilpotent maximal ideal and is of finite length, that is, artinian; the converse holds.}
\entree[ex]{1556}{0:25:56}{For $n\geq2$, $k[t]/(t^n)$ is local of length $n$ and $\mathbb{Z}/p^n\mathbb{Z}$ is local artinian; their spectra are points that are not spectra of fields.}
\entree[prop]{1948}{0:32:28}{An irreducible scheme is connected.}
\entree[ex]{2011}{0:33:31}{For $G$ an infinite set, the constant scheme $G_k$ is a disjoint union of infinitely many nonempty open sets, hence is not quasi-compact.}
\entree[def]{2202}{0:36:42}{A scheme is noetherian if it is quasi-compact and covered by finitely many affine open sets $\mathrm{Spec}\,A_i$ with $A_i$ noetherian; then $\Gamma(U,\mathcal{O}_X)$ is noetherian for every affine open set $U$.}
\entree[def]{2270}{0:37:50}{A scheme is locally noetherian if $\Gamma(U,\mathcal{O}_X)$ is noetherian for every affine open set $U$; every quasi-compact open set is then a noetherian scheme.}
\entree[prop]{2525}{0:42:05}{The localization of a noetherian ring is noetherian, as is the completion of a noetherian local ring with respect to the powers of an ideal; but $\widehat{A}\otimes_A\widehat{A}$ is not noetherian in general.}
\entree[def]{2589}{0:43:09}{A scheme is artinian if it is affine with artinian ring (of finite length as a module).}
\entree[prop]{2728}{0:45:28}{An artinian scheme is noetherian and discrete, hence finite: it is a finite sum of spectra of artinian local rings, and its ring is the product of these rings.}
\entree[prop]{2791}{0:46:31}{A scheme is artinian if and only if it is noetherian with discrete underlying space.}
\face{Tape 2}
\topic{Finiteness conditions on morphisms}
\entree[prop]{2904}{0:48:24}{Quasi-compact morphisms are stable under base change.}
\entree[prop]{2968}{0:49:28}{Quasi-compact morphisms are stable under composition.}
\entree[prop]{3033}{0:50:33}{If $g\circ f$ is quasi-compact and $g$ quasi-separated (quasi-compact diagonal), then $f$ is quasi-compact.}
\entree[def]{3158}{0:52:38}{$f$ is locally of finite type if, for all affine open sets $\mathrm{Spec}\,A\subset S$ and $\mathrm{Spec}\,B\subset X$ above it, $B$ is a finitely generated $A$-algebra; of finite type if it is moreover quasi-compact.}
\entree[def]{3227}{0:53:47}{$f$ is locally of finite presentation if, for such affine open sets, $B$ is finitely presented over $A$; this notion and the preceding one are local on $X$.}
\entree[prop]{3292}{0:54:52}{For $X=\mathrm{Spec}\,B$ and $S=\mathrm{Spec}\,A$, $X\to S$ is locally of finite type (resp. finite presentation) if and only if $B$ is finitely generated (resp. finitely presented) over $A$.}
\entree[prop]{3359}{0:55:59}{If $S$ is locally noetherian, \q{locally of finite type} and \q{locally of finite presentation} coincide.}
\entree[prop]{3421}{0:57:01}{Quasi-compact, finite type, locally finite type and locally finite presentation morphisms are stable under base change and under composition.}
\topic{Open immersions, closed immersions and immersions}
\entree[def]{3486}{0:58:06}{An open immersion is a morphism that is a homeomorphism onto an open set of $Y$, the structure sheaf of $X$ being induced by that of $Y$.}
\entree[prop]{3549}{0:59:09}{The subobjects of $Y$ defined by open immersions correspond to the open sets of $Y$; open immersions are stable under composition and base change.}
\entree[def]{3678}{1:01:18}{$f\colon X\to Y$ is a closed immersion if, for a covering of $Y$ by affine open sets $Y_i=\mathrm{Spec}\,A_i$, $f^{-1}(Y_i)=\mathrm{Spec}\,B_i$ is affine and $A_i\to B_i$ is surjective.}
\entree[prop]{3740}{1:02:20}{Intrinsic characterization: $f$ is a closed immersion if and only if it is a homeomorphism onto a closed set of $Y$ and $\mathcal{O}_Y\to f_*\mathcal{O}_X$ is surjective.}
\entree[disc]{3868}{1:04:28}{A question whether the sheaf condition forces a morphism to be affine; the lecturer first agrees, then retracts, saying this holds only if it remains true after any base change.}
\entree[prop]{3941}{1:05:41}{A closed immersion is a monomorphism.}
\entree[prop]{4070}{1:07:50}{The closed subschemes of $Y$ correspond to the quasi-coherent ideals $J\subset\mathcal{O}_Y$, via $J\mapsto V(J)=\mathrm{Spec}(\mathcal{O}_Y/J)$.}
\entree[prop]{4200}{1:10:00}{Closed immersions are stable under base change: $V(J)\times_YY'=V(J\mathcal{O}_{Y'})$, since forming the spectrum of an algebra commutes with base change.}
\entree[prop]{4265}{1:11:05}{For $X=V(J)$, the quasi-coherent ideals of $\mathcal{O}_X$ correspond to the quasi-coherent ideals $K\supset J$ of $\mathcal{O}_Y$ via $K\mapsto K/J$, and $V(K/J)=V(K)$.}
\entree[prop]{4328}{1:12:08}{Closed immersions are stable under composition.}
\entree[prop]{4647}{1:17:27}{If $f$ is quasi-compact and quasi-separated, the direct image of a quasi-coherent $\mathcal{O}_X$-module is quasi-coherent; the kernel of $\mathcal{O}_Y\to f_*\mathcal{O}_X$ for a closed immersion is therefore quasi-coherent.}
\face{Tape 2, from 1:18:29}
\noindent\textit{Probably the beginning of the session of 26 June, recorded immediately afterwards on the same tape.}\par\smallskip
\entree[prop]{4903}{1:21:43}{For $Y=\mathrm{Spec}\,A$ and $J$ an ideal of $A$, $V(J)=\mathrm{Spec}\,A/J$; the map $J\mapsto V(J)$ reverses order.}
\entree[prop]{4968}{1:22:48}{Open subschemes of $Y$ are ordered like the open sets; closed subschemes are ordered like the corresponding quasi-coherent ideals, with the order reversed.}
\entree[prop]{5038}{1:23:58}{The intersection (fibre product) of two closed subschemes $V(J_1)$ and $V(J_2)$ is $V(J_1+J_2)$, since $\mathcal{O}_Y/J_1\otimes_{\mathcal{O}_Y}\mathcal{O}_Y/J_2=\mathcal{O}_Y/(J_1+J_2)$.}
\entree[def]{5165}{1:26:05}{An immersion is a morphism that factors as a closed immersion $X\to U$ followed by an open immersion $U\to Y$.}
\entree[prop]{5297}{1:28:17}{An immersion is a homeomorphism onto a locally closed subset $X$ of $Y$; the largest open set in which $X$ is closed is $Y\setminus(\overline{X}\setminus X)$.}
\entree[prop]{5361}{1:29:21}{Intrinsic characterization: $f$ is an immersion if and only if it is a homeomorphism onto a locally closed subset of $Y$ and $f^{-1}\mathcal{O}_Y\to\mathcal{O}_X$ is surjective.}
\entree[prop]{5426}{1:30:26}{Equivalently, every $y\in f(X)$ has an open neighbourhood $V$ such that $f^{-1}(V)\to V$ is a closed immersion; for $V$ affine, $f^{-1}(V)$ is affine and the ring map is surjective.}
\entree[prop]{5617}{1:33:37}{Immersions are stable under base change and under composition.}
\entree[prop]{5682}{1:34:42}{Every diagonal morphism $X\to X\times_SX$ is an immersion; equivalently, every section of an $S$-scheme is an immersion, as is the kernel of two morphisms $u,v\colon X\to Y$ of $S$-schemes.}
\face{Tape 3, side A}
\topic{Discrete valuation rings; reduced, integral and normal schemes}
\entree[ex]{5841}{1:37:21}{A discrete valuation ring (for instance $\mathbb{Z}_{(p)}$, $k[t]_{(p)}$, $k[[t]]$) is a local integral domain with a single nonzero prime ideal; its spectrum has a generic point and a special point, lying in the closure of the former.}
\entree[cex]{5903}{1:38:23}{A discrete valuation ring is not artinian ($\mathfrak m\supsetneq\mathfrak m^2\supsetneq\cdots$): its spectrum is finite and noetherian without being artinian.}
\entree[prop]{5968}{1:39:28}{The infinitesimal neighbourhood of order $n$ of a point $x$, $\mathrm{Spec}\,\mathcal{O}_{X,x}/\mathfrak m_x^n$, is an artinian local scheme if $X$ is locally noetherian; for $I$ finite, the constant scheme $I_k$ is artinian.}
\entree[prop]{6032}{1:40:32}{For $G$ a group scheme of finite type over a field $k$, the identity component $G^0$ is a subgroup and $G/G^0$ is artinian.}
\entree[def]{6097}{1:41:37}{$X$ is reduced if $\mathcal{O}_X$ has no nonzero nilpotent section over any open set; equivalently, its local rings are reduced.}
\entree[prop]{6161}{1:42:41}{Being reduced is a local property, and $\mathrm{Spec}\,A$ is reduced if and only if $A$ is reduced.}
\entree[def]{6224}{1:43:44}{A scheme is integral if it is reduced and irreducible.}
\entree[prop]{6355}{1:45:55}{If $X$ is covered by nonempty open sets $X_i$, $X$ is integral if and only if the $X_i$ are integral and pairwise meet.}
\entree[prop]{6430}{1:47:10}{An affine scheme $\mathrm{Spec}\,A$ is integral if and only if $A$ is an integral domain.}
\entree[def]{6568}{1:49:28}{A scheme is normal if all its local rings are integrally closed integral domains.}
\entree[prop]{6634}{1:50:34}{A locally noetherian normal scheme is a disjoint sum of integral schemes.}
\entree[prop]{6697}{1:51:37}{An affine integral scheme $\mathrm{Spec}\,A$ is normal if and only if $A$ is integrally closed in its field of fractions.}
\entree[prop]{6899}{1:54:59}{For $X$ integral with generic point $\eta$, the residue field at $\eta$ is the field $K$ of rational functions of $X$.}
\entree[prop]{7023}{1:57:03}{A morphism $X'\to X$ of finite type between integral schemes, $X$ noetherian, is called birational if it induces an isomorphism of function fields; it is then an isomorphism over a nonempty open set of $X$.}
\entree[thm]{7150}{1:59:10}{Zariski's Main Theorem: if $X$ is normal noetherian integral, every separated birational morphism of finite type $X'\to X$ with finite fibres is an open immersion.}
\entree[ex]{7216}{2:00:16}{The blow-up of a point of the affine plane is a birational morphism, with exceptional fibre a projective line, which is not an open immersion.}
\entree[ex]{7280}{2:01:20}{The normalization of a curve with an ordinary double point is birational, with finite fibres, and is not an isomorphism (the local ring at the double point is not normal).}
\entree[ex]{7347}{2:02:27}{The normalization of a curve with a cusp is birational and bijective without being an isomorphism (the local ring is not normal).}
\entree[prop]{7413}{2:03:33}{The normalization morphism $X'\to X$ is birational with finite fibres, and is an isomorphism if and only if $X$ is normal.}
\topic{Regular schemes}
\entree[def]{7602}{2:06:42}{A scheme is regular if all its local rings are regular.}
\entree[disc]{7730}{2:08:50}{For a regular local ring over an algebraically closed field the completion is a power series ring; asked about perfect fields, he explains residue field extensions and generic points break this.}
\entree[def]{7860}{2:11:00}{A noetherian local ring $A$ is regular if its maximal ideal is generated by an $A$-regular sequence $x_1,\dots,x_n$; $x_1$ is a nonzerodivisor in $A$, $x_2$ in $A/x_1A$, and so on.}
\entree[prop]{7990}{2:13:10}{For an $M$-regular sequence (defined likewise on a module $M$), the associated Koszul complex is acyclic.}
\entree[prop]{8055}{2:14:15}{A field is regular; if $A$ is regular, so is $A[t]$; hence $k[t_1,\dots,t_n]$ is regular, as are $\mathbb{G}_{m,k}$ and $\mathrm{GL}_{n,k}$ (open sets of affine space).}
\entree[thm][Serre]{8121}{2:15:21}{A noetherian local ring is regular if and only if its global dimension is finite; the localization of a regular local ring at a prime is therefore regular.}
\entree[prop]{8186}{2:16:26}{A ring is called regular if its localizations at all primes are; principal ideal rings and Dedekind rings are regular (for example $\mathbb{Z}$ and rings of integers of number fields), as are their polynomial rings.}
\entree[prop]{8250}{2:17:30}{A noetherian local ring is regular if and only if its completion $\widehat{A}$ is.}
\entree[prop]{8313}{2:18:33}{A complete local algebra $\mathcal{O}$ over a field $k$, with residue field $k$, is regular if and only if $\mathcal{O}\simeq k[[t_1,\dots,t_n]]$ as a $k$-algebra; if $A$ is regular, so is $A[[t]]$.}
\entree[disc]{8586}{2:23:06}{A question on monoid schemes of finite presentation: a monomorphic endomorphism is an isomorphism, so left cancellation would make a monoid scheme a group, with a parallel statement for modules.}
\entree[prop]{8653}{2:24:13}{An endomorphism of a scheme of finite presentation over $S$ that is a monomorphism is an isomorphism (left as an exercise); consequently, a monoid scheme of finite presentation over the base whose left translations are monomorphisms (left cancellation on all points) is a group object.}
\face{Tape 3, side B}
\topic{Surjective endomorphisms and classes of morphisms}
\entree[prop]{8805}{2:26:45}{A surjective endomorphism of a finitely generated module over a commutative ring is bijective.}
\entree[disc]{8805}{2:26:45}{The audience supplies the determinant argument that a surjective endomorphism of a finitely generated module is bijective; the lecturer contrasts it with his own limit argument for algebras of finite presentation.}
\entree[prop]{9207}{2:33:27}{A surjective endomorphism of a finitely presented algebra over a ring is an isomorphism (noetherian case: the kernels of the iterates stabilize; general case: pass to the limit).}
\entree[rem]{9336}{2:35:36}{If \q{surjective} is replaced by \q{epimorphism} in this statement, it becomes subtler: the noetherian property of the rings $k_i$ and arguments with Artin rings are then needed.}
\entree[prop]{9661}{2:41:01}{Affine morphisms contain the isomorphisms and are stable under composition and under base change.}
\entree[prop]{9795}{2:43:15}{A class of morphisms containing the isomorphisms and stable under composition and base change is stable under fibre product: $f_1\times_Sf_2\in M$ if $f_1,f_2\in M$.}
\entree[prop]{10079}{2:47:59}{For such a class $M$: if $g\circ f\in M$ and the diagonal $\Delta_g\colon Y\to Y\times_ZY$ is in $M$, then $f\in M$.}
\entree[prop]{10336}{2:52:16}{Giving a morphism $f\colon X\to Y$ over $Z$ is equivalent to giving a section $\Gamma_f=(\mathrm{id},f)$ of $X\times_ZY\to X$; $f$ is the composite of $\Gamma_f$ and the second projection.}
\entree[prop]{10408}{2:53:28}{The graph $\Gamma_f$ is a monomorphism and is deduced from $\Delta_g\colon Y\to Y\times_ZY$ by the base change $f\times\mathrm{id}_Y$.}
\entree[prop]{10536}{2:55:36}{For $g\colon Y\to S$ and a class $M$ as above, the following are equivalent: $\Delta_g\in M$; every graph $\Gamma_f\colon X\to X\times_SY$ is in $M$; for all $u,v\colon X\to Y$, the kernel of $(u,v)$ in $X$ is in $M$.}
\entree[prop]{10860}{3:01:00}{These conditions are also equivalent to: every section of every base change $Y\times_SS'\to S'$ is in $M$.}
\entree[ex]{10987}{3:03:07}{The affine plane with doubled origin (two copies of the plane glued along the complement of the origin) has a non-affine diagonal.}
\entree[prop]{11051}{3:04:11}{Monomorphisms are stable under base change and composition; every diagonal morphism is a monomorphism; if $g\circ f$ is a monomorphism, so is $f$.}
\entree[prop]{11114}{3:05:14}{Surjective morphisms (in the set-theoretic sense) are stable under composition and base change.}
\entree[prop]{11246}{3:07:26}{Epimorphisms of schemes are stable under composition but not under base change.}
\entree[def]{11579}{3:12:59}{$f\colon X\to S$ is quasi-compact if the preimage of every quasi-compact open set of $S$ is quasi-compact.}
\entree[prop]{11643}{3:14:03}{Quasi-compactness of a morphism is local on the base: if $S$ is covered by affine open sets $S_i$, $f$ is quasi-compact if and only if each $f^{-1}(S_i)$ is.}
\session{26 June 1973}
\noindent\textit{Video 3/4 (\href{https://www.youtube.com/watch?v=Oqhr-Gt66es}{\texttt{Oqhr-Gt66es}}): tape 1 (3:16:07--4:04:57); tape 2, side A (4:04:57--4:36:34); tape 3, side A (4:36:34--5:24:51); tape 3, side B (5:24:51--5:56:13); tape 1 (5:56:13--6:44:54).}\par\medskip
\video{Oqhr-Gt66es}
\face{Tape 1 (3:16:07--4:04:57)}
\topic{Diagonals and separated morphisms}
\entree[prop]{12226}{3:23:46}{The diagonal of an affine scheme $\mathrm{Spec}\,B$ over $\mathrm{Spec}\,A$ is the closed immersion given by $B\otimes_AB\to B$, $x\otimes y\mapsto xy$; it follows that the diagonal of every $S$-scheme is an immersion.}
\entree[prop]{12294}{3:24:54}{If $g\circ f$ is an immersion, then $f$ is one.}
\entree[prop]{12483}{3:28:03}{An immersion is locally of finite type; the diagonal of every morphism is therefore locally of finite type.}
\entree[prop]{12547}{3:29:07}{An immersion is locally of finite presentation if and only if the quasi-coherent ideal $J$ defining it on an open set is of finite type.}
\entree[prop]{12611}{3:30:11}{If $f$ is locally of finite type, its diagonal is locally of finite presentation.}
\entree[prop]{12804}{3:33:24}{If $B$ is generated as an $A$-algebra by elements $b_i$, the kernel of $B\otimes_AB\to B$ is generated by the $b_i\otimes1-1\otimes b_i$.}
\entree[def]{12999}{3:36:39}{$f\colon X\to S$ is separated if its diagonal, which is an immersion, is a closed immersion, that is, if $\Delta_f(X)$ is closed in $X\times_SX$.}
\entree[prop]{13063}{3:37:43}{Equivalent: after any base change, every section has closed image; the graph of every $S$-morphism $Y\to X$ is closed; the kernel of every pair $(u,v)\colon Y\to X$ is a closed subscheme of $Y$.}
\entree[prop]{13255}{3:40:55}{Extension of identities: if $X$ is separated over $S$, $Y$ reduced, and two $S$-morphisms $u,v\colon Y\to X$ coincide on a dense open set $U$ of $Y$, then $u=v$.}
\entree[prop]{13384}{3:43:04}{Same conclusion if $u\circ i=v\circ i$ for a morphism $i\colon U\to Y$ with dense image.}
\entree[prop]{13576}{3:46:16}{A surjective immersion is a closed immersion defined by a nil ideal; if $Y$ is reduced, it is an isomorphism.}
\entree[prop]{13639}{3:47:19}{Every monomorphism is separated (its diagonal is an isomorphism), in particular every immersion is separated.}
\entree[prop]{13704}{3:48:24}{If $g\circ f$ is separated, then $f$ is separated.}
\entree[def]{13770}{3:49:30}{A scheme is separated if it is separated over $\mathrm{Spec}\,\mathbb{Z}$; an affine scheme is separated, since $A\otimes A\to A$ is surjective.}
\entree[prop]{13832}{3:50:32}{If $S$ is separated, an $S$-scheme $X$ is separated if and only if it is separated over $S$.}
\entree[prop]{13961}{3:52:41}{If $S$ is separated, every morphism from an affine scheme $X$ to $S$ is affine; in particular $X\times_SY$ is affine for $X,Y$ affine over $S$.}
\entree[prop]{14155}{3:55:55}{In a separated scheme, the intersection of two affine open sets is affine; the converse is false.}
\entree[cex]{14219}{3:56:59}{The affine line with doubled origin is not separated, although its diagonal is affine: two morphisms from the line to it coincide away from the origin without being equal.}
\face{Tape 2, side A}
\entree[ex]{14698}{4:04:58}{The complement of the zero section of $\mathbb{A}^n_S$ is affine for $n=1$ and is not for $n\geq2$ ($S$ nonempty affine).}
\entree[cex]{14762}{4:06:02}{The plane with doubled origin contains two affine open sets whose intersection, the complement of the origin, is not affine.}
\topic{Quasi-separated and coherent morphisms}
\entree[def]{14833}{4:07:13}{$f$ is quasi-separated if its diagonal is quasi-compact.}
\entree[prop]{14898}{4:08:18}{Equivalent: after any base change, every section is quasi-compact; or the kernel of any pair $(u,v)\colon Y\to X$ of $S$-morphisms is quasi-compact over $Y$.}
\entree[prop]{14967}{4:09:27}{A separated morphism is quasi-separated; every monomorphism, hence every diagonal morphism, is quasi-separated; if $g\circ f$ is quasi-separated, $f$ is.}
\entree[def]{15031}{4:10:31}{A scheme is quasi-separated if it is so over $\mathrm{Spec}\,\mathbb{Z}$; it is then so over every scheme beneath it.}
\entree[prop]{15094}{4:11:34}{If $S$ is quasi-separated, an $S$-scheme is quasi-separated if and only if it is so over $S$.}
\entree[prop]{15158}{4:12:38}{A scheme is quasi-separated if and only if the intersection of two quasi-compact open sets is quasi-compact.}
\entree[prop]{15289}{4:14:49}{$S$ is quasi-separated if and only if every morphism $X\to S$ from a quasi-compact scheme is quasi-compact.}
\entree[prop]{15416}{4:16:56}{$S$ is quasi-separated if and only if $X\times_SY$ is quasi-compact for all $X,Y$ quasi-compact over $S$.}
\entree[prop]{15613}{4:20:13}{If $S$ is covered by affine open sets $S_i$, $S$ is quasi-separated if and only if the $S_i\cap S_j$ are quasi-compact.}
\entree[cex]{15677}{4:21:17}{The complement of the origin in $\mathrm{Spec}\,k[t_1,t_2,\dots]$ is not quasi-compact; two copies glued along this open set form a scheme that is not quasi-separated.}
\entree[disc]{15804}{4:23:24}{Asked what must be added to quasi-separatedness to get separatedness, the lecturer finds nothing beyond the definition itself, after an example of a non-quasi-compact intersection.}
\entree[prop]{16003}{4:26:43}{If $g\circ f$ is quasi-compact (resp. of finite type) and $g$ quasi-separated, then $f$ is quasi-compact (resp. of finite type); the diagonal of a morphism locally of finite type is quasi-compact if and only if it is of finite type.}
\entree[def]{16068}{4:27:48}{A morphism is coherent if it is quasi-compact and quasi-separated; a scheme is coherent if it is so over $\mathrm{Spec}\,\mathbb{Z}$.}
\entree[prop]{16201}{4:30:01}{Coherent morphisms are stable under base change and under composition.}
\entree[prop]{16267}{4:31:07}{If $g\circ f$ is coherent and $g$ quasi-separated, then $f$ is coherent.}
\entree[rem]{16329}{4:32:09}{Nothing can be concluded about $g$ from $g\circ f$ coherent alone (take $f$ from the empty scheme): $f$ must have a reasonably large image.}
\entree[prop]{16396}{4:33:16}{If $S$ is quasi-separated, an $S$-scheme is coherent if and only if it is coherent over $S$.}
\entree[prop]{16474}{4:34:34}{If $f\colon X\to S$ is coherent and $\mathcal{F}$ a quasi-coherent $\mathcal{O}_X$-module, $f_*\mathcal{F}$ is quasi-coherent.}
\entree[prop]{16538}{4:35:38}{Under the same conditions, the higher direct images $R^qf_*\mathcal{F}$ are quasi-coherent.}
\face{Tape 3, side A}
\topic{Geometric reducedness and irreducibility}
\entree[prop]{16790}{4:39:50}{For a field $k$, the following are equivalent: $k$ is perfect; $A\otimes_kB$ is reduced for all reduced $k$-algebras $A,B$; $A\otimes_kk'$ is reduced for $A$ reduced and $k'/k$ an extension; $k'\otimes_kk'$ is reduced for every finite purely inseparable extension $k'/k$.}
\entree[cex]{17048}{4:44:08}{If $x\in k\setminus k^p$ and $k'=k[t]/(t^p-x)$, then $k'\otimes_kk'\simeq k'[t]/(t-\alpha)^p$, with $\alpha^p=x$, is not reduced.}
\entree[prop]{17116}{4:45:16}{If $k$ is perfect and $K,L$ are two extensions of $k$, then $K\otimes_kL$ is reduced.}
\entree[prop]{17370}{4:49:30}{For a $k$-algebra $A$ and a perfect extension $K_1$ of $k$, the following are equivalent ($A$ is then called geometrically reduced, or separable): $A\otimes_kB$ reduced for every reduced $B$; $A\otimes_kk'$ reduced for every extension $k'$; $A\otimes_kK_1$ reduced; $A\otimes_kk'$ reduced for every finite purely inseparable extension $k'$.}
\entree[def]{17834}{4:57:14}{A $k$-scheme $X$ is geometrically reduced, or separable, if $X\times_kY$ is reduced for every reduced $k$-scheme $Y$, which is equivalent to the analogous conditions on extensions; one reduces to $X$ affine.}
\entree[prop]{17899}{4:58:19}{A scheme separable over $k$ is reduced; the converse holds for every $X$ if and only if $k$ is perfect.}
\entree[prop]{17970}{4:59:30}{For every extension $k'/k$, $X$ is separable over $k$ if and only if $X_{k'}$ is over $k'$.}
\entree[prop]{18099}{5:01:39}{A field $k$ is separably closed if every algebraic extension of $k$ is purely inseparable; in characteristic zero, every field is perfect and separably closed is equivalent to algebraically closed.}
\entree[prop]{18295}{5:04:55}{For a field $k$, the following are equivalent: $k$ is separably closed; for $A,B$ with irreducible spectra, $\mathrm{Spec}(A\otimes_kB)$ is irreducible; likewise for $B$ an extension; $k'\otimes_kk'$ is an integral domain for every finite Galois extension $k'$.}
\entree[prop]{18360}{5:06:00}{If $k$ is separably closed and $K,L$ are extensions of $k$, the spectrum of $K\otimes_kL$ is irreducible.}
\entree[prop]{18552}{5:09:12}{For a $k$-scheme $X$ and $K_1$ separably closed, the following are equivalent ($X$ is then called geometrically irreducible): $X\times_kY$ irreducible for every irreducible $Y$; $X_{k'}$ irreducible for every extension $k'$; $X_{K_1}$ irreducible; $X_{k'}$ irreducible for $k'/k$ finite Galois.}
\entree[prop]{19341}{5:22:21}{A geometrically irreducible scheme is irreducible; the converse holds for every $X$ if and only if $k$ is separably closed.}
\entree[prop]{19407}{5:23:27}{For every extension $k'/k$, $X$ is geometrically irreducible over $k$ if and only if $X_{k'}$ is over $k'$.}
\face{Tape 3, side B}
\topic{Geometric integrality and connectedness}
\entree[prop]{19493}{5:24:53}{The product of two separable (resp. geometrically irreducible) $k$-schemes is separable (resp. geometrically irreducible).}
\entree[prop]{19691}{5:28:11}{For a field $k$, the following are equivalent: $k$ is algebraically closed; the tensor product of two integral $k$-algebras is integral; $A\otimes_kk'$ is integral for $A$ integral and $k'/k$ an extension; $k'\otimes_kk'$ is integral for every finite extension $k'/k$.}
\entree[prop]{19761}{5:29:21}{For a $k$-scheme $X$ and $K_1$ algebraically closed, the following are equivalent: $X\times_kY$ integral for every integral $Y$; $X_{k'}$ integral for every extension $k'$; $X_{K_1}$ integral; $X_{k'}$ integral for every finite extension $k'$.}
\entree[def]{19828}{5:30:28}{$X$ is geometrically integral if it is separable and geometrically irreducible; it is then integral, and the converse holds for every $X$ if and only if $k$ is algebraically closed.}
\entree[prop]{19891}{5:31:31}{Geometric integrality is invariant under extension of the base field and stable under the product of two $k$-schemes.}
\entree[prop]{19954}{5:32:34}{If $k$ is separably closed, the product over $k$ of two connected schemes is connected.}
\entree[prop]{20091}{5:34:51}{If $k'/k$ is a purely inseparable extension, the projection $X_{k'}\to X$ is a homeomorphism for every $k$-scheme $X$.}
\entree[prop]{20413}{5:40:13}{For $Y$ nonempty, the projection $X\times_kY\to X$ is surjective and open, so the topology of $X$ is a quotient of that of $X\times_kY$.}
\entree[prop]{20476}{5:41:16}{If $k$ is separably closed and $X$ connected nonempty, then $X_{k'}$ is connected nonempty for every extension $k'/k$.}
\entree[prop]{20602}{5:43:22}{The fibre of $X\times_kY\to X$ at $x$ is $\mathrm{Spec}(\kappa(x)\otimes_k\mathcal{O}_Y)$; in particular that of $X_{k'}\to X$ at $x$ is $\mathrm{Spec}(\kappa(x)\otimes_kk')$; that of $X\times_kY\to Y$ at $y$ is $X_{\kappa(y)}$.}
\entree[prop]{20796}{5:46:36}{If the product of two connected schemes over $k$ is always connected, then $k$ is separably closed: for $k'/k$ finite Galois with group $\Gamma$, $\mathrm{Spec}(k'\otimes_kk')$ is a sum of $|\Gamma|$ points.}
\entree[prop]{20993}{5:49:53}{For a $k$-scheme $X$ and $K_1$ separably closed, the following are equivalent ($X$ is then called geometrically connected): $X\times_kY$ connected for every connected $Y$; $X_{k'}$ connected for every extension $k'$; $X_{K_1}$ connected; $X_{k'}$ connected for $k'/k$ finite Galois.}
\entree[prop]{20993}{5:49:53}{A geometrically connected scheme (satisfying these conditions) is connected; the converse holds for every $X$ if and only if $k$ is separably closed.}
\entree[prop]{21188}{5:53:08}{A filtered projective limit of connected schemes, affine over a quasi-compact scheme $X$, is connected.}
\entree[prop]{21315}{5:55:15}{Geometric connectedness is invariant under extension of the base field and stable under the product of two $k$-schemes.}
\face{Tape 1 (5:56:13--6:44:54)}
\noindent\textit{Dated 26 June by the video description; probably the session of 27 June.}\par\smallskip
\topic{Complements on subschemes and coherence}
\entree[prop]{21601}{6:00:01}{For subsheaves $F_i$ of a sheaf $F$, the supremum is the image of $\coprod_iF_i\to F$, the smallest subsheaf containing the $F_i$; its sections are those that belong locally to one of the $F_i$.}
\entree[ex]{22316}{6:11:56}{The scheme obtained by gluing infinitely many (indexed by $\mathbb{Z}$) copies of a curve $S$ along $S\setminus\{s_0\}$ is connected, locally of finite type over $S$, non-separated and non-quasi-compact.}
\entree[def]{22510}{6:15:10}{A module $\mathcal{F}$ on a ringed space is coherent if it is of finite type and the kernel of every morphism $\mathcal{O}_U^n\to\mathcal{F}$ is of finite type.}
\entree[thm][Serre]{22572}{6:16:12}{Coherent modules are stable under finite sums, kernels, cokernels, images and extensions.}
\entree[prop]{22637}{6:17:17}{If $\mathcal{O}_X$ is coherent, the coherent modules are exactly those of finite presentation; for $X$ locally noetherian, $\mathcal{O}_X$ is coherent.}
\entree[prop]{22896}{6:21:36}{A surjective immersion is an open immersion if and only if it is an isomorphism; there exist surjective immersions with target $Y$ that are not isomorphisms if and only if $Y$ is not reduced.}
\entree[prop]{23161}{6:26:01}{For $Y$ closed in $X$, the ideal $J_Y$ of sections vanishing at every point of $Y$ is quasi-coherent; it is the largest ideal defining $Y$, and $V(J_Y)$ is the unique reduced closed subscheme with underlying space $Y$.}
\entree[prop]{23290}{6:28:10}{For a closed subscheme $V(J)$, the $V(J^{n+1})$ have the same underlying space; they are its infinitesimal neighbourhoods of order $n$.}
\entree[prop]{23353}{6:29:13}{An immersion whose underlying map is open is a closed immersion defined by a nil ideal on an open set of $Y$; it is not an open immersion in general.}
\entree[prop]{23421}{6:30:21}{Every locally noetherian scheme is quasi-separated.}
\entree[prop]{23618}{6:33:38}{A scheme is coherent if and only if it is a finite union of affine open sets $X_i$ whose intersections $X_i\cap X_j$ are quasi-compact.}
\entree[cex]{24017}{6:40:17}{The quasi-compactness hypothesis is necessary: for $X$ a disjoint sum of infinitely many copies of an affine scheme $S$, $f_*\mathcal{F}$ is the product of the $\mathcal{F}_i$, and an infinite product of quasi-coherent modules is not quasi-coherent in general.}
\entree[disc]{24017}{6:40:17}{Why quasi-compactness is needed for direct images of quasi-coherent sheaves: an infinite disjoint union gives infinite products of quasi-coherent sheaves, which fail to be quasi-coherent.}
\entree[def]{24145}{6:42:25}{An open set is retrocompact if its intersection with every quasi-compact open set is quasi-compact, that is, if the inclusion is quasi-compact, hence coherent.}
\entree[prop]{24207}{6:43:27}{If $U\subset X$ is a retrocompact open set, every quasi-coherent $\mathcal{O}_U$-module $\mathcal{F}$ is the restriction to $U$ of a quasi-coherent module on $X$, for instance $i_*\mathcal{F}$.}
\session{2 July 1973}
\noindent\textit{Video 3/4 (\href{https://www.youtube.com/watch?v=Oqhr-Gt66es}{\texttt{Oqhr-Gt66es}}): tape 1, side A (6:44:54--7:33:37); tape 1, side B (7:33:37--8:22:13); tape 2, side A (8:22:13--8:28:48).}\par\medskip
\video{Oqhr-Gt66es}
\face{Tape 1, side A}
\topic{Finite presentation and limits}
\entree[prop]{24432}{6:47:12}{If $A=\varinjlim A_i$ is a filtered limit of rings, a quasi-compact quasi-separated scheme of finite presentation over $A$ comes by base change from a scheme of finite presentation over $A_i$, for $i$ large enough.}
\topic{Projective and ample morphisms}
\entree[def][naive]{25093}{6:58:13}{$X\to Y$ is projective if there is a closed immersion of $X$ into $\mathbb{P}(E)$, with $E$ quasi-coherent of finite type over $Y$; quasi-projective if one only has an immersion.}
\entree[prop]{25222}{7:00:22}{If $E$ is quasi-coherent of finite type, $\mathbb{P}(E)$ is of finite type over $Y$; $\mathbb{P}(E)$ is separated over $Y$.}
\entree[prop]{25350}{7:02:30}{The naive notion of projective morphism is stable under base change, since $\mathbb{P}(E)\times_YY'=\mathbb{P}(E_{Y'})$.}
\entree[cex]{25477}{7:04:37}{The naive notion of projective morphism is not local on the base: even for $Y$ quasi-compact, a morphism that is projective over each open set of a cover need not embed into some $\mathbb{P}(E)$ over $Y$ (counterexamples exist), the embeddings over the open sets not gluing.}
\entree[rem]{25603}{7:06:43}{Making the notion local on the base (locally projective morphisms) does not repair it: the composite of two locally projective morphisms need not be locally projective.}
\entree[prop]{25734}{7:08:54}{The $Y$-morphisms $X\to\mathbb{P}(E)$ correspond to the invertible quotients $L$ of the module $E_X$, the inverse image of $E$ on $X$.}
\entree[def]{25800}{7:10:00}{An invertible sheaf $L$ on $X$ is very ample relative to $Y$ if there exist a quasi-coherent module $E$ on $Y$ and an epimorphism $f^*E\to L$ whose associated morphism $X\to\mathbb{P}(E)$ is an immersion.}
\entree[prop]{25928}{7:12:08}{If $X$ has a sheaf very ample relative to $Y$, then $X$ is separated over $Y$; if moreover $X$ is of finite type over $Y$ coherent, one may choose $E$ of finite type.}
\entree[prop]{26255}{7:17:35}{If $Y$ is coherent, $X\to Y$ admits an immersion into $\mathbb{P}(E)$ with $E$ quasi-coherent of finite type if and only if $X$ is of finite type over $Y$ and has an invertible sheaf very ample relative to $Y$.}
\entree[prop]{26447}{7:20:47}{The property of being very ample is stable under base change.}
\entree[prop]{26511}{7:21:51}{The tensor product of two very ample sheaves is very ample (Segre embedding $\mathbb{P}(E)\times\mathbb{P}(F)\to\mathbb{P}(E\otimes F)$); hence $L^{\otimes n}$ is very ample for $n>0$ if $L$ is.}
\entree[prop]{26648}{7:24:08}{If $Y$ is quasi-compact and $X$ of finite type over $Y$, $L$ is ample relative to $Y$ if and only if a power $L^{\otimes n}$ ($n>0$) is very ample; then $L^{\otimes n}$ is very ample for all $n\geq n_0$.}
\entree[prop]{26840}{7:27:20}{Ample sheaves are stable under base change and tensor product; for $n>0$, $L$ is ample if and only if $L^{\otimes n}$ is.}
\entree[disc]{26908}{7:28:28}{A question whether the inverse of an ample sheaf is ample leads to the Picard group, the cone of ample classes, and the conclusion that this happens only when the structure sheaf is ample.}
\entree[prop]{27035}{7:30:35}{In $\mathrm{Pic}(X)\otimes\mathbb{Q}$, the classes some positive multiple of which is ample form a convex cone.}
\entree[prop]{27098}{7:31:38}{For $Y$ quasi-compact, $\mathcal{O}_X$ is ample relative to $Y$ if and only if $X$ is quasi-affine over $Y$; if moreover $X$ is proper over $Y$, then $X\to Y$ is finite.}
\face{Tape 1, side B}
\entree[prop]{27219}{7:33:39}{If $Y$ is quasi-compact, $L$ ample and $M$ any invertible sheaf, then $L^{\otimes n}\otimes M$ is ample for $n$ large enough; the ample cone therefore spans $\mathrm{Pic}(X)\otimes\mathbb{Q}$ as soon as it is nonempty.}
\topic{Quasi-projective and proper morphisms}
\entree[def]{27285}{7:34:45}{$f\colon X\to Y$ is quasi-projective if it is of finite type and there exists an invertible sheaf on $X$ ample relative to $Y$.}
\entree[def]{27355}{7:35:55}{$f$ is projective if it is quasi-projective and proper; $f$ is proper if it is of finite type, separated and universally closed.}
\entree[prop]{27481}{7:38:01}{If $Y$ is coherent, $f$ is quasi-projective (resp. projective) if and only if there is an immersion (resp. closed immersion) of $X$ into $\mathbb{P}(E)$ with $E$ quasi-coherent of finite type.}
\entree[prop]{27544}{7:39:04}{Quasi-projective and projective morphisms are stable under base change.}
\entree[prop]{28150}{7:49:10}{Proper morphisms are stable under base change and under composition.}
\entree[prop]{28218}{7:50:18}{If $X$ is proper over $Y$ and $Z$ separated over $Y$, every $Y$-morphism $X\to Z$ is closed.}
\entree[prop]{28347}{7:52:27}{If $X$ is of finite type over $Y$ and $g\colon X\to Z$ a $Y$-immersion into a proper $Y$-scheme $Z$, then $X$ is proper over $Y$ if and only if $g$ is a closed immersion; closed immersions are proper.}
\entree[prop]{28411}{7:53:31}{Projective morphisms are proper; over an algebraically closed field, there exist proper schemes that are not projective.}
\entree[thm]{28486}{7:54:46}{Chow's lemma: every proper scheme over a field is dominated by a projective scheme by means of a birational morphism.}
\entree[disc]{28486}{7:54:46}{Proper does not imply projective, though Chow's lemma gives a birational reduction; projective implies proper is not trivial and uses reduction to Noetherian base and valuative criteria.}
\entree[prop]{28613}{7:56:53}{If $E$ is quasi-coherent of finite type, $\mathbb{P}(E)$ is proper over $Y$; one reduces to $\mathbb{P}^n_{\mathbb{Z}}$ over $\mathrm{Spec}\,\mathbb{Z}$.}
\topic{Valuative criteria}
\entree[thm]{28743}{7:59:03}{Valuative criterion of properness: for $Y$ locally noetherian and $f$ of finite type, $f$ is proper if and only if, for every discrete valuation ring $A$ with fraction field $K$ and every $\mathrm{Spec}\,A\to Y$, the sections of $X$ over $\mathrm{Spec}\,A$ correspond bijectively to those over $\mathrm{Spec}\,K$.}
\entree[prop]{28870}{8:01:10}{Injectivity of this map is equivalent to separatedness of $f$ (valuative criterion of separatedness).}
\entree[prop]{29317}{8:08:37}{For $A$ a discrete valuation ring with fraction field $K$ and $M$ a finitely generated $A$-module, the quotients of $M$ that are free of rank one correspond bijectively to the one-dimensional quotients of $M_K$, via $H\mapsto M\cap H_K$; hence $\mathbb{P}(M)$ satisfies the criterion.}
\topic{Integral extensions}
\entree[prop]{29521}{8:12:01}{For an $A$-algebra $B$ and $x\in B$, the following are equivalent: $x$ satisfies a monic polynomial with coefficients in $A$ (integral over $A$); $A[x]$ is a finitely generated $A$-module; $x$ belongs to a subalgebra of $B$ finite over $A$.}
\entree[prop]{29585}{8:13:05}{An $A$-algebra $B$ is integral if and only if it is a union (inductive limit) of subalgebras finite over $A$.}
\entree[prop]{29651}{8:14:11}{Subalgebras and quotient algebras of an integral algebra are integral; the tensor product of two integral algebras is integral; the elements of $B$ integral over $A$ form a subalgebra, the integral closure of $A$ in $B$.}
\entree[prop]{29718}{8:15:18}{Transitivity: if $B$ is integral over $A$ and $C$ integral over $B$, then $C$ is integral over $A$.}
\topic{Finite and integral morphisms}
\entree[def]{29849}{8:17:29}{$f\colon X\to Y$ is finite (resp. integral) if it is affine, $X=\mathrm{Spec}\,\mathcal{B}$, with $\mathcal{B}$ a quasi-coherent $\mathcal{O}_Y$-algebra finitely generated as a module (resp. integral over $\mathcal{O}_Y$).}
\entree[prop]{29981}{8:19:41}{Finite morphisms and integral morphisms are stable under base change and under composition.}
\face{Tape 2, side A}
\entree[prop]{30136}{8:22:16}{A morphism is finite if and only if it is integral and of finite type; an algebra $B$ is finite over $A$ if and only if it is finitely generated and integral.}
\entree[prop]{30201}{8:23:21}{An integral morphism is separated (being affine) and universally closed; a finite morphism is therefore proper.}
\entree[prop]{30267}{8:24:27}{Every closed immersion is finite, since $A/J$ is finite over $A$.}
\entree[thm]{30399}{8:26:39}{Lying-over (Cohen--Seidenberg): if $A\subset B$ is integral, every prime ideal of $A$ is induced by a prime ideal of $B$; this is the key point for integral morphisms to be closed.}
\entree[prop]{30462}{8:27:42}{For $k'/k$ algebraic and $X$ a $k$-scheme, the projection $X_{k'}\to X$ is closed.}
\session{3 July 1973}
\noindent\textit{Video 3/4 (\href{https://www.youtube.com/watch?v=Oqhr-Gt66es}{\texttt{Oqhr-Gt66es}}): tape 1, side A (8:28:48--9:16:43); tape 1, side B (9:16:43--10:05:05); tape 2 (10:05:05--10:53:45).}\par\medskip
\video{Oqhr-Gt66es}
\face{Tape 1, side A}
\entree[prop]{30706}{8:31:46}{A morphism is finite if and only if it is integral and of finite type; a morphism is integral if and only if it is affine and universally closed.}
\entree[prop]{30771}{8:32:51}{A projective morphism is proper; a proper morphism is separated and universally closed.}
\entree[prop]{30835}{8:33:55}{Closed immersions are finite, finite morphisms are projective, projective ones are quasi-projective.}
\entree[prop]{30898}{8:34:58}{Affine morphisms and quasi-compact immersions are quasi-affine; quasi-affine morphisms of finite type are quasi-projective.}
\entree[prop]{31028}{8:37:08}{A morphism is finite if and only if it is proper with finite fibres.}
\entree[prop]{31094}{8:38:14}{A morphism of schemes of finite type over $\mathbb{C}$ is proper if and only if the induced map between the analytic spaces of complex points is proper.}
\topic{Ample sheaves}
\entree[prop]{31422}{8:43:42}{For $X$ coherent and $L$ invertible, the following are equivalent: the $X_f$, $f\in\Gamma(X,L^{\otimes n})$, $n>0$, form a base of the topology; those that are affine form a base.}
\entree[prop]{31560}{8:46:00}{For $X$ coherent, these conditions are also equivalent to: for every quasi-coherent $F$ of finite type, $F\otimes L^{\otimes n}$ is generated by its sections for $n$ large; the same condition for quasi-coherent ideals of finite type.}
\entree[prop]{31756}{8:49:16}{If $X$ is coherent over an affine $Y=\mathrm{Spec}\,A$ and $L^{\otimes n}$ is very ample relative to $Y$ for some $n>0$, these conditions are satisfied; the converse holds if $X$ is of finite type over $Y$.}
\entree[def]{31946}{8:52:26}{On a coherent scheme, an invertible sheaf $L$ is ample if it satisfies the preceding equivalent conditions.}
\entree[def]{32012}{8:53:32}{$L$ is ample relative to $Y$ if there is a covering of $Y$ by affine open sets $Y_i$ such that $L$ restricted to $f^{-1}(Y_i)$ is ample; it is then so over every affine open set of $Y$.}
\entree[prop]{32075}{8:54:35}{If there is an invertible sheaf ample relative to $Y$, then $X$ is separated over $Y$.}
\entree[prop]{32204}{8:56:44}{For $f$ coherent and $Y$ quasi-compact, $L$ is ample relative to $Y$ if and only if, for every quasi-coherent module of finite type $F$ on $X$, $f^*f_*(F\otimes L^{\otimes n})\to F\otimes L^{\otimes n}$ is surjective for $n$ large.}
\entree[prop]{32333}{8:58:53}{If a power $L^{\otimes n}$ ($n>0$) is very ample relative to $Y$, then $L$ is ample; the converse holds if $f$ is of finite type and $Y$ quasi-compact.}
\entree[prop]{32398}{8:59:58}{For $f\colon X\to Y$, $g\colon Y\to Z$, $Z$ quasi-compact, $L$ ample (resp. very ample) relative to $f$ and $M$ ample (resp. very ample) relative to $g$, the sheaf $L\otimes f^*M^{\otimes n}$ is ample (resp. very ample) relative to $g\circ f$ for $n\geq n_0$.}
\entree[prop]{32732}{9:05:32}{If $X\subset\mathbb{P}^n_Y$ and $Y\subset\mathbb{P}^m_Z$, $X$ embeds into $\mathbb{P}^n_Z\times_Z\mathbb{P}^m_Z$, then by Segre into $\mathbb{P}^{(n+1)(m+1)-1}_Z$, the induced sheaf being $L\otimes f^*M$.}
\entree[prop]{32997}{9:09:57}{For $E$ quasi-coherent and $I$ invertible, $\mathbb{P}(E)\simeq\mathbb{P}(E\otimes I)$, the fundamental sheaf of $\mathbb{P}(E\otimes I)$ being that of $\mathbb{P}(E)$ twisted by the inverse image of $I$.}
\entree[cex]{33188}{9:13:08}{The notions of quasi-projective morphism and projective morphism are not local on the base.}
\topic{Quasi-affine morphisms}
\entree[def]{33252}{9:14:12}{A scheme is quasi-affine if it is isomorphic to a quasi-compact open set of an affine scheme; it is then separated and coherent.}
\entree[prop]{33316}{9:15:16}{Every morphism from a scheme $X$ to an affine scheme $\mathrm{Spec}\,B$ factors uniquely through $X\to\mathrm{Spec}\,\Gamma(X,\mathcal{O}_X)$.}
\face{Tape 1, side B}
\entree[prop]{33470}{9:17:50}{For $X$ coherent and $A=\Gamma(X,\mathcal{O}_X)$, the following are equivalent: $X$ is quasi-affine; $X\to\mathrm{Spec}\,A$ is an open immersion; $\mathcal{O}_X$ is ample; the $X_f$, $f\in A$, that are affine form a base of the topology; every quasi-coherent module is generated by its sections.}
\entree[prop]{33600}{9:20:00}{These conditions are also equivalent to the existence of an invertible sheaf $L$ such that $L$ and $L^{-1}$ are ample.}
\entree[def]{33740}{9:22:20}{$f\colon X\to Y$ is quasi-affine if there is a covering of $Y$ by affine open sets $Y_i$ such that $f^{-1}(Y_i)$ is quasi-affine; this is equivalent to $\mathcal{O}_X$ ample relative to $Y$.}
\entree[prop]{33808}{9:23:28}{A morphism is quasi-affine if and only if it is coherent and, for every quasi-coherent module $F$ on $X$, $f^*f_*F\to F$ is surjective.}
\entree[prop]{33870}{9:24:30}{Quasi-affine morphisms are stable under base change and under composition, with no quasi-compactness hypothesis.}
\entree[prop]{33935}{9:25:35}{An immersion is quasi-affine if and only if it is quasi-compact; if $g\circ f$ is quasi-affine and $g$ quasi-separated, then $f$ is quasi-affine.}
\entree[prop]{34060}{9:27:40}{A separated morphism of finite type with finite fibres is quasi-affine.}
\entree[prop]{34132}{9:28:52}{If $Y$ is coherent, such a morphism factors as a quasi-compact open immersion followed by a finite morphism (modern form of Zariski's Main Theorem); of finite type with finite fibres is called quasi-finite.}
\topic{Radicial morphisms}
\entree[prop]{34487}{9:34:47}{The following are equivalent: $f$ is universally injective (injective after any base change); each fibre has at most one point $x$ above $y$, with $\kappa(x)/\kappa(y)$ purely inseparable. One then says $f$ is radicial.}
\entree[prop]{34689}{9:38:09}{$f$ is radicial if and only if $X(K)\to Y(K)$ is injective for every field $K$, if and only if the diagonal $X\to X\times_YX$ is surjective.}
\entree[prop]{34883}{9:41:23}{Radicial morphisms are stable under base change and under composition; every monomorphism is radicial.}
\entree[prop]{35012}{9:43:32}{An integral, radicial and surjective morphism is a universal homeomorphism; conversely a universal homeomorphism is integral, radicial and surjective.}
\entree[ex]{35146}{9:45:46}{For $k'/k$ purely inseparable (for instance the perfect closure of $k$), $\mathrm{Spec}\,k'\to\mathrm{Spec}\,k$ is a universal homeomorphism, so $X_{k'}\to X$ is a homeomorphism.}
\entree[ex]{35209}{9:46:49}{The normalization of a curve with a cusp is integral, radicial and surjective, hence a universal homeomorphism.}
\entree[disc]{35272}{9:47:52}{A nodal curve minus one preimage of the double point in its normalization is a homeomorphism onto the curve but not a universal homeomorphism, as base change to the normalization shows.}
\entree[cex]{35338}{9:48:58}{The normalization of a curve with an ordinary double point, with one of the two points above the double point removed, is a homeomorphism that is not universal.}
\topic{Open and generalizing morphisms}
\entree[def]{35659}{9:54:19}{$f$ is open if the underlying map is, universally open if it remains so after any base change.}
\entree[prop]{35806}{9:56:46}{Open morphisms are stable under composition but not under base change; universally open morphisms are stable under both.}
\entree[ex]{35806}{9:56:46}{The normalization of a curve with an ordinary double point, with one of the two points above it removed, is bijective and a homeomorphism, hence open, but not universally open.}
\entree[prop]{35869}{9:57:49}{An open morphism is generalizing: if $f(x)=y$, every generization of $y$ lifts to a generization of $x$, that is, $\mathrm{Spec}\,\mathcal{O}_{X,x}\to\mathrm{Spec}\,\mathcal{O}_{Y,y}$ is surjective.}
\face{Tape 2}
\entree[prop]{36375}{10:06:15}{A filtered inductive limit of rings is zero if and only if one of these rings is zero.}
\entree[prop]{36440}{10:07:20}{A universally open morphism is universally generalizing; generalizing morphisms are stable under composition, universally generalizing ones under composition and base change.}
\entree[cex]{36503}{10:08:23}{For $\eta$ a maximal point of $X$, $\mathrm{Spec}\,\kappa(\eta)\to X$ is generalizing; it is not open in general (for example for the affine line over a field).}
\entree[prop]{36565}{10:09:25}{If $f$ is locally of finite presentation, $f$ is open if and only if it is generalizing, and universally open if and only if it is universally generalizing.}
\entree[prop]{36696}{10:11:36}{The image of a morphism of finite type between noetherian schemes is constructible, that is, a finite union of locally closed subsets.}
\entree[prop]{36760}{10:12:40}{A constructible subset of a scheme is open if and only if it is stable under generization.}
\entree[prop]{36951}{10:15:51}{If $f$ is locally of finite presentation, $f$ is open at $x$ (the image of every neighbourhood of $x$ is a neighbourhood of $f(x)$) if and only if it is generalizing at $x$.}
\entree[prop]{37016}{10:16:56}{Every scheme over a field $k$ is universally open over $k$.}
\topic{Flat and faithfully flat morphisms}
\entree[def]{37143}{10:19:03}{A module $M$ over $A$ is flat if the functor $M\otimes_A-$ is exact, faithfully flat if moreover it is faithful, that is, if $M\otimes_AN=0$ implies $N=0$.}
\entree[prop]{37209}{10:20:09}{A flat module $M$ is faithfully flat if and only if $u$ is injective (resp. surjective, bijective) as soon as $M\otimes u$ is.}
\entree[prop]{37402}{10:23:22}{A module $M$ is faithfully flat if and only if it is flat and $M/\mathfrak mM\neq0$ for every maximal ideal $\mathfrak m$ of $A$.}
\entree[prop]{37726}{10:28:46}{An $A$-algebra $B$ is faithfully flat if and only if it is flat and $\mathrm{Spec}\,B\to\mathrm{Spec}\,A$ is surjective.}
\entree[prop]{37918}{10:31:58}{If $A\to B$ is a local homomorphism of local rings, $B$ flat over $A$ is equivalent to $B$ faithfully flat; the map of spectra is then surjective.}
\entree[prop]{38052}{10:34:12}{Flatness of an algebra is stable under localization: if $B$ is flat over $A$, $B[T^{-1}]$ is flat over $A[S^{-1}]$; the local rings of $B$ are therefore flat over those of $A$.}
\entree[prop]{38114}{10:35:14}{If $B$ is flat over $A$, $\mathrm{Spec}\,B\to\mathrm{Spec}\,A$ is generalizing.}
\entree[prop]{38243}{10:37:23}{A module on a ringed space is flat if and only if its fibres are flat as modules over the fibres of the sheaf of rings.}
\entree[prop]{38374}{10:39:34}{A quasi-coherent module $\widetilde{M}$ on $\mathrm{Spec}\,A$ is flat if and only if the $A$-module $M$ is flat.}
\entree[prop]{38502}{10:41:42}{A module $M$ on $X$ is flat relative to $f\colon X\to Y$ if each fibre $M_x$ is flat over $\mathcal{O}_{Y,f(x)}$.}
\entree[def]{38630}{10:43:50}{$f$ is flat if $\mathcal{O}_X$ is flat relative to $Y$; a notion local on $X$ and $Y$, equivalent to $\mathcal{O}_{X,x}$ flat over $\mathcal{O}_{Y,f(x)}$ for all $x$; for $X,Y$ affine, to $B$ flat over $A$.}
\entree[prop]{38767}{10:46:07}{A flat morphism is generalizing, hence universally generalizing; flat morphisms contain the isomorphisms and are stable under composition and base change.}
\entree[def]{38963}{10:49:23}{$f$ is faithfully flat if it is flat and surjective; these morphisms are stable under composition and base change.}
\entree[prop]{39030}{10:50:30}{A morphism $f$ is faithfully flat if and only if the inverse image functor $f^*$, from modules on $Y$ to modules on $X$, is exact and faithful.}
\entree[prop]{39165}{10:52:45}{If $f$ is flat, the closure of the image of every irreducible component of $X$ is an irreducible component of $Y$.}
\session{9 July 1973}
\noindent\textit{Video 4/4 (\href{https://www.youtube.com/watch?v=vuys5SNs2eQ}{\texttt{vuys5SNs2eQ}}): tape 1 (0:00:00--0:48:32); tape 2, side A (0:48:32--1:38:09); tape 2, side B (1:38:09--2:27:17); tape 3 (2:27:17--3:15:45).}\par\medskip
\video{vuys5SNs2eQ}
\face{Tape 1}
\topic{Flatness criteria}
\entree[prop]{2}{0:00:02}{A generalizing morphism sends the maximal points of $X$ (generic points of the irreducible components) to maximal points of $Y$.}
\entree[prop]{131}{0:02:11}{A morphism $f\colon X\to Y$ is generalizing if and only if, for every irreducible closed set $Y'$ of $Y$, every irreducible component of $f^{-1}(Y')$ dominates $Y'$.}
\entree[prop]{131}{0:02:11}{For a field extension $k'/k$ and a $k$-scheme $X$, every irreducible component of $X_{k'}$ dominates an irreducible component of $X$.}
\entree[disc]{259}{0:04:19}{Asked for a topological characterisation of flatness, the lecturer gives cases over a regular curve or Dedekind base, and for reduced schemes, flatness means components dominate.}
\entree[prop]{327}{0:05:27}{Over a principal ideal ring, a module is flat if and only if it is torsion-free.}
\entree[prop]{390}{0:06:30}{Over an integral domain $A$, every flat $A$-module is torsion-free.}
\entree[prop]{455}{0:07:35}{Over a Dedekind ring, every torsion-free module is flat (one reduces by localization to a discrete valuation ring, hence a principal ring).}
\entree[prop]{520}{0:08:40}{A filtered inductive limit of flat modules is flat.}
\entree[prop]{584}{0:09:44}{Let $Y$ be integral, regular of dimension $1$, and $X$ reduced over $Y$: $X$ is flat over $Y$ if and only if every irreducible component of $X$ dominates $Y$.}
\entree[prop]{587}{0:09:47}{Without assuming $X$ reduced (locally noetherian), $X$ is flat over $Y$ if and only if all the associated points of $X$ map to the generic point of $Y$.}
\entree[prop]{712}{0:11:52}{If $Y$ is regular, $X$ Cohen--Macaulay and $f$ locally of finite type, $f$ is flat if and only if $\dim_x X=\dim_{f(x)}Y+\dim_x X_{f(x)}$ for all $x$.}
\entree[prop]{900}{0:15:00}{If $Y$ is noetherian normal, $f$ of finite presentation, universally open with geometrically reduced fibres, then $f$ is flat.}
\topic{Faithfully flat base change}
\entree[prop]{1166}{0:19:26}{If $f\colon X\to Y$ is faithfully flat, the map $Y'\mapsto f^{-1}(Y')$ from subschemes of $Y$ to those of $X$ is injective.}
\entree[prop]{1230}{0:20:30}{If $f$ is faithfully flat, the inverse image of quasi-coherent modules is exact and faithful; two submodules of a module that become equal after inverse image are equal.}
\entree[prop]{1299}{0:21:39}{If $S'\to S$ is faithfully flat, the base change functor from $S$-schemes to $S'$-schemes is faithful: $\mathrm{Hom}_S(X,Y)\to\mathrm{Hom}_{S'}(X',Y')$ is injective.}
\entree[prop]{1625}{0:27:05}{For a commutative square of ringed spaces, there is a canonical morphism $g^*f_*F\to f'_*F'$, defined by adjunction.}
\entree[prop]{1827}{0:30:27}{If $f$ is coherent (quasi-compact, quasi-separated), $g$ flat and $F$ quasi-coherent, $g^*f_*F\to f'_*g'^*F$ is an isomorphism; likewise for $R^if_*$.}
\entree[prop]{2099}{0:34:59}{If $A\to B$ is a local homomorphism of local rings and $B$ is flat over $A$, then $B$ is faithfully flat and $\mathrm{Spec}\,B\to\mathrm{Spec}\,A$ is surjective.}
\entree[prop]{2164}{0:36:04}{A flat morphism is generalizing, hence universally generalizing.}
\entree[prop]{2356}{0:39:16}{The canonical morphism $g^*f_*F\to f'_*F'$ is an isomorphism with no flatness hypothesis if $f$ is affine.}
\entree[def]{2620}{0:43:40}{$f$ is radicial if it is universally injective, that is, remains injective after any base change.}
\entree[prop]{2620}{0:43:40}{If $g\colon S'\to S$ is surjective, $f$ is surjective (resp. radicial) if and only if $f'$ is.}
\entree[prop]{2686}{0:44:46}{A morphism is radicial if and only if its diagonal morphism is surjective.}
\entree[prop]{2750}{0:45:50}{If $g$ is surjective and quasi-compact, $f$ is quasi-compact (resp. quasi-separated, coherent) if and only if $f'$ is.}
\face{Tape 2, side A}
\topic{Descent of properties along faithfully flat morphisms}
\entree[prop]{2913}{0:48:33}{A morphism $f$ is quasi-separated if and only if its diagonal morphism $\Delta_f$ is quasi-compact; it is coherent if it is quasi-compact and quasi-separated.}
\entree[disc]{2913}{0:48:33}{A listener points out a mistake in a reduction step; the lecturer concedes that one can reduce to $Y$ affine but not to $Y'$ affine, and repairs the argument using quasi-compactness.}
\entree[prop]{3040}{0:50:40}{If $g\colon S'\to S$ is faithfully flat and quasi-compact, $f$ is locally of finite type (resp. locally of finite presentation) if and only if $f'$ is.}
\entree[prop]{3104}{0:51:44}{A morphism is of finite type if it is locally of finite type and quasi-compact, of finite presentation if it is locally of finite presentation and coherent; these properties also descend along $g$ faithfully flat quasi-compact.}
\entree[prop]{3297}{0:54:57}{If $A\to A'$ is faithfully flat, an $A$-algebra $B$ is finitely generated (resp. finitely presented) if and only if $B\otimes_A A'$ is over $A'$.}
\entree[prop]{3491}{0:58:11}{A quasi-coherent $\mathcal{O}_X$-module $F$ is of finite type (resp. finite presentation) if and only if $F'$ is, for $g$ faithfully flat quasi-compact.}
\entree[prop]{3556}{0:59:16}{Under the same hypotheses, $F$ is locally free of finite type, of given rank $n$, or invertible, if and only if $F'$ is.}
\entree[prop]{3624}{1:00:24}{For $Y'\to Y$ faithfully flat (without quasi-compactness), a quasi-coherent module $F$ on $X$ is flat over $Y$ if and only if $F'$ is flat over $Y'$; in particular $f$ is flat if and only if $f'$ is.}
\entree[prop]{3687}{1:01:27}{If $A\to A'$ is faithfully flat, an $A$-module $M$ is flat if and only if $M\otimes_A A'$ is flat over $A'$; likewise for faithful flatness.}
\entree[prop]{3752}{1:02:32}{A quasi-coherent module is locally free of finite type if and only if it is of finite presentation and flat; this property descends, with the rank.}
\topic{Topological descent}
\entree[prop]{3885}{1:04:45}{For $g\colon S'\to S$ and $Z\subset S$, one always has $\overline{g^{-1}(Z)}\subset g^{-1}(\overline{Z})$; equality for all $Z$ characterizes open morphisms.}
\entree[cex]{3950}{1:05:50}{The inclusion of the generic point of an integral scheme $S$ into $S$ is flat, but in general not open.}
\entree[prop]{4013}{1:06:53}{If $g\colon S'\to S$ is flat (or merely generalizing) and $f\colon Z\to S$ quasi-compact, then $g^{-1}(\overline{f(Z)})=\overline{g^{-1}(f(Z))}$.}
\entree[prop]{4142}{1:09:02}{If $g$ is generalizing and $f\colon Z\to S$ quasi-compact dominant, then $f'\colon Z'\to S'$ is dominant; the converse holds if $g$ is surjective.}
\entree[prop]{4464}{1:14:24}{A quasi-compact morphism $f\colon Z\to S$ is dominant if and only if its fibre at each maximal point of $S$ is nonempty.}
\entree[prop]{4795}{1:19:55}{If $g\colon S'\to S$ is flat (or generalizing) and quasi-compact, the topology of $g(S')$ induced by $S$ is the quotient topology of $S'$; if $g$ is moreover surjective, $g$ is submersive.}
\entree[prop]{5252}{1:27:32}{A faithfully flat quasi-compact morphism is universally submersive.}
\entree[prop]{5323}{1:28:43}{For $g$ faithfully flat quasi-compact: if $f'$ is open (resp. closed, generalizing, a homeomorphism), so is $f$; likewise for the corresponding universal properties.}
\entree[prop]{5590}{1:33:10}{If $g$ is surjective, $f$ is surjective if and only if $f'$ is, and $f'$ injective implies $f$ injective.}
\entree[prop]{5811}{1:36:51}{For $g$ faithfully flat quasi-compact, $f$ is separated (resp. proper) if and only if $f'$ is.}
\face{Tape 2, side B}
\entree[prop]{5892}{1:38:12}{For $g$ faithfully flat quasi-compact, $f$ is an isomorphism if and only if $f'$ is (flatness is necessary: a universal topological property does not suffice).}
\entree[prop]{6022}{1:40:22}{The base change functor along a faithfully flat quasi-compact morphism is faithful and conservative: an arrow over $S$ that becomes an isomorphism was one.}
\entree[prop]{6088}{1:41:28}{A morphism $f$ is a monomorphism if and only if $\Delta_f$ is an isomorphism; consequently $f$ is a monomorphism if and only if $f'$ is (for $g$ faithfully flat quasi-compact).}
\topic{Faithfully flat descent of modules}
\entree[prop]{6217}{1:43:37}{If $g$ is faithfully flat, $M\mapsto M'$ is exact, faithful and conservative on quasi-coherent modules; $\mathrm{Hom}(M,N)\to\mathrm{Hom}(M',N')$ is injective.}
\entree[prop]{6408}{1:46:48}{With $S''=S'\times_S S'$ and the projections $p_1,p_2$, the sequence $\mathrm{Hom}(M,N)\to\mathrm{Hom}(M',N')\rightrightarrows\mathrm{Hom}(M'',N'')$ is a complex; if $g$ is faithfully flat, the first arrow is injective.}
\entree[def]{6470}{1:47:50}{$g$ is a descent morphism for quasi-coherent modules (descent of rank $1$) if this sequence is exact for all quasi-coherent $M$, $N$.}
\entree[thm]{6533}{1:48:53}{A faithfully flat quasi-compact morphism $g$ is a descent morphism, and even of effective descent, for quasi-coherent modules.}
\entree[def]{6663}{1:51:03}{A gluing datum on a module $M'$ of $S'$ is an isomorphism $c\colon p_1^*M'\simeq p_2^*M'$; $M\mapsto(M',c)$ defines a functor from modules of $S$ to modules of $S'$ equipped with gluing.}
\entree[prop]{6732}{1:52:12}{This functor is fully faithful if and only if $g$ is a descent morphism.}
\entree[prop]{6865}{1:54:25}{If $(M',c)$ comes from a module on $S$, then, on $S'''=S'\times_S S'\times_S S'$, one has $q_{13}^*c=q_{23}^*c\circ q_{12}^*c$ (cocycle condition).}
\entree[prop]{7642}{2:07:22}{The morphism $g$ is of descent of rank $0$ if the functor to modules with gluing is faithful, of rank $1$ if it is fully faithful; a faithfully flat morphism is of descent of rank $0$.}
\entree[disc]{7705}{2:08:25}{A discussion of terminology: 'clutching' versus 'gluing' data, with a remark that clutching is standard in vector bundle theory and K-theory, as in bundles on hemispheres.}
\entree[def]{7841}{2:10:41}{A gluing datum satisfying the cocycle condition is a descent datum.}
\entree[def]{7906}{2:11:46}{$g$ is of descent of rank $2$ (effective descent) if the functor $\mathrm{Mod}(S)\to\mathrm{Mod}(S')$ with descent data is an equivalence of categories.}
\entree[thm]{7970}{2:12:50}{If $g$ is faithfully flat and quasi-compact, $g^*$ induces an equivalence between quasi-coherent modules on $S$ and quasi-coherent modules on $S'$ equipped with a descent datum.}
\entree[read]{8175}{2:16:15}{The general nonsense of descent is in Giraud’s memoir, which he calls pretty hard reading. \cite{girauddescent}}
\entree[prop]{8629}{2:23:49}{If $g$ is of descent of rank $1$ over every open set, effectiveness of a descent datum is local on $S$: the modules $M_\alpha$ on an open covering glue by uniqueness.}
\entree[prop]{8760}{2:26:00}{For an open covering $(S_i)$ of $S$, the morphism $\coprod_i S_i\to S$ is of effective descent for quasi-coherent modules.}
\face{Tape 3}
\entree[prop]{8840}{2:27:20}{If a descent (resp. effective descent) morphism $\bar S\to S$ factors through $S'\to S$, then $S'\to S$ is of descent (resp. effective descent) for quasi-coherent modules.}
\entree[prop]{8903}{2:28:23}{A morphism $S'\to S$ admitting a section is of effective descent for quasi-coherent modules.}
\entree[thm]{9876}{2:44:36}{If $A\to A'$ is faithfully flat, the functor $M\mapsto M\otimes_A A'$ is an equivalence from $A$-modules to $A'$-modules equipped with a descent datum.}
\entree[prop]{10004}{2:46:44}{The module descended from an $A'$-module $M'$ with descent datum is the kernel of $M'\rightrightarrows M''$, with $M''=M'\otimes_A A'$; one then has $M\otimes_A A'\simeq M'$.}
\topic{Descent of subobjects}
\entree[prop]{10195}{2:49:55}{If $g\colon S'\to S$ is faithfully flat, $G\mapsto G'$ is injective from quasi-coherent submodules of $F$ to those of $F'$.}
\entree[prop]{10195}{2:49:55}{If $g$ is moreover quasi-compact, the sequence of quasi-coherent submodules of $F$, $F'$, $F''$ is exact: a submodule of $F'$ descends if and only if its two inverse images in $F''$ coincide.}
\entree[prop]{10517}{2:55:17}{If $g$ is surjective, a subset of $S'$ comes from a subset of $S$ if and only if its two inverse images in $S''$ coincide.}
\entree[prop]{10644}{2:57:24}{If $g$ is faithfully flat and quasi-compact, the open (resp. closed) sets of $S$ are identified with the open (resp. closed) sets of $S'$ having the same inverse image in $S''$ under $p_1$ and $p_2$.}
\entree[prop]{10707}{2:58:27}{The same exactness holds for open subschemes, for closed subschemes (quasi-coherent ideals), and for retrocompact subschemes (quasi-compact immersions).}
\entree[prop]{10835}{3:00:35}{If $g$ is faithfully flat quasi-compact, $T\mapsto g^{-1}(T)$ is injective on subschemes of $S$.}
\entree[disc]{10903}{3:01:43}{Subobjects versus subschemes: the lecturer gives monomorphisms that are not immersions, such as a point removed and then added back, or an infinitesimal neighbourhood, or a curve with a double point.}
\entree[prop]{10967}{3:02:47}{Likewise for subobjects (schemes equipped with a monomorphism to $S$): $T'_1\subset T'_2$ implies $T_1\subset T_2$, by descent of isomorphisms.}
\entree[cex]{10967}{3:02:47}{The morphism $(\mathbb{A}^1\setminus\{0\})\sqcup\{0\}\to\mathbb{A}^1$ is a monomorphism that is not an immersion.}
\entree[cex]{11096}{3:04:56}{The normalization of a curve with an ordinary double point, with one of the two points above the double point removed, is a bijective monomorphism that is not an immersion.}
\entree[prop]{11349}{3:09:09}{If $g$ is faithfully flat quasi-compact, the retrocompact locally closed subsets of $S$ are identified with those of $S'$ having the same inverse image in $S''$.}
\entree[thm]{11483}{3:11:23}{If $g$ is faithfully flat quasi-compact, the sequence $\mathrm{Hom}_S(X,Y)\to\mathrm{Hom}_{S'}(X',Y')\rightrightarrows{}$ $\mathrm{Hom}_{S''}(X'',Y'')$ is exact for all $S$-schemes $X$, $Y$ (descent of rank $1$ for schemes).}
\session{10 July 1973}
\noindent\textit{Video 4/4 (\href{https://www.youtube.com/watch?v=vuys5SNs2eQ}{\texttt{vuys5SNs2eQ}}): tape 1 (3:15:45--4:03:41); tape 2 (4:03:41--4:52:26); tape 1 (4:52:26--5:41:49).}\par\medskip
\video{vuys5SNs2eQ}
\face{Tape 1 (3:15:45--4:03:41)}
\topic{Descent of morphisms and of affine schemes}
\entree[prop]{11824}{3:17:04}{If $g$ is faithfully flat quasi-compact and $Y$ quasi-separated over $S$, a morphism $X'\to Y'$ compatible with the descent data descends to an $S$-morphism $X\to Y$ (via its graph, a retrocompact subscheme of $X'\times_{S'}Y'$).}
\entree[prop]{12082}{3:21:22}{If $S'\to S$ is faithfully flat quasi-compact, the functor from $S$-schemes to $S'$-schemes with descent data is fully faithful.}
\entree[prop]{12211}{3:23:31}{If $S'\to S$ is faithfully flat quasi-compact, the category of affine $S$-schemes is equivalent to that of affine $S'$-schemes with descent data.}
\entree[prop]{12211}{3:23:31}{If $g$ is faithfully flat quasi-compact, $f$ is an open (resp. closed, resp. quasi-compact) immersion if and only if $f'$ is.}
\entree[disc]{12537}{3:28:57}{Whether a filtering union of locally closed pieces is locally closed in the target is first left uncertain, then refuted with a dense subset of the line.}
\entree[cex]{12606}{3:30:06}{A space that is a filtered union of open sets each mapped homeomorphically onto a locally closed subset of $Y$ need not map onto a locally closed subset: a dense subset of a line (for schemes, the disjoint union of all the closed points of a line mapped into it).}
\entree[q]{12671}{3:31:11}{Do immersions that are not quasi-compact descend along a faithfully flat quasi-compact morphism? Believed so, but not proved here; the attempt through a filtered union of quasi-compact open sets fails.}
\topic{Constructible sets and Chevalley's theorem}
\entree[prop]{12865}{3:34:25}{If $f\colon Z\to S$ is quasi-compact, $s\in\overline{f(Z)}$ if and only if $s$ has a generization in $f(Z)$.}
\entree[prop]{12993}{3:36:33}{If $f\colon Z\to S$ is quasi-compact, $s$ is an interior point of the complement of $f(Z)$ if and only if all generizations of $s$ lie outside $f(Z)$.}
\entree[thm]{13186}{3:39:46}{If $f\colon Z\to S$ is of finite presentation, $f(Z)$ is locally constructible in $S$ (Chevalley's theorem); the proof reduces to the noetherian case.}
\entree[prop]{13380}{3:43:00}{A generalizing morphism that is locally of finite presentation is open.}
\topic{Descent of properties of morphisms}
\entree[cex]{13514}{3:45:14}{Being projective (resp. quasi-projective) does not descend: it is not local on the base for the Zariski topology.}
\entree[prop]{14094}{3:54:54}{For $g$ faithfully flat quasi-compact, quasi-coherent quotients of $M$ descend like submodules, which gives descent of closed subschemes.}
\entree[cex]{14220}{3:57:00}{For schemes, the functor to schemes with descent data is fully faithful but not an equivalence in general: a descent datum is not always effective.}
\entree[prop]{14416}{4:00:16}{For $g$ of descent of rank $1$ (resp. $2$) for quasi-coherent modules, it is so for quasi-coherent algebras; product, associativity, commutativity and unit descend.}
\face{Tape 2}
\entree[prop]{14623}{4:03:43}{If $g$ is faithfully flat quasi-compact, $f$ is affine (resp. quasi-affine, integral, finite, quasi-finite) if and only if $f'$ is.}
\entree[prop]{14816}{4:06:56}{A coherent morphism $f\colon X\to Y$ is quasi-affine if $X\to\mathrm{Spec}\,f_*\mathcal{O}_X$ is an open immersion.}
\entree[prop]{14946}{4:09:06}{If $A\to A'$ is faithfully flat and $B$ an $A$-algebra, $B$ is integral (resp. finite) over $A$ if and only if $B\otimes_A A'$ is over $A'$.}
\entree[prop]{15009}{4:10:09}{An algebra $B$ is integral over $A$ if and only if its finitely generated $A$-subalgebras are finite over $A$.}
\entree[def]{15077}{4:11:17}{A morphism is quasi-finite if it is of finite type with finite fibres.}
\entree[prop]{15139}{4:12:19}{If $g$ is faithfully flat quasi-compact, an invertible sheaf $L$ on $X$ is ample (resp. very ample) relative to $f$ if and only if $L'$ is relative to $f'$.}
\entree[def]{15264}{4:14:24}{$L$ is very ample relative to $f$ if $f^*f_*L\to L$ is surjective and the induced morphism $X\to\mathbf{P}(f_*L)$ is an immersion.}
\topic{Effective descent for schemes}
\entree[prop]{15735}{4:22:15}{A gluing datum on $X'$ is a descent datum if and only if $X''\rightrightarrows X'$ (over $p_1$, $p_2$, with cartesian squares) is an equivalence relation.}
\entree[prop]{16055}{4:27:35}{A descent datum on $X'$ is effective if and only if the equivalence relation $X''\rightrightarrows X'$ is universally effective: the quotient $X$ exists, $X''=X'\times_X X'$, and this remains true after any base change on $X$.}
\entree[disc]{16247}{4:30:47}{A listener notes a missing hypothesis; the lecturer agrees and corrects: the equivalence relation must be effective and universal, so that the quotient stays a quotient after base change.}
\entree[prop]{16510}{4:35:10}{An open set of $X'$ is stable under the descent datum if and only if its underlying subset has the same inverse image in $X''$ under $q_1$ and $q_2$ (set-theoretic condition).}
\entree[prop]{16640}{4:37:20}{If $X'$ is a union of stable open sets $X'_i$, the descent datum on $X'$ is effective if and only if the induced data on the $X'_i$ are.}
\face{Tape 1 (4:52:26--5:41:49)}
\entree[read]{17357}{4:49:17}{For effective equivalence relations he points to Demazure. \cite{demazure}}
\entree[prop]{17552}{4:52:32}{If $g$ is faithfully flat quasi-compact and $X'$ is covered by open sets $X'_i$ stable under the descent datum and affine over $S'$, the descent datum on $X'$ is effective.}
\entree[prop]{17754}{4:55:54}{A faithfully flat, quasi-compact and radicial morphism is a universal homeomorphism.}
\entree[prop]{18155}{5:02:35}{If $g$ is faithfully flat quasi-compact radicial, every subset of $X'$ is stable under a descent datum, which is therefore effective; for example for a purely inseparable field extension.}
\entree[cex]{18219}{5:03:39}{For a separable extension $k'/k$ of degree $2$, there is a proper smooth scheme of dimension $3$ over $k'$, not projective, equipped with a non-effective descent datum.}
\entree[prop]{18281}{5:04:41}{If $g$ is faithfully flat quasi-compact and $L'$ is ample relative to $S'$ with compatible descent datum on $(X',L')$, the descent datum on $X'$ is effective.}
\entree[rem]{18420}{5:07:00}{If $X'$ is of finite type over $S'$, an ample sheaf $L'$ has a very ample power, so the proof reduces to the very ample case.}
\entree[prop]{18554}{5:09:14}{If $f'$ is coherent and $(X',L')$ carries a descent datum, then $E'=f'_*L'$ inherits a descent datum, because direct image commutes with flat base change.}
\entree[thm]{18808}{5:13:28}{A faithfully flat quasi-compact morphism is of descent of rank $2$ for pairs $(X',L')$, $L'$ very ample relative to the base; the descended $L$ is very ample.}
\entree[prop]{18872}{5:14:32}{If $S'\to S$ is finite locally free of nowhere zero rank, every quasi-projective $S'$-scheme equipped with a descent datum relative to $S$ is effective, and the descended scheme is quasi-projective over $S$.}
\entree[prop]{19389}{5:23:09}{For a locally free algebra $B$ over $\mathcal{O}_T$, the norm $N\colon B\to\mathcal{O}_T$ (determinant of multiplication) is multiplicative, whence a norm of invertible sheaves from $T'=\mathrm{Spec}\,B$ to $T$.}
\entree[prop]{19516}{5:25:16}{Ampleness of an invertible sheaf is preserved under inverse image and under passing to the norm.}
\entree[prop]{19653}{5:27:33}{If $X$ is a scheme over a field $k$ and $X\otimes_k k'$ is quasi-projective (resp. projective) for some extension $k'$, then $X$ is quasi-projective (resp. projective).}
\entree[prop]{19843}{5:30:43}{Finite extension principle: a property stable under base change and passing to the limit, true after some extension of $k$, is true after a finite extension of $k$.}
\entree[prop]{19909}{5:31:49}{Two $k$-schemes of finite presentation that are isomorphic over an extension of $k$ are isomorphic over a finite extension of $k$.}
\topic{Descent and group actions}
\entree[prop]{20176}{5:36:16}{If $G$ acts on $S'$ with $S'\times G\simeq S'\times_S S'$, a descent datum on an $S'$-object $X'$ is equivalent to an action of $G$ on $X'$ compatible with the projection to $S'$.}
\entree[prop]{20242}{5:37:22}{For $G$ the constant group defined by a group $\Gamma$, a descent datum is equivalent to an action of $\Gamma$ on $X'$ compatible with that on $S'$.}
\entree[prop]{20367}{5:39:27}{If $S'\to S$ is of descent of rank $1$, $X\mapsto X'$ equipped with the action of $\Gamma$ is a fully faithful functor.}
\session{11 July 1973}
\noindent\textit{Video 2/2 of the groups course (\href{https://www.youtube.com/watch?v=3ULFIl_2Wis}{\texttt{3ULFIl\_2Wis}}): tape 2, side A (5:30:34--6:18:36); tape 2, side B (6:18:36--7:01:38).}\par\medskip
\video{3ULFIl_2Wis}
\face{Tape 2, side A}
\noindent\textit{Recorded on video 2/2 of the groups course (\href{https://www.youtube.com/watch?v=3ULFIl_2Wis}{\texttt{3ULFIl\_2Wis}}); by its content it belongs to the algebraic geometry course.}\par\smallskip
\entree[prop]{19839}{5:30:39}{If $S'\to S$ is a torsor under a group scheme $G$, a descent datum on an $S'$-scheme $X'$ is equivalent to an action of $G$ on $X'$ compatible with that on $S'$.}
\entree[ex]{19904}{5:31:44}{For $k'/k$ finite Galois with group $\Gamma$, $\mathrm{Spec}\,k'\to\mathrm{Spec}\,k$ is a torsor under the constant group $\Gamma$.}
\entree[prop]{19972}{5:32:52}{For $k'/k$ finite Galois with group $\Gamma$ and $X$, $Y$ $k$-schemes, $\mathrm{Hom}_k(X,Y)$ is the set of $\Gamma$-invariant $k'$-morphisms $X'\to Y'$.}
\entree[prop]{20037}{5:33:57}{A descent datum on a $k'$-scheme $X'$ relative to $k'/k$ is equivalent to an action of $\Gamma$ on $X'$ compatible with that on $k'$; it is effective if $X'$ is quasi-projective.}
\entree[prop]{20102}{5:35:02}{For $k'/k$ infinite Galois, $\Gamma_k=\varprojlim\Gamma_{i,k}$ (finite constant groups of the finite Galois subextensions) is an affine group scheme, and $\mathrm{Spec}\,k'$ is a torsor under $\Gamma_k$.}
\entree[prop]{20167}{5:36:07}{One has $\Gamma_k(k)=\Gamma$, and a descent datum on $X'$ relative to $k'/k$ is equivalent to an action of $\Gamma_k$ on $X'$ compatible with that on $k'$.}
\entree[prop]{20447}{5:40:47}{An action of $\Gamma_k$ on $X'$ is determined by the induced action of the abstract group $\Gamma$, the morphism from the constant group $\Gamma$ to $\Gamma_k$ being schematically dominant.}
\entree[prop]{20510}{5:41:50}{For arbitrary Galois $k'/k$, $\mathrm{Hom}_k(X,Y)$ is the set of $k'$-morphisms $X'\to Y'$ commuting with the action of $\Gamma$.}
\topic{Descent for families of morphisms}
\entree[def]{20702}{5:45:02}{Descent of degree $0$ for $(S'_i\to S)$: a morphism $X\to Y$ over $S$ is determined by its inverse images on the $S'_i$; degree $1$: morphisms over the $S'_i$ compatible on the $S'_i\times_SS'_j$ come from a morphism over $S$.}
\entree[def]{20768}{5:46:08}{Descent of degree $2$: objects $X'_i$ over the $S'_i$, equipped with isomorphisms over the $S'_i\times_SS'_j$ satisfying the cocycle condition on triples, come from an object over $S$ (effective descent).}
\entree[prop]{20831}{5:47:11}{For a family $(S'_i\to S)$ and $S'=\coprod S'_i$, one has $S'\times_SS'=\coprod S'_i\times_SS'_j$ and the $S'$-schemes are the families of $S'_i$-schemes; descent for the family therefore reduces to that of $S'\to S$.}
\entree[prop]{21089}{5:51:29}{A gluing datum on $X'/S'$ is equivalent to the data, for every $S$-scheme $T$ and every pair $u,v\colon T\to S'$, of an isomorphism $u^*X'\simeq v^*X'$ functorial in $T$; the cocycle condition means that these isomorphisms are transitive.}
\entree[prop]{21346}{5:55:46}{The objects of $S'$ equipped with a descent datum form the category of cartesian sections of the fibred category of relative schemes over the full subcategory of $S$-schemes that factor through $S'$.}
\entree[prop]{21411}{5:56:51}{For a family $(S'_i\to S)$, this category depends only on the sieve generated by the family.}
\entree[prop]{21606}{6:00:06}{For a faithfully flat quasi-compact base change $S'\to S$, being an isomorphism, an immersion or proper are properties of morphisms that descend.}
\entree[prop]{21668}{6:01:08}{Morphisms $X_i\to Y_i$ over an open covering $(S_i)$ of $S$, compatible on the intersections, glue to a unique morphism $X\to Y$; the same holds for a faithfully flat quasi-compact morphism $S'\to S$.}
\entree[prop]{21730}{6:02:10}{Faithfully flat quasi-compact descent is effective for quasi-coherent modules and affine relative schemes, but not for arbitrary relative schemes.}
\topic{Grothendieck topologies: fpqc, fppf, \'etale}
\entree[cex]{21857}{6:04:17}{The morphism $\mathbb{G}_m\to\mathbb{G}_m$, $\lambda\mapsto\lambda^n$ ($n\geq2$) is faithfully flat quasi-compact, with equivalence relation defined by $\mu_n$, but has no section locally for the Zariski topology.}
\entree[prop]{22049}{6:07:29}{A faithfully flat quasi-compact morphism $T\to S$ has a section after the base change $T\to S$ itself, hence locally for the fpqc topology.}
\entree[def]{22241}{6:10:41}{A family $(S'_i\to S)$ is covering for the fpqc topology if $S$ is a union of open sets $U$ for which there exist finitely many flat quasi-compact morphisms $U'_\alpha\to U$, with images covering $U$, each factoring through some $S'_i$.}
\entree[prop]{22369}{6:12:49}{If faithfully flat quasi-compact morphisms and open coverings are of descent of degree $0$, $1$ or $2$ for a category of objects, so are the fpqc covering families.}
\entree[prop]{22433}{6:13:53}{fpqc covering families are stable under composition and base change: the category of schemes is a site for this topology.}
\entree[def]{22497}{6:14:57}{A presheaf $F$ on schemes is an fpqc sheaf if, for every covering family $(S_i\to S)$, $F(S)\to\prod F(S_i)\rightrightarrows\prod F(S_i\times_SS_j)$ is exact.}
\entree[prop]{22562}{6:16:02}{A presheaf is an fpqc sheaf if and only if it is a Zariski sheaf and the exactness condition holds for every faithfully flat quasi-compact $S'\to S$.}
\entree[prop]{22691}{6:18:11}{In sheaves on a site, epimorphisms are effective and stable under base change, sums are distributive and base change commutes with inductive limits.}
\face{Tape 2, side B}
\noindent\textit{Recorded on video 2/2 of the groups course (\href{https://www.youtube.com/watch?v=3ULFIl_2Wis}{\texttt{3ULFIl\_2Wis}}); by its content it belongs to the algebraic geometry course.}\par\smallskip
\entree[prop]{22719}{6:18:39}{If $H$ is a subsheaf of groups of $G$, the quotient sheaf $G/H$ by the relation of right translations exists and $G$ acts on it by left translations; if $H$ is normal, $G/H$ is a sheaf of groups.}
\entree[prop]{22847}{6:20:47}{For every scheme $X$, the functor $S\mapsto\mathrm{Hom}(S,X)$ is a sheaf for the fpqc topology.}
\entree[prop]{22976}{6:22:56}{A faithfully flat quasi-compact morphism $S'\to S$ is a universal effective epimorphism: for every scheme $X$, $\mathrm{Hom}(S,X)\to\mathrm{Hom}(S',X)\rightrightarrows\mathrm{Hom}(S'',X)$ is exact, $S''=S'\times_SS'$.}
\entree[prop]{23180}{6:26:20}{The finest topology for which representable functors are sheaves has as covering families the universally strictly epimorphic families.}
\entree[def]{23307}{6:28:27}{fppf topology: in the definition of the fpqc topology, replace quasi-compact by locally of finite presentation; a fundamental system of coverings is given by the families of flat morphisms locally of finite presentation whose images cover $S$.}
\entree[def]{23437}{6:30:37}{Discrete topology: a family $(S_i\to S)$ is covering if one of the $S_i\to S$ admits a section.}
\entree[def]{23437}{6:30:37}{Zariski topology: a family is covering if every point of $S$ has an open neighbourhood $U$ over which one of the $S_i\times_SU\to U$ admits a section.}
\entree[def]{23499}{6:31:39}{\'Etale topology: a fundamental system of coverings is given by the families of \'etale morphisms whose images cover $S$.}
\topic{Quotients and algebraic spaces}
\entree[def]{23629}{6:33:49}{An algebraic space (Artin) is obtained by gluing schemes for the \'etale topology rather than the Zariski topology, that is, by quotienting by an equivalence relation whose projections are \'etale.}
\entree[prop]{23692}{6:34:52}{A representable functor $F$ is an fpqc sheaf and satisfies $F(A)=\varprojlim F(A/\mathfrak{m}^n)$ for every complete noetherian local ring $A$; these necessary conditions are not sufficient.}
\entree[thm]{23822}{6:37:02}{A compact reduced irreducible complex analytic space whose field of meromorphic functions has transcendence degree equal to its dimension is an algebraic space (Artin).}
\entree[prop]{24213}{6:43:33}{If $R\rightrightarrows X$ is an equivalence relation whose projections are faithfully flat quasi-compact, a scheme $Y$ is the quotient of $X$ by $R$ in the sense of fpqc sheaves if and only if $X\to Y$ is faithfully flat quasi-compact and $R=X\times_YX$.}
\entree[prop]{24410}{6:46:50}{For every equivalence relation $R$ on $X$, the quotient sheaf $X/R$ exists, whether or not it is representable by a scheme.}
\entree[thm]{24539}{6:48:59}{If the projections of the equivalence relation $R$ are finite locally free, $X/R$ is representable by a scheme $Y$, with $X\to Y$ finite locally free surjective and $R=X\times_YX$, if and only if $X$ is a union of affine open sets stable under $R$.}
\entree[prop]{24670}{6:51:10}{This condition is equivalent to each equivalence class of $R$, a finite set, being contained in an affine open set.}
\entree[prop]{24799}{6:53:19}{If $X$ is quasi-projective over an affine base, or carries an ample invertible sheaf, every finite subset of $X$ lies in an affine open set; the condition is then satisfied.}
\entree[prop]{24862}{6:54:22}{For a finite group $G$ acting freely on $X$, the quotient $X/G$ exists as a scheme if and only if every orbit is contained in an affine open set.}
\entree[prop]{24930}{6:55:30}{Over a field, for a flat equivalence relation $R$ on a scheme $X$ of finite type, there is a dense open set of $X$, stable under $R$, whose quotient by the induced relation exists as a scheme.}
\entree[thm]{25060}{6:57:40}{If $G$ is an algebraic group over a field and $H\subset G$ a subgroup, the relation defined by $H$ is flat and the quotient $G/H$ is a scheme, a homogeneous space under $G$ by translations.}

\section*{References}\addcontentsline{toc}{section}{References}
\noindent Numbers in square brackets refer to this list. It gathers what the lecturer tells the audience to read, with the moment he says it (video/timestamp), and the other sources used. A ``?'' marks a title the captions do not make certain.
\begingroup\renewcommand{\section}[2]{}\setlength{\itemsep}{1pt}\small
\begin{thebibliography}{99}
\bibitem{bourbaki} N.~Bourbaki, \emph{Algèbre commutative}, Hermann, 1961--1965: ``more than adequate'' for the commutative algebra needed (1/4, 0:03:27; also 2:07:43).
\bibitem{zariskisamuel} O.~Zariski, P.~Samuel, \emph{Commutative Algebra}, Van Nostrand, 1958--1960: ``slightly old-time'', but with the substance (1/4, 0:03:27).
\bibitem{godement} R.~Godement, \emph{Topologie algébrique et théorie des faisceaux}, Hermann, 1958: the standard reference for sheaves (1/4, 0:06:42).
\bibitem{sga4} M.~Artin, A.~Grothendieck, J.-L.~Verdier, \emph{SGA~4, Théorie des topos et cohomologie étale des schémas}, LNM 269, 270, 305, Springer, 1972--1973: exp.~I--II, for generalized sheaves, ``easy reading'' (1/4, 0:06:42--0:07:47).
\bibitem{nagata} M.~Nagata, \emph{Local Rings}, Interscience, 1962: the principle of idealization (1/4, 5:23:45).
\bibitem{fga} A.~Grothendieck, ``Technique de descente et théorèmes d'existence en géométrie algébrique'', Séminaire Bourbaki (\emph{FGA}), 1959--1962: ``the yoga of descent'' (2/4, 2:26:42).
\bibitem{cartan13} A.~Grothendieck, exposés in the \emph{Séminaire Henri Cartan} 13 (1960/61), ``Techniques de construction en géométrie analytique'' (? the seminar's name is uncertain in the captions): the same yoga for analytic geometry (2/4, 8:52:22).
\bibitem{girauddescent} J.~Giraud, ``Méthode de la descente'', \emph{Mém. Soc. Math. France} 2 (1964) (? the name is uncertain in the captions): the general nonsense of descent, ``pretty hard reading'' (4/4, 2:16:15).
\bibitem{demazure} M.~Demazure (? seminar named with a doubt), for effective equivalence relations (4/4, 4:49:17).
\bibitem{gaeta} Alexander Grothendieck, \emph{Introduction to functorial algebraic geometry}, Part~1: \emph{Affine algebraic geometry}, lecture notes written by Federico Gaeta, State University of New York at Buffalo, Department of Mathematics, 168~pp., \url{https://webusers.imj-prg.fr/~leila.schneps/grothendieckcircle/FuncAlg.pdf}.
\end{thebibliography}
\endgroup

\end{document}
